\section{The Galois link II: Selmer and Bloch--Kato bounds}\label{sec:pf:linkN-sel}

\providecommand{\linkNselvcc}{v_{\mathrm{cc}}}
\providecommand{\linkNselsP}{s_{\Pi}}
\providecommand{\linkNseleta}{\eta_{\mathrm{loc}}}
\providecommand{\linkNselmult}{\mathrm{mult}}
\providecommand{\linkNselZg}{Z^{\mathrm{good}}}
\providecommand{\linkNsellst}{\ell^{*}}
\providecommand{\linkNsellen}{\operatorname{len}_{\ell}}
\providecommand{\linkNselSel}{\operatorname{Sel}}
\providecommand{\linkNselG}{\mathcal G}
\providecommand{\linkNselP}{\mathcal P_0}
\providecommand{\linkNselexc}{\mathrm{exc}}
\providecommand{\linkNselExc}{\mathrm{Exc}}
\providecommand{\linkNselerr}{\mathrm{err}}
\providecommand{\linkNselB}{\mathrm{B}}
\providecommand{\linkNselT}{\mathcal T}
\providecommand{\linkNselbsm}{\beta_{\mathrm{sm}}}
\providecommand{\linkNselbqd}{\beta_{qd}}

\subsection{Introduction}\label{sec:pf:linkN-sel:intro}

This section proves the theorem which, in the proof of the Galois link (\cref{thm:linkN}(a)), produces the
\emph{table of local bounds} (step 2 of the outline given there). Its root is \cref{pf:linkN-sel:main}.

\paragraph{What is proved.}
Let $\mathcal F$ be a family (\cref{def:families:closed}) whose members $(D_{x,i},L_{x,i})$ carry the Galois data of
\cref{sec:families:galois}. For a local index $\iota$ of a member, with prime $\ell=\ell_\iota$, and a good eigensystem
$\pi\in\linkNselZg_\iota$, the basic local inequality is
\begin{equation}\label{eq:pf:linkN-sel:bd}
\linkNselvcc(\iota,\pi)\ \le\ \linkNselsP(\iota,\rho,\pi)+\beta+VV(\ell)+v_\ell(h_E)+\ell^{-2}:
\end{equation}
the normalised length of the congruence module of $\pi$ is bounded by the normalised length of a Selmer group of a Galois
representation $\rho_{\iota,\pi}$, plus an error term $\beta$, the lengths of the Tate--Shafarevich groups
$\Sha^1(G_{E,S_\ell},\Ql/\Zl(-e))$ ($0<e<t$), the $\ell$-adic valuation of the class number $h_E$ and $\ell^{-2}$. The main theorem
says: under the literature package $\sfname{Hypotheses}_2$, for a family with good twists, with the norm-choice facts and two extra
conditions at every member (with constants $C,C_T$), and with a bound at the primes $2$ and $3$, there is one constant $C_s$ such
that eventually every member with finite odd Tate--Shafarevich groups admits a \emph{hub datum} $X$ (Galois representations and
local length tables) satisfying fifteen statements, the \emph{table of local bounds} $\linkNselT(L,P,X,C_s)$
(\cref{pf:linkN-sel:concl}). Among them are \eqref{eq:pf:linkN-sel:bd} with $\beta=0$ at all primes $\ell>\linkNsellst$ outside an
exceptional set of controlled size, with explicit $\beta$ at the two auxiliary primes $q,d$, at the small primes and at $2$ and $3$;
a bound for the lattice indices $\linkNseleta(\iota)$ in terms of congruence lengths; a bound for the divisibility defect at the two
generator primes $g_0,g_1$; and numerical side conditions. The summation of this table into the link $\sfname{LinkN}$ of
\cref{def:families:linkN} is not part of this section.

\paragraph{The arguments.}
The middle range of primes is treated by \lref{GlobalTheory.T.star_ad_and_bb_le_of_joint_cond}, proved in another section of
Part II, which we use as a black box; we prove here the two inputs that it needs (\cref{sec:pf:linkN-sel:thmE,sec:pf:linkN-sel:bcd}). The main arguments of this section are:
\begin{enumerate}[label=(\alph*)]
\item a lattice-index bound (\cref{sec:pf:linkN-sel:pass}): if an endomorphism of $M_\iota$ commuting with the Hecke algebra acts
by an integer $A$ on the good eigenspaces and by $0$ on the other eigenspaces of $Z_\iota$, then
$\linkNseleta(\iota)\le\sum_\pi\linkNselmult(\iota,\pi)\bigl(\linkNselvcc(\iota,\pi)+v_\ell(A)\bigr)$;
\item the congruence bound at large primes (\cref{sec:pf:linkN-sel:large}): a global argument bounding the congruence ideal by the
lengths of modules of cocycles attached to the blocks $\rho$ and $\chi_c$ ($1\le c\le t$) of a reducible situation, these being
bounded by the Selmer group, the class number (through Artin reciprocity) and the Tate--Shafarevich groups; the local input is a
theorem on twisted lattice cohomology of Bloch--Kato type (``Theorem TOP'');
\item the bound at $q$ and $d$ (``Theorem B'', \cref{sec:pf:linkN-sel:thmB}), deduced from a global lattice index theorem and a
local lifting theorem (``Theorem A'') by an explicit closing inequality;
\item the bound at the generator primes (\cref{sec:pf:linkN-sel:typ}), deduced from a paired character condition.
\end{enumerate}
The names ``Theorem A, B, E, TOP'' and ``Corollary A'' are those of the formal development \cite{LeanSource2026}; they do not refer
to published results. The shape of \eqref{eq:pf:linkN-sel:bd}, a congruence module bounded by a Selmer group, is that of the
numerical criterion of Wiles and Lenstra \cite{Wiles1995,Lenstra1995}, and the construction of cohomology classes from congruences in
(b) is in the spirit of Ribet \cite{Ribet1976} and of Bella{\"\i}che--Chenevier \cite{BellaicheChenevier2009}; the formal development
does not refer to these works.

\paragraph{Organisation.}
\Cref{sec:pf:linkN-sel:setting} fixes notation and states the main theorem. \Cref{sec:pf:linkN-sel:elem} collects the elementary
parts of the table (monotonicity, small primes, exceptional primes, numerical row conditions, setups).
\Cref{sec:pf:linkN-sel:pass,sec:pf:linkN-sel:typ,sec:pf:linkN-sel:thmB,sec:pf:linkN-sel:large} prove the bounds (a)--(d).
\Cref{sec:pf:linkN-sel:thmE,sec:pf:linkN-sel:bcd} prove the inputs of the middle range, and \cref{sec:pf:linkN-sel:proof}
assembles the proof of the main theorem. \Cref{sec:pf:linkN-sel:local} proves local results used here and in \cref{sec:pf:linkN-adm}.
Several deep local statements of $p$-adic Hodge theory are stated precisely and their proofs are deferred.

% ---------------------------------------------------------------------------
\subsection{Setting and statement}\label{sec:pf:linkN-sel:setting}

We use the data $D$ and size data $L$ of \cref{sec:families}. We write $\linkNselG=\{g_0,g_1\}$ for the set of generator primes of
the host and $\linkNselP$ for its finite set of small primes ($2\in\linkNselP$).

\begin{definition}[Sizes and the large-prime threshold]\label{pf:linkN-sel:sizes}
The scale of $L$ consists of $x>1$, $0<\eps<1/52$ and $J\ge1$; put $\Lambda=x^{\eps}$ and $b=e^{\Lambda}$. The size data provide
positive reals $S_c$ and $\eta_{\mathrm{sz}}$, a real $C_t$, and further quantities ($b'$, $w''$, $b_T(\tau)$ for $\tau$ in a
finite set $T_T$, $\mathrm{de}_2$) which enter only through the numerical lemmas below. The ranks of $D$ are natural numbers
$k,m,t,f,n,r,n_V$ with $t$ odd, $f=2n$, $n_V=r+t$, $r+5\le t$, $m\le r$ and $1\le f$. The host provides a CM field $E$ with class
number $h_E$ and primes $q,d>4tr+1$. The \emph{large-prime threshold} is
\[
\linkNsellst=\linkNsellst(L)=\frac{m\,b}{n\,t^2(N_m+t)\Lambda^2+k\,n^2(S_c+t+\Lambda)},
\]
where $N_m$ is a rank invariant. For $C\in\R$ and a prime $\ell$ put
\begin{gather*}
\linkNselbsm(L,C)=C f m\,b\log x+C x^C\Lambda^2,\\
\linkNselbqd(L,C,\ell)=Cf\bigl(r(b/\ell+S_c+t+\eta_{\mathrm{sz}}+\Lambda)+mb\bigr)+C(frt)^C\Lambda .
\end{gather*}
\formal{Formalism.BasicNotions.LevelStarExponents.lStar}
\formal{GlobalTheory.NLink.midF}
\formal{Formalism.MainNotions.TwistedPlaces.midG}
\end{definition}

\begin{definition}[Local indices and their invariants]\label{pf:linkN-sel:inv}
A local index $\iota$ of $D$ carries (\cref{sec:families:galois}) a prime $\ell=\ell_\iota$, a complete discrete valuation ring
$\OO=\OO_\iota$ finite over $\Zl$ with fraction field $K$, a finite set $X_\iota$ of eigensystems with subsets
$\linkNselZg_\iota\subseteq Z_\iota\subseteq X_\iota$, a Hecke algebra $T_\iota\subseteq\OO^{X_\iota}$ which is a local ring and
separates the points of $X_\iota$, a free $\OO$-module $M_\iota$ of finite rank with a faithful action of $T_\iota$, a lift index
$1\le j_\iota\le m$, a place $w_\iota$ of $\mathit{LL}$ above $\ell$, characters $\chi_{\iota,c}\colon G_E\to\OO^\times$ with
$\chi_{\iota,c}=\chi_{\iota,1}\chi_{\mathrm{cyc}}^{1-c}$ for $1\le c\le t$, and disjoint finite sets $S_\iota,Q_\iota$ of split places
of $E$ prime to $\ell$, with elements $\zeta_{\iota,v}\in\OO$ for $v\in Q_\iota$. All places $w_\iota$ are twisted. For $x\in X_\iota$
let $V_x\subseteq M_\iota\otimes_\OO K$ be the common eigenspace of $T_\iota$ with eigenvalues $(f(x))_{f\in T_\iota}$, and
$V_S=\sum_{x\in S}V_x$. For an $\OO$-module $N$ the \emph{normalised length} is
$\linkNsellen(N)=\len_\OO(N)/\len_\OO(\OO/\ell\OO)\in\Q$, where a module of infinite length has normalised length $0$. Then:
\begin{itemize}
\item $\linkNseleta(\iota)=\linkNsellen\bigl(((M_\iota+V_{X_\iota\smallsetminus\linkNselZg_\iota})\cap V_{\linkNselZg_\iota})/
(M_\iota\cap V_{\linkNselZg_\iota})\bigr)$, with $M_\iota$ identified with its image in $M_\iota\otimes K$, and
$\eta_j(w)=\sum_{\iota:(j_\iota,w_\iota)=(j,w)}\linkNseleta(\iota)$ for $1\le j\le m$;
\item $\linkNselmult(\iota,\pi)=\dim_KV_\pi$;
\item $\linkNselvcc(\iota,\pi)=\linkNsellen(\OO/\mathfrak c_\pi)$, where $\mathfrak c_\pi=\{c(\pi):c\in T_\iota,\ c(y)=0\text{ for all }
y\notin Z_\iota\}$;
\item for a family $\rho$ of continuous representations $\rho_{\iota,\pi}\colon G_E\to\GL_r(\OO_\iota)$,
$\linkNselsP(\iota,\rho,\pi)=\linkNsellen(\linkNselSel(\iota,\rho,\pi))$, where $\linkNselSel(\iota,\rho,\pi)\subseteq
H^1(G_E,W_{\iota,\rho,\pi})$ is the Selmer group with the local conditions at $\ell$ given by the local theory of $D$ and relaxed at
$S_\iota\cup Q_\iota$, and $W_{\iota,\rho,\pi}$ is $(K/\OO)^r$ with $G_E$ acting through $\rho_{\iota,\pi}\chi_{\iota,1}^{-1}$; the
Selmer group is \emph{finite} if it has finite length;
\item $\mathrm{sh}(\ell,e)\ge0$ ($e\in\N$) are rational data of $D$ and $VV(\ell)=\sum_{1\le e<t}\mathrm{sh}(\ell,e)$; we say that $D$
has \emph{finite odd Tate--Shafarevich groups} if $\Sha^1(G_{E,S_p},\Q_p/\Z_p(-e))$ has finite length for all primes $p$ and
$0<e<t$, in which case $\mathrm{sh}(p,e)$ equals this length (a field of $D$);
\item $\eps_\ell=\ell^{-2}$; $\nu_L(\ell)\ge0$ is a cyclotomic invariant of $D$ at $\ell$; $\#G\ge1$ is an integer attached to the
rows of $D$; $\delta_\iota=\sum_{0<k<t}(s_0(\iota)+v_\ell(k))$, where $s_0(\iota)$ is the local level exponent of $\iota$ (equal to $1$ if
$\ell\notin\linkNselP$).
\end{itemize}
We write $\mathrm{lw}(w)=\log\Nm(w)/[\mathit{LL}:\Q]$ and $e(w)$ for the ramification index of a place $w$ of $\mathit{LL}$, $d_{Z_0}(j)$
for the dimension of the $j$-th zero space of $L$, and $n_\lambda$ for the local degree of a place $\lambda$ of $E$.
\formal{Formalism.BasicNotions.lenUnit}
\formal{Formalism.BasicNotions.NData.etaLoc}
\formal{Formalism.BasicNotions.NData.vcc}
\formal{Formalism.MainNotions.N.Br.selPi}
\formal{Formalism.MainNotions.TwistedPlaces.ShaFinOdd}
\formal{Formalism.MainNotions.TwistedPlaces.nuL}
\end{definition}

By \lref{Formalism.MainNotions.TruncatedSeriesCountingIdentities.Br.etaLoc_eq} (another section), $\linkNseleta(\iota)$ is the
normalised length of the quotient displayed above; this is the description used below.

\begin{definition}[Hub data, killing, the local bound]\label{pf:linkN-sel:hub}
A \emph{hub datum} $X$ for $D$ consists of families of continuous representations $r_{\iota,\pi}\colon G_E\to\GL_{n_V}(\OO_\iota)$ and
$\rho_{\iota,\pi}\colon G_E\to\GL_r(\OO_\iota)$ ($\iota$ a local index, $\pi\in X_\iota$), real tables $\mathrm{len}_0(\iota,\pi,\lambda)$
($\lambda$ a place of $E$) and $\linkNselerr(\iota,\pi)$, a partition of the places above $\ell_\iota$ into rows and columns, a band
predicate on pairs $(\iota,\lambda)$, natural numbers $c_B(\iota),\dots$, a finite set $\linkNselexc(X)$ of primes, an integer $\eta(X)$
and twisting data. Its \emph{table} $\mathrm{tab}_X$ has $\mathrm{len}(\iota,\lambda)=\sup_{\pi\in\linkNselZg_\iota}
\mathrm{len}_0(\iota,\pi,\lambda)$ and $BB(\iota)=\sup_{\pi\in\linkNselZg_\iota}\linkNselerr(\iota,\pi)$ (both $0$ if
$\linkNselZg_\iota=\emptyset$). For a family $\rho$, the \emph{killing condition} $\mathrm{Kill}(\rho)$ means: for every local index
$\iota$ and $\pi\in\linkNselZg_\iota$ the Selmer group $\linkNselSel(\iota,\rho,\pi)$ is finite and
$\linkNselsP(\iota,\rho,\pi)\le\nu_L(\ell_\iota)+\eps_{\ell_\iota}$, and $t\sum_{v\in Q_\iota}\linkNsellen(\OO/(\zeta_{\iota,v}-1))\le
\eps_{\ell_\iota}$. For $\beta\in\R$, the \emph{local bound} $\linkNselB(\rho,\iota,\pi;\beta)$ is the inequality
\eqref{eq:pf:linkN-sel:bd}, i.e.\ $\linkNselvcc(\iota,\pi)\le\linkNselsP(\iota,\rho,\pi)+\beta+VV(\ell)+v_\ell(h_E)+\eps_\ell$ with
$\ell=\ell_\iota$.
\formal{Formalism.MainNotions.N.KillingConditions.HubData}
\formal{Formalism.MainNotions.N.KillingConditions.KillAt}
\formal{GlobalTheory.NLink.Bd}
\end{definition}

\begin{definition}[The numerical row conditions]\label{pf:linkN-sel:rownumdef}
Put $h_B=b+S_c+\eta_{\mathrm{sz}}+t+2$, $e_P(C,\ell)=C(1+h_B/\ell+\log_\ell h_B)$, $\eta_3=\eta_{\mathrm{sz}}+t+2$,
$w_B=\max\bigl(w''+\eta_{\mathrm{sz}},\max_{\tau}b_T(\tau)+b'/x+\eta_{\mathrm{sz}}\bigr)$, $w_P=w_B+(t+1)/2$, $G_P=w_P+S_c+\eta_3$,
$\ell_A=k\Lambda/(m\,\mathrm{de}_2)$ and $j_A(\ell)=\max(1,\lceil\log_{\ell-1}\ell_A\rceil)$, and
$\mathrm{flat}(\iota,\lambda,\ell)=(w_P+\eta_3)/\ell+S_c+\Lambda+\eta_3\cdot[\text{$(\iota,\lambda)$ is a band pair}]$. Let $L_w$ be a
host constant and $\mathrm{ord}_H(\ell)$ the least $o>0$ with $\ell^o\equiv\pm1\pmod{2^{e+2}}$ ($e$ a rank parameter); primes are
\emph{split} or \emph{inert} according to the class assigned to them by the host. For a table $\mathrm{tb}$ and $C\in\R$ the numerical
row conditions are (R1)--(R4), required for every local index $\iota$ with $\ell=\ell_\iota\notin\linkNselG$ and every $\lambda\mid\ell$,
and (R5), required for every band pair:
\begin{enumerate}[label=(R\arabic*)]
\item if $\ell\ge5$ is split and ($\ell\le\linkNsellst$ or $\ell\in\linkNselexc$):
$2\sum_{\lambda\mid\ell}\mathrm{len}(\iota,\lambda)\le\sum_{\lambda\mid\ell}(2n_\lambda m+1)e_P(C,\ell)$;
\item if $\ell$ is inert, $\ell<\ell_A$ and $4L_w<\mathrm{ord}_H(\ell)$:
$\mathrm{len}(\iota,\lambda)\le C(n_\lambda+1)r\bigl(h_B/(\ell-1)^{j_A(\ell)+1}+(G_P+1)(j_A(\ell)+1)+\Lambda\bigr)$;
\item if $\ell$ is inert, $\ell_A\le\ell$ and ($\ell\le\linkNsellst$ or $\ell\in\linkNselexc$):
$\mathrm{len}(\iota,\lambda)\le C(n_\lambda+1)r\bigl(h_B/\ell^2+\mathrm{flat}(\iota,\lambda,\ell)\bigr)+Cn_\lambda m\,e_P(C,\ell)$;
\item if $\ell$ is inert, $\ell<\ell_A$ and $\mathrm{ord}_H(\ell)\le4L_w$: the same bound as in (R3);
\item for every band pair $(\iota,\lambda)$ and every prime $\ell$ with $\lambda\mid\ell$ there is $s$ with $w''\le s\le w_B$ and
$|\ell s-b|\le C(\ell\eta_3+\ell S_c+b/\ell+s+\ell\Lambda)$.
\end{enumerate}
\formal{GlobalTheory.NLink.RowNum}
\end{definition}

\begin{definition}[The table of local bounds]\label{pf:linkN-sel:concl}
Let $P$ be a parameter choice for $(D,L)$, $X$ a hub datum and $C\in\R$; write $\rho=X.\rho$ and $\ell=\ell_\iota$. The \emph{table
of local bounds} $\linkNselT(L,P,X,C)$ is the conjunction of:
\begin{description}[style=nextline,leftmargin=1.5em,font=\normalfont\itshape]
\item[pos] $C\ge0$;
\item[pass] for every $\iota$ with $\ell\notin\linkNselG$:
$\linkNseleta(\iota)\le\sum_{\pi\in\linkNselZg_\iota}\linkNselmult(\iota,\pi)\bigl(\linkNselvcc(\iota,\pi)+v_\ell(\#G)\bigr)$;
\item[mult] for $1\le j\le m$ and every twisted place $w$:
$\sum_{\iota:(j_\iota,w_\iota)=(j,w)}\sum_{\pi\in\linkNselZg_\iota}\linkNselmult(\iota,\pi)$ is at most the rank of the $j$-th zero space
of $L$ at $w$;
\item[rank] that rank equals $d_{Z_0}(j)$ for $1\le j\le m$ and all twisted $w$;
\item[kill] $\mathrm{Kill}(\rho)$;
\item[num] the numerical row conditions (R1)--(R5) for $\mathrm{tab}_X$ and $C$;
\item[mid] for $\ell\notin\linkNselG$ in the \emph{middle range} ($\ell\notin\linkNselP$, $\ell\ne3$, $\ell\notin\{q,d\}$, and not both
$\ell>\linkNsellst$ and $\ell\notin\linkNselexc(X)$) and $\pi\in\linkNselZg_\iota$: $\linkNselvcc\le\linkNselsP+\sum_{\lambda\mid\ell}
\mathrm{len}_0(\iota,\pi,\lambda)+VV(\ell)+v_\ell(h_E)+\linkNselerr(\iota,\pi)+\eps_\ell$;
\item[bb] in the same range, $\linkNselerr(\iota,\pi)\le C f n t^2 n_V\Lambda^2$;
\item[large] $\linkNselB(\rho,\iota,\pi;0)$ whenever $\ell>\linkNsellst$, $\ell\notin\linkNselG$, $\ell\notin\linkNselexc(X)$ and
$\pi\in\linkNselZg_\iota$;
\item[exc] $\#\linkNselexc(X)\le(\#T_T+2)\bigl(S_c+3\eta_{\mathrm{sz}}+t+8b/\linkNsellst+9\bigr)$, and every element of $\linkNselexc(X)$
is at most $3b$;
\item[two] there is a table $\mathrm{tf}$ with $\sum_{\lambda\mid\ell}\mathrm{len}_{\mathrm{tf}}(\iota,\lambda)+BB_{\mathrm{tf}}(\iota)
\le\linkNselbsm(L,C)$ and $\linkNselB\bigl(\rho,\iota,\pi;\sum_{\lambda\mid\ell}\mathrm{len}_{\mathrm{tf}}(\iota,\lambda)+BB_{\mathrm{tf}}(\iota)\bigr)$
for all $\iota$ with $\ell\in\linkNselP$ or $\ell=3$ and all $\pi\in\linkNselZg_\iota$;
\item[ram] $\linkNselB(\rho,\iota,\pi;\linkNselbqd(L,C,\ell))$ for $\ell\in\{q,d\}$ and $\pi\in\linkNselZg_\iota$;
\item[typ] for $1\le j\le m$ and $\ell\in\linkNselG$:
$\sum_{w}\mathrm{lw}(w)e(w)\eta_j(w)\le\log\ell\cdot d_{Z_0}(j)\bigl(Cx^C\eta_{\mathrm{sz}}+2\eps_\ell+VV(\ell)+v_\ell(h_E)\bigr)$, the sum
over the twisted places $w$ of the support of $L$ above $\ell$;
\item[par] the parameters are polynomially bounded with constant $C$: $P.C_{Lw}\le C$, $P.C_{YL}\le C$, $k,n,t\le x^C$,
$x^{-\eps}\le\mathrm{de}_2$, $\log h_E\le Cx^C$ and $\log\#G\le Cx^C\Lambda$;
\item[small] $3,q,d$ and all small primes are at most $\linkNsellst$, the prime $3$ is not inert, and $\linkNselbqd(L,C,\ell)\le2Cfmb$
for $\ell\in\{q,d\}$.
\end{description}
A family $\mathcal F$ \emph{admits tables of local bounds}, $\linkNselT(\mathcal F)$, if there are $C,x_0\in\R$ such that for all
$x\ge x_0$ with $x>x_1$ and all $i\in I_x$ for which $D_{x,i}$ has finite odd Tate--Shafarevich groups there is a hub datum $X$ for
$D_{x,i}$ with $\linkNselT(L_{x,i},P_{x,i},X,C)$, $P_{x,i}$ being the parameter choice of the member.
\formal{GlobalTheory.NLink.Concl}
\formal{GlobalTheory.NLink.ConclF}
\end{definition}

\begin{definition}[Setups and the family hypotheses]\label{pf:linkN-sel:hyps}
$\sfname{Hypotheses}_2$ is the conjunction of seven literature statements: Fontaine's crystalline period ring
\cite{Fontaine1982,Fontaine1994}, global Artin reciprocity \cite{Artin1927,Neukirch1999}, the local Bloch--Kato theory (dimension
formula, rational and integral duality for $H^1_f$) \cite{BlochKato1990}, local Tate duality, the local Euler characteristic formula
and local finiteness for finite extensions of $\Ql$ \cite{Tate1963,Serre1997,NeukirchSchmidtWingberg2008}, and an integral comparison
of Bhatt--Scholze type \cite{BhattScholze2022,BhattScholze2023}. Let $D$ be a datum with a norm choice $\mathrm{nch}$
(\cref{sec:families:galois}), local Hodge--Witt data $HW$, a hub datum $X$ and $C_2\in\R$, and let $R(C_2)=(C_2,320,25,23,32,32)$ (the
\emph{row constants}). A \emph{setup} $(L,\mathrm{nch},HW,X,C_2)$ consists of: the norm-choice facts for $(D,\mathrm{nch})$; the
statement that $X$ is \emph{canonical} for $R(C_2)$, i.e.\ its representations form a \emph{point family} (each $r_\iota$ is good,
each $\rho_{\iota,\pi}$ is continuous and, for $\pi\in Z_\iota$, a complement of the characters $\chi_{\iota,c}$ in $r_{\iota,\pi}$) and
its tables are the canonical tables built with $R(C_2)$; the \emph{admissibility of the parameters} ($\eta_{\mathrm{sz}}+t+S_c<w''$,
$3w''\le b$, $b_T(\tau)+b'/x\le b$, $t\le\eta_{\mathrm{sz}}$); the inequality $S_c+\eta_{\mathrm{sz}}+t+2\le b$; $\Lambda\ge8$; a host
split bound $\ge60$; $\mathrm{Kill}(X.\rho)$; and finite odd Tate--Shafarevich groups for $D$. For a family $\mathcal F$ and $C\in\R$,
$C_T\in\N$ we consider:
\begin{itemize}
\item the \emph{facts hypothesis} $(\mathrm F)_{C,C_T}$: for all $x>x_1$ and $i\in I_x$ the norm-choice facts hold for
$(D_{x,i},\mathrm{nch}_{x,i})$, and $D_{x,i}$ satisfies the \emph{commuting endomorphism condition} (for every $\iota$ with
$\ell_\iota\notin\linkNselG$ there is an $\OO_\iota$-endomorphism of $M_\iota$ commuting with $T_\iota$ which acts on $V_x$ by $\#G$ for
$x\in\linkNselZg_\iota$ and by $0$ for $x\in Z_\iota\smallsetminus\linkNselZg_\iota$) and the paired character condition
$\mathrm{PC}(C,C_T)$ of \cref{pf:linkN-sel:pcc};
\item $\sfname{TwGood}(\mathcal F)$: the twisting characters of the norm choices are eventually good (\cref{sec:families:galois});
\item the \emph{bound at $2$ and $3$}, $(\mathrm B_{23})$: there are $C_a,x_a$ such that for all $x\ge x_a$, $x>x_1$, $i\in I_x$ and every
setup $(L_{x,i},\mathrm{nch}_{x,i},HW,X,C_2)$ (with any Fontaine and comparison data), $\linkNselB(X.\rho,\iota,\pi;\linkNselbsm(L_{x,i},C_a))$
holds for every $\iota$ with $\ell_\iota=2$ or $\ell_\iota=3\notin\linkNselP$ and every $\pi\in\linkNselZg_\iota$.
\end{itemize}
\formal{GlobalTheory.NLink.Hypotheses₂}
\formal{GlobalTheory.NLink.SetupAlt}
\formal{GlobalTheory.NLink.FactsFAlt}
\formal{GlobalTheory.NLink.BoundAtTwoThreeFamilyFonAlt}
\end{definition}

\begin{theorem}[Table of local bounds]\label{pf:linkN-sel:main}
Assume $\sfname{Hypotheses}_2$. Let $\eps\in\R$, $\mathcal F\in\sfname{FamP}(\eps)$, $C\ge0$ and $C_T\in\N$, and suppose that
$\mathcal F$ satisfies $\sfname{TwGood}$, $(\mathrm F)_{C,C_T}$ and $(\mathrm B_{23})$. Then $\mathcal F$ admits tables of local bounds.
More precisely, let $C_0$ be the constant of \cref{pf:linkN-sel:thmE}, $c_1=2080$ and $Q=260X_0^2X_2^2(X_1+X_2)$
(\cref{pf:linkN-sel:thmBmain}), $C_1\ge1$ the constant of the parameter bound along $\mathcal F$
(\lref{GlobalTheory.T.admissible_parameters_family_const}), $C_a$ the constant of
$(\mathrm B_{23})$ and $C_b$ that of \cref{pf:linkN-sel:starS}. Then the conclusion holds with the single constant
\begin{equation}\label{eq:pf:linkN-sel:Cs}
C_s=\max\bigl(3C_2+1,\ 22365,\ C_G,\ 2000(C+1)+5C_1,\ C_1,\ \max(0,C_a,C_b),\ C_T\bigr),
\end{equation}
where $C_2=\max(C_0,1)$, $C_G=\max(c_1,\Sigma(Q),\deg Q,1)$ and $\Sigma(Q)$ is the sum of the coefficients of $Q$.
\formallabel{GlobalTheory.NLinkClose.main_alt_alt}
\end{theorem}

\begin{remark}\label{pf:linkN-sel:remmain}
(i) All seven fields of $\sfname{Hypotheses}_2$ are used. (ii) The hypothesis $(\mathrm B_{23})$ is used in an essential way: it gives the
field \emph{two} at $\ell=2$ and at $\ell=3\notin\linkNselP$. (iii) The representations $\rho$ of the conclusion are part of
the hub datum $X$, which is chosen for each member; the local bounds concern the Selmer groups of these representations.
(iv) Nothing is asserted about members some of whose groups $\Sha^1(G_{E,S_p},\Q_p/\Z_p(-e))$, $0<e<t$, are infinite. (v) Since normalised
lengths of infinite modules are $0$, the local bounds are meaningful only together with the finiteness of the Selmer groups, which is
part of the field \emph{kill}.
\end{remark}

% ---------------------------------------------------------------------------
\subsection{Elementary parts of the table}\label{sec:pf:linkN-sel:elem}

\subsubsection{Monotonicity in the constant}

\begin{lemma}\label{pf:linkN-sel:mono}
Let $D$, $L$, a parameter choice $P$ and a hub datum $X$ be given, and let $C\le C'$ with $C'\ge0$. Then
$\linkNselbsm(L,C)\le\linkNselbsm(L,C')$, and each of the fields \emph{two}, \emph{typ}, \emph{bb} and \emph{par} of
\cref{pf:linkN-sel:concl} for the constant $C$ implies the same field for $C'$.
\formallabel{GlobalTheory.T.sizeConditions_mono}
\formallabel{GlobalTheory.NLinkMono.main}
\formallabel{GlobalTheory.NLink.midF_mono}
\end{lemma}

\begin{proof}
Since $x>1$ we have $\log x\ge0$, and $b,\Lambda>0$. Moreover $Cx^C\le C'x^{C'}$: if $C\ge0$ this is the product of $C\le C'$ and
$x^C\le x^{C'}$ between nonnegative numbers; if $C<0$ then $Cx^C\le0\le C'x^{C'}$. Multiplying $C\le C'$ by $fmb\log x\ge0$ and
$Cx^C\le C'x^{C'}$ by $\Lambda^2$ gives $\linkNselbsm(L,C)\le\linkNselbsm(L,C')$. For \emph{two} keep the same table: its bound
$\le\linkNselbsm(L,C)$ implies $\le\linkNselbsm(L,C')$. For \emph{typ}, multiply $Cx^C\eta_{\mathrm{sz}}\le C'x^{C'}\eta_{\mathrm{sz}}$
by $\log\ell\cdot d_{Z_0}(j)\ge0$. For \emph{bb}, multiply $C\le C'$ by $fnt^2n_V\Lambda^2\ge0$. For \emph{par}, the bounds on
$P.C_{Lw},P.C_{YL}$ pass to $C'$, $x^C\le x^{C'}$ handles $k,n,t$, and $Cx^C\le C'x^{C'}$ (times $\Lambda\ge0$ for $\log\#G$) handles
$\log h_E$ and $\log\#G$; the condition on $\mathrm{de}_2$ does not involve $C$.
\end{proof}

\begin{lemma}\label{pf:linkN-sel:starF}
Let $\rho$ be a family of representations and $C\in\R$, and suppose that $\linkNselB(\rho,\iota,\pi;\linkNselbsm(L,C))$ holds for all
$\iota$ with $\ell_\iota=2$, or $\ell_\iota=3\notin\linkNselP$, or $\ell_\iota\in\linkNselP\smallsetminus\{2\}$, and all
$\pi\in\linkNselZg_\iota$. Then the field \emph{two} holds for $(L,\rho,C)$.
\formallabel{GlobalTheory.NLink.starF_of}
\end{lemma}

\begin{proof}
Take the table with $\mathrm{len}=0$ and $BB=\linkNselbsm(L,C)$; its length bound is an equality. If $\ell_\iota\in\linkNselP$ or
$\ell_\iota=3$, then $\ell_\iota=2$, or $\ell_\iota\in\linkNselP\smallsetminus\{2\}$, or $\ell_\iota=3\notin\linkNselP$, and the
hypothesis applies.
\end{proof}

\subsubsection{The small odd primes}

For $A,B\in\R$ and a small prime $\ell$ let $\mathrm{cost}(L,A,B,\ell)=Afm\,h_B\,s_\ell+Bf^2t^2n_V\,\lambda_\ell$, where $h_B$ is as in
\cref{pf:linkN-sel:rownumdef}, $s_\ell=\max(1,v_\ell(N_E))$ for an integer $N_E$ attached to the host, and $\lambda_\ell\ge0$ is an
explicit logarithmic quantity.

\begin{lemma}\label{pf:linkN-sel:starS}
Assume $\sfname{Hypotheses}_2$. For every family $\mathcal F$ there are $C_b,x_b\in\R$ such that for all $x\ge x_b$ with $x>x_1$, all
$i\in I_x$ and every setup $(L_{x,i},\mathrm{nch}_{x,i},HW,X,C_2)$, the local bound $\linkNselB(X.\rho,\iota,\pi;\linkNselbsm(L_{x,i},C_b))$
holds for all $\iota$ with $\ell_\iota\in\linkNselP\smallsetminus\{2\}$ and all $\pi\in\linkNselZg_\iota$.
\formallabel{UpperTheory.N.starSF_of_prints}
\formallabel{UpperTheory.N.starSF_of_rowC_alt}
\formallabel{UpperTheory.N.exists_nonneg_c_constants_of_prints}
\formallabel{UpperTheory.N.n_at_bounds_of_setup}
\end{lemma}

\begin{proof}
\emph{Step 1: the cost bound at every setup.} Let a setup be given, $\iota$ with $\ell=\ell_\iota\in\linkNselP$, $\ell\ne2$, and
$\pi\in\linkNselZg_\iota$. The host conditions imply that $\ell$ is not a generator prime and, since $\ell$ divides the level of the
host, that $-d$ is a square modulo $\ell$. The lemma \lref{UpperTheory.N.decl_split} gives a decomposition of the places
above $\ell$ adapted to this splitting, and the lemma \lref{UpperTheory.N.conclusion_split_alt} gives local length bounds at
$\pi$ for it (from $\sfname{Hypotheses}_2$, the setup, and the local statements \lref{UpperTheory.N.locAll} at all primes). The
lemma \lref{UpperTheory.N.bound_smallPrime_of_splitBound_cyclotomic} states that such local bounds, together with the
killing condition, finite odd Tate--Shafarevich groups, the column bound \lref{UpperTheory.N.h0_le} and the elementary inequalities
$\lfloor h_B\rfloor\le h_B$, $b\le h_B$, $\lfloor2b\rfloor\le2b$ and $1\le s_0(\iota)$, imply
$\linkNselB(X.\rho,\iota,\pi;\mathrm{cost}(L,6000,30000,\ell))$.

\emph{Step 2: uniformity.} Let $C_0\ge1$ and $x_p$ be the constants of the parameter bound
\lref{GlobalTheory.T.admissible_parameters_family_const}, which states that along $\mathcal F$ the field \emph{par} holds
eventually with a constant $C_0\ge1$ (in particular $n,t\le x^{C_0}$). Put $A=6000$, $B=30000$, $C_b=18AC_0+2B(3C_0+9)+5C_0+1$ and
$x_b=\max(3,x_p)$. For $x\ge x_b$ we have $\log x\ge1$, and the estimate \lref{UpperTheory.N.Num.cost_le} (deferred) states that
then $\mathrm{cost}(L,A,B,\ell)\le\linkNselbsm(L,C)$ for every small prime $\ell$ and every $C\ge C_b$. Since the local bound is monotone
in $\beta$, Step 1 gives the claim.
\end{proof}

\subsubsection{Exceptional primes}

\begin{definition}[Exceptional primes]\label{pf:linkN-sel:excdef}
For size data $L$ put $\sigma^+=\lfloor2b/\linkNsellst\rfloor$, call \emph{base sizes} the reals $W\in\{b,w''\}\cup\{b_T(\tau)+
\lfloor b'/x\rfloor:\tau\in T_T\}$, and let $\mathrm{win}(W)$ be the set of integers $n$ with $W-\eta_{\mathrm{sz}}-2\sigma^+-3\le n
\le W+S_c+2\eta_{\mathrm{sz}}+2\sigma^++4$. The \emph{exceptional set} $\linkNselExc(L)$ is the set of primes $\ell>\linkNsellst$ dividing
some element of $\mathrm{win}(W)$ for some base size $W$.
\formal{GlobalTheory.IndefiniteRowSideFullHecke.BaseSizes.Exc}
\end{definition}

\begin{lemma}\label{pf:linkN-sel:exc}
Let $L$ be size data of a member of a family. Then $\linkNselExc(L)$ is finite,
$\#\linkNselExc(L)\le(\#T_T+2)\bigl(S_c+3\eta_{\mathrm{sz}}+t+8b/\linkNsellst+9\bigr)$, and every element of $\linkNselExc(L)$ is at
most $3b$. Consequently, for every family and all $x>x_1$, $i\in I_x$, the field \emph{exc} holds for $\linkNselexc=\linkNselExc(L_{x,i})$.
\formallabel{GlobalTheory.NLink.ExcBound.excOK_at}
\formallabel{GlobalTheory.NLink.excOK_family}
\formallabel{GlobalTheory.T.excludedPrimes_few_and_small_along_family}
\end{lemma}

\begin{proof}
The size conditions of the members give: $\linkNsellst>0$; $2\linkNsellst<b$; $2\eta_{\mathrm{sz}}+t+8b/\linkNsellst+7<\linkNsellst$;
$b_T(\tau)\le\frac54b'$, $b'\le b$ and $w''\le b$; and all elements of the windows are positive (numerical lemmas with deferred proofs,
\lref{GlobalTheory.Survey.CongruenceBoundParameters.two_lStar_lt_b} for the second and
\lref{GlobalTheory.Survey.CongruenceBoundParameters.window_pos} for the last). From the third inequality,
$8b/\linkNsellst<\linkNsellst$, i.e.\ $8b<(\linkNsellst)^2$. Every base size satisfies $W\le2b$, and $S_c+2\eta_{\mathrm{sz}}+2\sigma^++4<b$;
hence every $n\in\mathrm{win}(W)$ satisfies $0<n<3b$.

If $n\in\mathrm{win}(W)$ had two distinct prime factors $p,p'>\linkNsellst$, then $pp'\mid n$ and $n\ge pp'>(\linkNsellst)^2>8b>3b$, a
contradiction. Hence each $n$ contributes at most one prime to $\linkNselExc(L)$, and
\[
\#\linkNselExc(L)\le\sum_W\#\mathrm{win}(W)\le(\#T_T+2)\bigl(S_c+3\eta_{\mathrm{sz}}+4\sigma^++8\bigr),
\]
since there are at most $\#T_T+2$ base sizes and an interval of length $S_c+3\eta_{\mathrm{sz}}+4\sigma^++7$ contains at most
$S_c+3\eta_{\mathrm{sz}}+4\sigma^++8$ integers. As $4\sigma^+\le8b/\linkNsellst$ and $t\ge0$, the stated bound follows; in particular
$\linkNselExc(L)$ is finite. An element $\ell\in\linkNselExc(L)$ divides some $n$ with $0<n<3b$, so $\ell\le3b$. The family version holds
from the threshold $x_1$ on.
\end{proof}

\subsubsection{The special primes}

\begin{lemma}\label{pf:linkN-sel:small}
For every family $\mathcal F$ and every $C\ge1$ there is $x_0$ such that for all $x\ge x_0$ with $x>x_1$ and all $i\in I_x$: the field
\emph{small} holds for $(L_{x,i},C)$, $\Lambda\ge8$, the host split bound of $D_{x,i}$ is $\ge60$, and $C\le b$.
\formallabel{GlobalTheory.T.specialPrimes_small_along_family}
\formallabel{GlobalTheory.SpecialSmall.special_primes_small}
\formallabel{GlobalTheory.SpecialSmall.at_one}
\end{lemma}

\begin{proof}
For $x>x_1$ the scale of $L_{x,i}$ has parameter $x$ and the fixed $\eps>0$ of the family, and $C_t$ and $\eps_z=\eps/J$ are family
constants. Since $x^{\eps}\to\infty$ and $e^{x^\eps}$ grows faster than any power of $x$, for large $x$: $x\ge3$, $x^\eps\ge\max(16,C)$,
$x^{\eps_z/5}\ge60$, $4\cdot2^Cx^{(11+|C_t|)C+\eps}\le e^{x^\eps}$ and $32x^{5+\eps}\le e^{x^\eps}$. Assume these. The split bound is at least
$\lfloor x^{\eps_z/5}\rfloor\ge60$ by a tie of the size data, $\Lambda=x^\eps\ge16$, and $C\le x^\eps\le e^{x^\eps}=b$. The size conditions
give $k\le x^5$ and $t\le x^{|C_t|+1}$, and $f\le k$, $r\le2k$; hence $frt\le2x^{11+|C_t|}$ and $(frt)^C\Lambda\le2^Cx^{(11+|C_t|)C+\eps}\le b/4$.
They also give $S_c,\eta_{\mathrm{sz}},t\le b'$ and $16rb'\le mb$, while $32x^{5+\eps}\le b$ gives $16r\Lambda\le mb$; together
$r(S_c+t+\eta_{\mathrm{sz}}+\Lambda)\le mb/4$. For $\ell\in\{q,d\}$ we have $\ell>4tr+1\ge2r$, so $rb/\ell\le mb/2$, and
\[
\linkNselbqd(L,C,\ell)\le Cf\bigl(\tfrac{mb}2+\tfrac{mb}4+mb\bigr)+C\tfrac{b}4\le2Cfmb .
\]
The remaining parts of the field \emph{small} ($3,q,d$ and the small primes are at most $\linkNsellst$, and $3$ is not inert) are size
conditions of the member.
\end{proof}

\subsubsection{The numerical row conditions}

\begin{lemma}\label{pf:linkN-sel:rownum}
Let $R=(C_2,C_B,C_C,C_A,C_B',C_K)$ be row constants and $X$ a hub datum whose tables are the canonical tables for $R$. Let $C$ satisfy
$1\le C_2$, $C_B,C_C,C_A\ge0$, $1\le C$, $3C_2+1\le C$, $K_1\le C$, $K_2+1\le C$, $K_3\le C$, $6C_B+1\le C$ and $2(3C_A+1)\le C$, where
$K_1=C_B+48$, $K_2=6C_B+12+6C_C$ and $K_3=3C_AC_B+12C_A+9$ (we then say that $R$ is \emph{admissible for} $C$). Then (R1)--(R4) hold for
$\mathrm{tab}_X$ and $C$; if moreover $S_c+\eta_{\mathrm{sz}}+t+2\le b$, then (R5) holds as well.
\formallabel{GlobalTheory.NLink.RowA.main}
\formallabel{GlobalTheory.NLink.RowA.two_eA_le}
\formallabel{GlobalTheory.NLink.localBoundTableFromData_rowB_rowEntries}
\formallabel{GlobalTheory.NLink.RowBSize.lenRowB_le}
\formallabel{GlobalTheory.T.mediumInert_smallInert_lengthBound_tabX}
\formallabel{GlobalTheory.NLink.localBoundTable_rowsCD_admissible}
\formallabel{GlobalTheory.T.localBoundTableFromData_bandCondition}
\formallabel{GlobalTheory.NBand.bandPlaceWeightApproximation_tabX}
\end{lemma}

\begin{proof}
All canonical tables are explicit, and each condition is a direct comparison.

(R1) At a split $\ell\ge5$ the canonical table is $\mathrm{len}_0(\iota,\pi,\lambda)=(2n_\lambda m+1)e_A(\ell)$ if $\lambda$ is a column
place and $0$ otherwise, where $e_A(\ell)=\lceil C_A(2+2h_N/\ell+\log_\ell(2h_N+1))\rceil$ and $h_N=\lfloor h_B\rfloor$; it does not depend
on $\pi$. Since $h_N\le h_B$, $1\le h_B$ and $2h_N+1\le\ell h_B$, we get $2+2h_N/\ell+\log_\ell(2h_N+1)\le3(1+h_B/\ell+\log_\ell h_B)$, hence
$e_A(\ell)\le(3C_A+1)(1+h_B/\ell+\log_\ell h_B)$ and $2e_A(\ell)\le e_P(C,\ell)$ because $2(3C_A+1)\le C$. Multiplying by $2n_\lambda m+1\ge0$
and summing over $\lambda\mid\ell$ gives (R1).

(R2) At an inert $\ell$ in the range of (R2) the canonical table is a rational number chosen strictly between $3C_2\,c_B(\ell,\lambda)$ and
$(3C_2+1)c_B(\ell,\lambda)$, where $c_B(\ell,\lambda)=(n_\lambda+1)r\bigl(h_B/(\ell-1)^{j_A+1}+(G_P+1)(j_A+1)+\Lambda\bigr)>0$; since
$3C_2+1\le C$ it is at most $C\,c_B(\ell,\lambda)$.

(R3), (R4) At the other inert primes the canonical table is a rational number chosen strictly between two explicit expressions
$\mathrm{bd}(K_1,K_2,K_3)<\mathrm{bd}(K_1,K_2+1,K_3)$ built from $h_B,\mathrm{flat},e_P$; the estimate
\lref{GlobalTheory.NLink.RowCDSize.len_le} (deferred) states that the larger one is at most the right-hand side of (R3) when $K_1$, $K_2+1$
and $K_3$ are at most $C$, which gives (R3) and (R4).

(R5) A band pair $(\iota,\lambda)$ comes, by definition of the canonical band predicate, with $\ell\ge11$, a good eigensystem and word
data at $\lambda$ whose box conditions provide $s$ with $w''\le s\le w_B$ and an inequality $|\ell s-b|\le C_B(\cdots)$ in the word
parameters; these parameters are bounded by $b,\eta_{\mathrm{sz}},S_c$, and $\log_\ell(2h)\le\Lambda$ for the relevant $h\le4e^\Lambda$.
An elementary inequality then gives (R5), using $6C_B+1\le C$, $S_c+\eta_{\mathrm{sz}}+t+2\le b$, $w''\ge0$ and $\Lambda\ge1$.
\end{proof}

\subsubsection{Setups}

\begin{lemma}[Killing by the twist]\label{pf:linkN-sel:kill}
Let $D$ be a datum with a norm choice, twisting data and a family $\rho$. Suppose that the norm choice is tied to $D$, that the twists
are good at all twisted places, and that the signature choice condition holds for $\rho$ (for every $\iota$ and good $\pi$ there is a good
signature $s$ of the twist at $w_\iota$ such that finiteness of the Selmer group of the Bloch--Kato reduction of the twist for $s$ implies
finiteness of $\linkNselSel(\iota,\rho,\pi)$ and $\linkNselsP(\iota,\rho,\pi)$ is at most the length of that Selmer group, and
$\sum_{v\in Q_\iota}\linkNsellen(\OO/(\zeta_{\iota,v}-1))$ is at most the tame valuation of the twist at its auxiliary places). Then
$\mathrm{Kill}(\rho)$.
\formallabel{GlobalTheory.T.chosenTwist_annihilates_selmerGroups}
\formallabel{GlobalTheory.NLink.twist_kills_selmer}
\end{lemma}

\begin{proof}
Let $\iota$ be a local index and $w=w_\iota$. Goodness of the twist at $w$ provides nonnegative reals $\eps_T(s)$, indexed by the good
signatures $s$, with $\sum_s\eps_T(s)\le\eps(w)$ and with the length of the Selmer group of the Bloch--Kato reduction for $s$ at most
$\nu+\eps_T(s)$, where $\nu$ is the cyclotomic constant of the twist; it also provides the bound $t\cdot(\text{tame valuation})\le\eps(w)$.
Since the norm choice is tied to $D$, the prime of the twist is $\ell_\iota$, so $\eps(w)=\eps_{\ell_\iota}$ and $\nu=\nu_L(\ell_\iota)$. The
signature choice condition gives, for each $\pi\in\linkNselZg_\iota$, a good signature $s$ for which the Selmer group of $\pi$ is finite and
$\linkNselsP(\iota,\rho,\pi)\le\nu_L(\ell_\iota)+\eps_T(s)\le\nu_L(\ell_\iota)+\eps_{\ell_\iota}$, and it gives
$t\sum_v\linkNsellen(\OO/(\zeta_{\iota,v}-1))\le\eps_{\ell_\iota}$.
\end{proof}

\begin{lemma}[Canonical hub data]\label{pf:linkN-sel:canon}
Assume the degree-one Chebotarev statement. Let $D,L,HW$ and row constants $R$ be given, and suppose that $D$ satisfies the two
conditions of its norm-choice facts which provide Galois representations at the cuspidal eigensystems (the condition on the
attached representations and the forty-first provision). Then for all choices of exponents $e'$, of exceptional sets (as functions of the
representations) and of twisting data $I$, there is a hub datum $X$ which is canonical for $R$, with the prescribed $e'$,
$\linkNselexc(X)$ and $I$, with good $r$ at every local index, with $\eta(X)\le2\eta_{\mathrm{sz}}$ when $\eta_{\mathrm{sz}}\ge0$, and with
$c_B(\iota)$ equal to its canonical value.
\formallabel{GlobalTheory.HubX.exists_canonical_hubData_with_bounds}
\end{lemma}

\begin{proof}
By the lemma \lref{GlobalTheory.PtCore.exists_ptFam} (deferred, like the two other lemmas cited in this proof), which states that Chebotarev, Galois representations at the
cuspidal eigensystems and the existence of a good family at every index give a point family, by
\lref{GlobalTheory.T.galoisRepsCusp_of_adm2}, which derives the Galois representations at the cuspidal eigensystems from the first
condition, and since the forty-first provision provides a good family at every index, there is a point family $(r,\rho)$.
Let $X$ have these representations, the canonical tables and the prescribed $e',\linkNselexc,I$. It is canonical for $R$ by the
verification \lref{GlobalTheory.HubX.canon_ourX}; its $\eta(X)$ is $\lfloor\eta_{\mathrm{sz}}\rfloor\le2\eta_{\mathrm{sz}}$.
\end{proof}

\begin{proposition}[Setups exist]\label{pf:linkN-sel:setup}
Assume $\sfname{Hypotheses}_2$. Let $\mathcal F\in\sfname{FamP}(\eps)$, $C,C_2\in\R$ and $C_T\in\N$, and suppose that $\mathcal F$ satisfies
$\sfname{TwGood}$ and $(\mathrm F)_{C,C_T}$. Then there is $x_0$ such that for all $x\ge x_0$, $x>x_1$, all $i\in I_x$ for which $D_{x,i}$
has finite odd Tate--Shafarevich groups, and every finite set $\linkNselexc$ of primes, there are $HW$, Fontaine data, comparison data and
a hub datum $X$ with $\linkNselexc(X)=\linkNselexc$, the admissibility of the jump parameters of $L_{x,i}$, and a setup
$(L_{x,i},\mathrm{nch}_{x,i},HW,X,C_2)$.
\formallabel{GlobalTheory.T.exists_hubData_admJ_setup}
\formallabel{GlobalTheory.NLink.exists_setup_large_alt}
\formallabel{GlobalTheory.HTW.admissibleParameters_of_large_scale}
\end{proposition}

\begin{proof}
Three thresholds are used: $x_a$ such that every size datum of scale $(x,\eps)$ with $x\ge x_a$ and parameter constant $C_{Lw}$ at most the
corresponding family constant has admissible parameters and satisfies $S_c+\eta_{\mathrm{sz}}+t+2\le b$ (the estimate
\lref{GlobalTheory.HTW.admissible_params_large}, deferred, states exactly this, up to rearranging its inequalities); $x_b$ from
\cref{pf:linkN-sel:small} with $C=1$; and $x_c$ from $\sfname{TwGood}(\mathcal F)$. Let $x_0=\max(x_a,x_b,x_c)$ and fix a member with
finite odd Tate--Shafarevich groups and a set $\linkNselexc$.

The Bhatt--Scholze field of $\sfname{Hypotheses}_2$ gives local data $HW$, Fontaine data and comparison data, by
\lref{GlobalTheory.T.exists_wachFunctor_comparison_of_bhattScholze}, which states that the integral comparison provides them for every
datum. The norm-choice facts contain the hypotheses of \cref{pf:linkN-sel:canon}, which gives $X$ canonical for $R(C_2)$ with the given
$\linkNselexc$ and the twisting data supplied by $\sfname{TwGood}$. \Cref{pf:linkN-sel:small} gives $\Lambda\ge8$ and the split bound. For the
killing condition we apply \cref{pf:linkN-sel:kill}: the norm choice is tied to $D$ and the twists are good by the norm-choice facts and
$\sfname{TwGood}$; the signature choice condition for $X.\rho$ is given by the statement \lref{GlobalTheory.JointSix.main}, which
derives it from Chebotarev, the point family and the norm-choice facts together with the good-point properties of the rows
(\lref{GlobalTheory.Proof.LocalGlobalAssembly.main}, which assembles them from three conditions of the norm-choice facts) and the
absolute irreducibility at good points (\lref{GlobalTheory.PtCore.residualAbsIrredGoodPoint_of_residualIrred}, which deduces it from
residual irreducibility). Finally the admissibility of the jump parameters follows from the norm-choice facts and the admissibility of the
parameters, and all parts of a setup are available.
\end{proof}

% ---------------------------------------------------------------------------
\subsection{The lattice index bound}\label{sec:pf:linkN-sel:pass}

In this subsection $\OO$ is a discrete valuation ring with fraction field $K$ and normalised valuation $v$, $X$ is a finite set, and
$T\subseteq\OO^X$ is an $\OO$-subalgebra which is a local ring and separates the points of $X$. Let $V$ be a finite-dimensional
$K$-vector space with a ring homomorphism $\mathrm{act}\colon T\to\End_K(V)$; for $x\in X$ let $V_x=\bigcap_{f\in T}\ker(\mathrm{act}(f)
-f(x))$ and for $S\subseteq X$ let $V_S=\sum_{x\in S}V_x$. For an $\OO$-submodule $L\subseteq V$ and $S\subseteq X$ the
\emph{congruence module} is
\[
\mathrm{CM}(L,S)=L\big/\bigl((V_S\cap L)+(V_{X\smallsetminus S}\cap L)\bigr).
\]
For $Z\subseteq X$ let $\mathrm{cc}(Z)=\{c\in T:c(y)=0\text{ for all }y\notin Z\}$, an ideal of $T$, and $v_Z(x)=\min_{c\in\mathrm{cc}(Z)}
v(c(x))\in\N\cup\{\infty\}$ (a minimum, as the values lie in a well-ordered set).

We use the following result proved in another section of Part II:
\lref{Formalism.NG.QuotientPowerLengths.congMod_le} states that if $\sum_xV_x=V$, $L$ is finitely generated, $S\subseteq X$ is
finite, $E$ is a set of $K$-endomorphisms of $V$ commuting with $\mathrm{act}(T)$ and mapping $L$ into $L\cap V_S$, and
$w\colon X\to\N\cup\{\infty\}$ is such that for every $x\in S$ some $g\in E$ acts on $V_x$ as a scalar $l\in\OO$ with $v(l)\le w(x)$, then
$\len_\OO\mathrm{CM}(L,S)\le\sum_{x\in S}\dim_KV_x\cdot w(x)$.

\begin{lemma}\label{pf:linkN-sel:etabound}
Assume $\sum_{x\in X}V_x=V$. Let $L\subseteq V$ be a finitely generated $\OO$-submodule which spans $V$ and is stable under
$\mathrm{act}(T)$, let $Z_g\subseteq Z\subseteq X$ be finite, and let $a\in\End_K(V)$ preserve $L$, commute with $\mathrm{act}(T)$, and act
on $V_x$ as multiplication by an integer $A$ for $x\in Z_g$ and as $0$ for $x\in Z\smallsetminus Z_g$. Then for every
$B\colon X\to\N\cup\{\infty\}$ with $v_Z\le B$ on $Z_g$,
\[
\len_\OO\mathrm{CM}(L,Z_g)\le\sum_{x\in Z_g}\dim_KV_x\,\bigl(B(x)+v(A)\bigr).
\]
\formallabel{Formalism.NG.etaBound}
\formallabel{Formalism.NG.congBoundGood}
\end{lemma}

\begin{proof}
It suffices to treat $B=v_Z$. Apply \lref{Formalism.NG.QuotientPowerLengths.congMod_le} with $S=Z_g$,
$E=\{a\circ\mathrm{act}(c):c\in\mathrm{cc}(Z)\}$ and $w(x)=v_Z(x)+v(A)$. Since $T$ is commutative, $a\,\mathrm{act}(c)\,\mathrm{act}(f)=
a\,\mathrm{act}(f)\,\mathrm{act}(c)=\mathrm{act}(f)\,a\,\mathrm{act}(c)$, so $E$ commutes with $\mathrm{act}(T)$, and $E$ preserves $L$ since
$\mathrm{act}(c)$ and $a$ do. For $c\in\mathrm{cc}(Z)$, $\mathrm{act}(c)$ acts on $V_y$ by $c(y)$, which vanishes for $y\notin Z$; as
$V=\sum_yV_y$, $\mathrm{act}(c)(V)\subseteq V_Z$. Next $a(V_Z)\subseteq V_{Z_g}$, because $a$ acts by $A$ on $V_y$ for $y\in Z_g$ and by
$0$ for $y\in Z\smallsetminus Z_g$. Hence every element of $E$ maps $L$ into $V_{Z_g}$. Finally, for $x\in Z_g$ choose $c\in\mathrm{cc}(Z)$
with $v(c(x))=v_Z(x)$; on $V_x$, $a\,\mathrm{act}(c)$ is multiplication by $A\,c(x)\in\OO$, and $v(A\,c(x))=v(A)+v_Z(x)=w(x)$.
\end{proof}

\begin{lemma}\label{pf:linkN-sel:heckelen}
Let $\ell$ be a prime, $M$ a finite free $\OO$-module with an $\OO$-algebra homomorphism $\mathrm{act}\colon T\to\End_\OO(M)$, and identify
$M$ with its image in $V=M\otimes_\OO K$. Let $Z_g\subseteq Z\subseteq X$, $A\in\Z\smallsetminus\{0\}$ and $\alpha\in\End_\OO(M)$ commuting
with $\mathrm{act}(T)$, such that $\alpha\otimes K$ acts on $V_x$ by $A$ for $x\in Z_g$ and by $0$ for $x\in Z\smallsetminus Z_g$. Then
\[
\linkNsellen\Bigl(\frac{(M+V_{X\smallsetminus Z_g})\cap V_{Z_g}}{M\cap V_{Z_g}}\Bigr)\ \le\ \sum_{x\in Z_g}\dim_KV_x
\bigl(\linkNsellen(\OO/\mathfrak c_x)+v_\ell(A)\bigr),\qquad \mathfrak c_x=\{c(x):c\in\mathrm{cc}(Z)\}.
\]
\formallabel{Formalism.MainNotions.TruncatedSeriesCountingIdentities.Br.hecke_length_bound}
\formallabel{Formalism.MainNotions.TruncatedSeriesCountingIdentities.Br.vcc_ne_top}
\end{lemma}

\begin{proof}
If $\ell\in\OO^\times$, then $\len_\OO(\OO/\ell\OO)=0$, all normalised lengths vanish by convention, and the right-hand side is
$\ge0$. Otherwise put $e=v(\ell)=\len_\OO(\OO/\ell\OO)\ge1$. Since $T$ separates points, $\sum_xV_x=V$ by
\lref{Formalism.MainNotions.TruncatedSeriesCountingIdentities.Br.eigSpan} (another section), which deduces this from the existence of
elements of $T$ that are a nonzero constant at one point and vanish elsewhere. The image $L$ of $M$ is finitely generated, spans $V$ and is
$T$-stable. By \lref{Formalism.MainNotions.TruncatedSeriesCountingIdentities.Br.len_le} (another section), which compares the two
modules by means of the projector $C^{-1}\mathrm{act}(\sum_{x\in Z_g}e_x)$ built from such elements $e_x$, the length of the quotient on the
left is at most $\len_\OO\mathrm{CM}(L,Z_g)$. \Cref{pf:linkN-sel:etabound} with $a=\alpha\otimes K$ bounds this by
$\sum_{x\in Z_g}\dim V_x(v_Z(x)+v(A))$. Now $v(A)=v_\ell(A)\,e$: write $A=\ell^{v_\ell(A)}m'$ with $\ell\nmid m'$; as $\ell$ and $m'$ are
coprime in $\Z$ and $\OO$ is local with $\ell\notin\OO^\times$, $m'$ is a unit of $\OO$ (this is the lemma
\lref{Formalism.MainNotions.TruncatedSeriesCountingIdentities.Br.addVal_int}). Moreover $\len_\OO(\OO/\mathfrak c_x)=v_Z(x)$, because
$\mathfrak c_x$ is generated by an element $c(x)$ of minimal valuation, and $v_Z(x)<\infty$ for $x\in Z$, because an element of $T$ which is a
nonzero constant at $x$ and vanishes elsewhere lies in $\mathrm{cc}(Z)$. Dividing by $e$ gives the claim.
\end{proof}

\begin{corollary}\label{pf:linkN-sel:eib}
Let $D$ be a datum and $\iota$ a local index. For every $Z'\subseteq Z_\iota$, every $A\in\Z\smallsetminus\{0\}$ and every
$\alpha\in\End_{\OO_\iota}(M_\iota)$ commuting with $T_\iota$ which acts by $A$ on $V_x$ for $x\in Z'$ and by $0$ on $V_x$ for
$x\in Z_\iota\smallsetminus Z'$,
\[
\linkNsellen\Bigl(\frac{(M_\iota+V_{X_\iota\smallsetminus Z'})\cap V_{Z'}}{M_\iota\cap V_{Z'}}\Bigr)\le\sum_{\pi\in Z'}
\linkNselmult(\iota,\pi)\bigl(\linkNselvcc(\iota,\pi)+v_\ell(A)\bigr).
\]
\formallabel{Formalism.MainNotions.TruncatedSeriesCountingIdentities.Br.hecke_length_bound_localised}
\formallabel{Formalism.MainNotions.TruncatedSeriesCountingIdentities.endomorphism_index_bound}
\formallabel{Formalism.MainNotions.TruncatedSeriesCountingIdentities.second_statement}
\end{corollary}

\begin{proof}
This is \cref{pf:linkN-sel:heckelen} for $\OO_\iota$, $T_\iota$ (local and separating), $M_\iota$, $Z_\iota$ and $Z_g=Z'$, since
$\linkNselvcc(\iota,\pi)=\linkNsellen(\OO/\mathfrak c_\pi)$ and $\linkNselmult(\iota,\pi)=\dim V_\pi$.
\end{proof}

\begin{corollary}[The field \emph{pass}]\label{pf:linkN-sel:pass}
If $D$ satisfies the commuting endomorphism condition, then for every local index $\iota$ with $\ell_\iota\notin\linkNselG$,
$\linkNseleta(\iota)\le\sum_{\pi\in\linkNselZg_\iota}\linkNselmult(\iota,\pi)\bigl(\linkNselvcc(\iota,\pi)+v_\ell(\#G)\bigr)$.
\formallabel{GlobalTheory.EtaPassage.main}
\end{corollary}

\begin{proof}
Take the endomorphism of the commuting endomorphism condition, $A=\#G\ge1$ and $Z'=\linkNselZg_\iota\subseteq Z_\iota$. By the description
of $\linkNseleta(\iota)$ in \cref{pf:linkN-sel:inv}, \cref{pf:linkN-sel:eib} is the claim.
\end{proof}

% ---------------------------------------------------------------------------
\subsection{The bound at the generator primes}\label{sec:pf:linkN-sel:typ}

\begin{definition}[Paired character condition]\label{pf:linkN-sel:pcc}
For $C\in\R$, $\eta\in\N$, a local index $\iota$, a family $\rho$ and $\pi\in X_\iota$ put
\begin{multline*}
\mathrm{rhs}(C,\eta,\iota,\rho,\pi)=4fn_V\,e_T+2fr^2+\linkNselsP(\iota,\rho,\pi)+\sum_{0<e<t}\mathrm{sh}(\ell,e)+v_\ell(h_E)\\
+4n_V\delta_\iota+t\sum_{v\in Q_\iota}\linkNsellen(\OO/(\zeta_{\iota,v}-1)),
\end{multline*}
with $e_T=Cn_V^2(n_V+\eta+t)$. For $C\in\R$ and $C_T\in\N$, the \emph{paired character condition} $\mathrm{PC}(C,C_T)$ says: for every local
index $\iota$ with $\ell_\iota=g_k$, the other generator prime $g_{k'}$, all families $r$ (good) and $\rho$ at $\iota$ with semisimple
decompositions $r_\pi\otimes K\simeq\rho_\pi\oplus\bigoplus_{c=1}^t\chi_{\iota,c}$ for all $\pi\in Z_\iota$, and all natural numbers
$\eta,s,S$, signs and finite sets of integers attached to the places of $E$ such that $\eta+t+1<s$, $\log g_{k'}\le2\log g_k$,
$C_T\bigl((s+S)(n_V+1)^2+(n_V+1)^3\bigr)<g_k$, a sign condition and the host jump condition at $\iota$ hold, the Selmer groups of the
good eigensystems are finite and $\Sha^1(G_{E,S_\ell},\Ql/\Zl(-e))$ is finite for $0<e<t$:
$\linkNseleta(\iota)\le\sum_{\pi\in\linkNselZg_\iota}\linkNselmult(\iota,\pi)\,\mathrm{rhs}(C,\eta,\iota,\rho,\pi)$.
\formal{Formalism.MainNotions.PairedCharacterCondition.PairedCharacterConditionAlt}
\formal{Formalism.MainNotions.PairedCharacterCondition.rhs}
\end{definition}

\begin{lemma}\label{pf:linkN-sel:starHlocal}
Let $C\in\R$ and put $M(\ell)=Cx^C\eta_{\mathrm{sz}}+2\eps_\ell+VV(\ell)+v_\ell(h_E)$. Suppose that $M(\ell)\ge0$ for $\ell\in\linkNselG$, that
the fields \emph{mult} and \emph{rank} hold, and that $\linkNseleta(\iota)\le\sum_{\pi\in\linkNselZg_\iota}\linkNselmult(\iota,\pi)M(\ell_\iota)$
for every $\iota$ with $\ell_\iota\in\linkNselG$. Then the field \emph{typ} holds with the constant $C$.
\formallabel{GlobalTheory.NLinkTyp.starH_of_local}
\end{lemma}

\begin{proof}
Fix $1\le j\le m$ and $\ell\in\linkNselG$. Let $w$ be a twisted place of the support of $L$ above $\ell$. Every $\iota$ with
$(j_\iota,w_\iota)=(j,w)$ has $\ell_\iota=\ell$, so by $\eta_j(w)=\sum_\iota\linkNseleta(\iota)$ and the hypothesis,
\[
\eta_j(w)\le\Bigl(\sum_{\iota:(j_\iota,w_\iota)=(j,w)}\sum_{\pi\in\linkNselZg_\iota}\linkNselmult(\iota,\pi)\Bigr)M(\ell)\le d_{Z_0}(j)\,M(\ell),
\]
by \emph{mult}, \emph{rank} and $M(\ell)\ge0$. Multiply by $\mathrm{lw}(w)e(w)\ge0$ and sum over these $w$. By
\lref{Formalism.MainNotions.TwistedPlaces.admissibleMeanCondition_holds} (another section), which states that
$\sum_{w}\mathrm{lw}(w)e(w)\le\log\ell$ over the twisted places $w$ of the support above $\ell$ (indeed
$\mathrm{lw}(w)e(w)=e(w)f(w)\log\ell/[\mathit{LL}:\Q]$), we get $\sum_w\mathrm{lw}(w)e(w)\eta_j(w)\le\log\ell\cdot d_{Z_0}(j)M(\ell)$.
\end{proof}

\begin{lemma}\label{pf:linkN-sel:rhsle}
Let $X$ be a hub datum with $\mathrm{Kill}(X.\rho)$ and $\eta(X)\le\eta_{\mathrm{sz}}$, let $P$ satisfy the field \emph{par} with a constant
$C_1\ge1$, and let $C'\ge0$. Let $\iota$ be a local index with $\ell=\ell_\iota\notin\linkNselP$ and $\nu_L(\ell)=0$. Then for
$\pi\in\linkNselZg_\iota$, with $C''=2000(C'+1)+5C_1$,
\[
\mathrm{rhs}(C',\eta(X),\iota,X.\rho,\pi)\le C''x^{C''}\eta_{\mathrm{sz}}+2\eps_\ell+VV(\ell)+v_\ell(h_E).
\]
Moreover $\nu_L(g_k)=0$ whenever $g_k\ne2$.
\formallabel{GlobalTheory.NLinkTyp.rhs_le}
\formallabel{GlobalTheory.NLinkTyp.typePrime_size_terms_le_elementary}
\formallabel{UpperTheory.KStar.nuL_eq_zero_of_g}
\end{lemma}

\begin{proof}
The killing condition gives $\linkNselsP(\iota,X.\rho,\pi)\le\nu_L(\ell)+\eps_\ell=\eps_\ell$ and
$t\sum_{v\in Q_\iota}\linkNsellen(\OO/(\zeta_{\iota,v}-1))\le\eps_\ell$. Since $\ell\notin\linkNselP$, $s_0(\iota)=1$, and the lemma
\lref{Formalism.MainNotions.LargeLevelLemmas.vdel_le}, which states $\delta_\iota\le t\,s_0(\iota)+v_\ell(t!)$, gives
$0\le\delta_\iota\le t+v_\ell(t!)\le2t$. It remains to bound $4fn_Ve_T+2fr^2+4n_V\delta_\iota$, which is elementary: with $f=2n$,
$n_V=r+t\le2t$, $\eta(X)\le\eta_{\mathrm{sz}}$, $1\le t\le\eta_{\mathrm{sz}}$, $1\le n$, $n,t\le x^{C_1}$ (from \emph{par}) and $x\ge1$,
$4fn_Ve_T\le4\cdot2n\cdot2t\cdot C'(2t)^2\cdot4\eta_{\mathrm{sz}}\le256C'x^{4C_1}\eta_{\mathrm{sz}}$, $2fr^2\le4x^{3C_1}\eta_{\mathrm{sz}}$ and
$4n_V\delta_\iota\le16x^{2C_1}\eta_{\mathrm{sz}}$, so the sum is at most $(256C'+20)x^{4C_1}\eta_{\mathrm{sz}}\le C''x^{C''}\eta_{\mathrm{sz}}$.
Finally, for a generator prime $g_k$ the invariant $\nu_L(g_k)$ is computed in a field in which $g_k$ splits (a condition of the datum),
and it vanishes by the lemma \lref{UpperTheory.Mu.nu_eq_zero_of_split} (deferred), which states this vanishing for split primes. The
hypothesis $g_k\ne2$ is not used.
\end{proof}

\begin{lemma}\label{pf:linkN-sel:etaLoc}
Assume $\sfname{Hypotheses}_2$. Let $(L,\mathrm{nch},HW,X,C_2)$ be a setup, $C'\in\R$ and $C_T\in\N$ with $C_T\le b$, and assume
$\mathrm{PC}(C',C_T)$ and the size condition on the generator primes ($(kb)^8\le g_k\le2g_{k'}$ and $\log g_k\le x^{\theta}\Lambda$ for
$k,k'\in\{0,1\}$, $\theta$ a size exponent). Then for every $\iota$ with $\ell_\iota=g_k$,
$\linkNseleta(\iota)\le\sum_{\pi\in\linkNselZg_\iota}\linkNselmult(\iota,\pi)\,\mathrm{rhs}(C',\eta(X),\iota,X.\rho,\pi)$.
\formallabel{GlobalTheory.NLinkTyp.etaLoc_le_rhs_alt}
\end{lemma}

\begin{proof}
If $\linkNselZg_\iota=\emptyset$, apply \cref{pf:linkN-sel:eib} with $Z'=\emptyset$, $\alpha=0$ and $A=1$: it gives $\linkNseleta(\iota)\le0$,
which is the claim. Otherwise we use the lemma \lref{GlobalTheory.T.split_bound_of_params_alt}: applied to the canonical hub
datum of the setup (with the Fontaine and finiteness fields of $\sfname{Hypotheses}_2$, Chebotarev, $C_T\le b$, $\Lambda\ge8$ and the
norm-choice facts) it provides, at the generator primes, the goodness of $X.r$, the properties of the eigensystems of $Z_\iota$ from which
the semisimple decompositions follow (by \lref{Formalism.MainNotions.N.semisimple_decomposition_of}, which deduces the decomposition
from a condition of the norm-choice facts and these properties), and integers $s,S$, signs and sets with $\eta(X)+t+1<s$,
$C_T((s+S)(n_V+1)^2+(n_V+1)^3)<g_k$, the sign condition and the host jump condition. From the size condition, $g_{k'}\le2g_k$ and $g_k\ge2$
give $\log g_{k'}\le\log2+\log g_k\le2\log g_k$. The Selmer groups of the good eigensystems are finite by the killing condition of the setup,
and the Tate--Shafarevich groups by the finiteness part of the setup. The paired character condition now gives the claim.
\end{proof}

\begin{proposition}[The field \emph{typ}]\label{pf:linkN-sel:typprop}
Assume $\sfname{Hypotheses}_2$ and a setup $(L,\mathrm{nch},HW,X,C_2)$. Let $C'\ge0$, $C_1\ge1$ and $C_T\in\N$ with $C_T\le b$, and assume
$\mathrm{PC}(C',C_T)$, the field \emph{mult}, the field \emph{par} with constant $C_1$, and the size condition on the generator primes. Then
the field \emph{typ} holds with the constant $2000(C'+1)+5C_1$.
\formallabel{GlobalTheory.T.star_bound_of_paired_character_condition}
\formallabel{GlobalTheory.NLinkTyp.main_alt_alt}
\end{proposition}

\begin{proof}
For the canonical hub datum, $\eta(X)=\lfloor\eta_{\mathrm{sz}}\rfloor\le\eta_{\mathrm{sz}}$ (as $\eta_{\mathrm{sz}}\ge t\ge1$). Let
$C''=2000(C'+1)+5C_1$; then $M(\ell)=C''x^{C''}\eta_{\mathrm{sz}}+2\eps_\ell+VV(\ell)+v_\ell(h_E)\ge0$, every term being nonnegative. For
$\iota$ with $\ell_\iota=g_k$: $g_k\ne2$ by the size condition, $g_k\notin\linkNselP$ by the host conditions, and $\nu_L(g_k)=0$ by
\cref{pf:linkN-sel:rhsle}; so \cref{pf:linkN-sel:etaLoc,pf:linkN-sel:rhsle} give $\linkNseleta(\iota)\le\sum_\pi\linkNselmult(\iota,\pi)
M(g_k)$. The field \emph{rank} is part of the norm-choice facts. Now apply \cref{pf:linkN-sel:starHlocal}.
\end{proof}

% ---------------------------------------------------------------------------
\subsection{The bound at \texorpdfstring{$q$}{q} and \texorpdfstring{$d$}{d}}\label{sec:pf:linkN-sel:thmB}

\subsubsection{Statements used}

A local index $\iota$ is \emph{of type $QD$} if $\ell_\iota\in\{q,d\}$ and $\linkNselZg_\iota\ne\emptyset$. For a place $w$ of $E$, $n_w$
is its local degree.

\begin{lemma}[Local fields at $q$ and $d$]\label{pf:linkN-sel:qdfields}
For every datum $D$: $q\equiv d\equiv3\pmod4$; every place $v$ of $E$ above $\ell\in\{q,d\}$ has ramification index $2$ and is not fixed
by complex conjugation; at $q$ the decomposition group of $v$ preserves the CM type of $E$, and at $d$ every nontrivial inertia element
exchanges the CM type and its complement. Moreover, at every local index with $\ell_\iota\in\{q,d\}$ and every $w\mid\ell_\iota$, the
completion $K_w$ is $K_{w,0}(\sqrt{-\ell_\iota})$ with $[K_w:K_{w,0}]=2$, and $K_w$ lies in an auxiliary field $K_s$ with
$[K_s:\Q_{\ell_\iota}]=\frac{\ell_\iota-1}2[K_w:\Q_{\ell_\iota}]$.
\end{lemma}

This is \lref{GlobalTheory.T.local_fields_at_auxiliary_primes}; its proof (about $560$ lines) is deferred. We also use three
consequences, whose proofs are deferred: \lref{GlobalTheory.T.theoremB_ramifiedPrimes_ne_typePrimes} states that $q,d\notin\linkNselG$;
\lref{GlobalTheory.T.mul_lt_finrank_cyclotomic_at_q_and_d} states that $tr<[E(\zeta_\ell):E]$ for $\ell\in\{q,d\}$; and
\lref{GlobalTheory.T.theoremB_localDegree_sum_bounds} states that every set $s$ of places above $\ell\in\{q,d\}$ satisfies
$\sum_{w\in s}n_w\le2f$ and $2\#s\le\sum_{w\in s}n_w$.

\begin{theorem}[Integrality from a metric]\label{pf:linkN-sel:intmetric}
Let $D$ be a datum with local data $HW$, Fontaine data and comparison data, and $C_A,C_B,C_K,C\in\R$ with $C\ge160$. Assume (F2) at the Wach
structures of all good triples (\cref{pf:linkN-sel:exact}); that $\dim_{\mathit{LL}}D_{\mathrm{dR},\tau}(T)=\rank T$ for every crystalline lattice
$T$ and embedding $\tau$ at the pairs $(K,\mathit{LL})$ admitting a good triple; and that the letter components of every crystalline lattice at a
good triple of level $1$ have full rank. If the thirteen \emph{admitted hypotheses} of the local lifting theorem with constants $(C_A,C_B,C_K)$
hold (among them the natural Herr classes of \cref{pf:linkN-sel:herr} and an Archimedean signature matching with these constants), then the
local theory of $D$ satisfies the conclusion of \cref{pf:linkN-sel:thmA} with the constant $C$.
\formallabel{Formalism.MainNotions.N.Ln34.integrality_from_metric_of_large_constant}
\end{theorem}

\begin{proof}
By \lref{Formalism.MainNotions.N.Ln34.lifting_theorem_of_uniform_link} (deferred; about $16{,}200$ lines, which assemble some thirty
statements on Wach-type lattices, frames, model lattices, exponent recursions and rank bounds), the conclusion holds for all such data as soon as
the \emph{first frame of models} holds: for every datum satisfying the first admitted hypothesis, the dimension and rank assumptions above, the
model statements for lattices and for filtrations imply the first uniform link statement. The first frame of models is
\lref{Formalism.MainNotions.N.Ln34.u1link_pf} (deferred; about $2{,}700$ lines).
\end{proof}

\begin{theorem}[Theorem A: local lifting]\label{pf:linkN-sel:thmA}
Assume the local Bloch--Kato theory, local Tate duality and the local Euler characteristic formula. For every datum $D$ with local data
$HW$, Fontaine data and comparison data, and every $C\ge160$, the local theory of $D$ satisfies the following for every prime $\ell\ne2$
and all good pairs $(K,\mathit{LL})$ of finite extensions of $\Ql$. Let $T_1,T_2$ be crystalline $\OO_{\mathit{LL}}$-lattices of $G_K$ and
$F\colon T_1\to T_2$ a surjective morphism, $r=\rank T_2$, and let $S,\eta,m,h\in\N$, $\mathrm{up}(\tau)\in\{\pm\}$ and $b(\tau)\in\N$
($\tau\colon K\to\mathit{LL}$) be such that the Hodge--Tate jumps of $T_1$ at $\tau$ lie in $[-h,h]\cap\bigl(I_{\mathrm{up}(\tau)}(b(\tau),S)
\cup[-\eta,\eta]\bigr)$, where $I_+(b,S)=[b,b+S]$ and $I_-(b,S)=[-b-S,-b]$, that $\eta<b(\tau)$, that $T_2$ has at most $m$ jumps (with
multiplicity) in $[-\eta,\eta]$ at each $\tau$, and that $[K:K_0]=2$ or $\mathrm{up}(\tau)$ depends only on $\tau|_{K_0}$. Then
\[
\mathrm{col}(F)\le\Bigl(C(n e+1)r\bigl(\tfrac{eh}{\ell}+eS+e\log_\ell h+1+\eta\bigr)+ne(e+1)mC(h+1)\Bigr)\len(\OO_{\mathit{LL}}/\ell),
\]
where $n=[K_0:\Ql]$, $e=[K:K_0]$, and $\mathrm{col}(F)$ is the length of the quotient of the module of finite (crystalline) cocycles of
$T_2$ by the image of the finite cocycles of $T_1$.
\formallabel{GlobalTheory.T.localLifting_of_blochKato_tateDuality_eulerChar}
\formallabel{GlobalTheory.theoremA_closing}
\formallabel{GlobalTheory.T.integralityFromMetric_of_160_le}
\formallabel{GlobalTheory.TheoremAFinal.integralityFromMetric_of_le}
\end{theorem}

\begin{proof}
By \cref{pf:linkN-sel:intmetric} (whose deeper steps, about $20{,}500$ lines through Wach-type lattices, Herr cocycles and universal lift
filtrations, compare \cite{Wach1996,Berger2004,Herr1998,FontaineLaffaille1982}, are deferred), for $C\ge160$ and any constants $(C_A,C_B,C_K)$
the displayed bound holds for every local theory satisfying the admitted hypotheses and three local inputs: exactness of the
Wach functor on short exact sequences of lattices, the dimension of $D_{\mathrm{dR}}$ of crystalline representations, and the rank of the
letter decomposition. The three inputs are \cref{pf:linkN-sel:exact} and the lemmas
\lref{GlobalTheory.CrysDdR.finrank_DdRE_of_isCrys} and \lref{GlobalTheory.NL.IotaLetterRank.finrank_iota_letter} (deferred), and
the admitted hypotheses, with $(C_A,C_B,C_K)=(23,32,32)$, are provided from the three literature statements by the lemma
\lref{GlobalTheory.T.localLifting_theoremA_hypotheses_admitted} (deferred).
\end{proof}

\begin{theorem}[Global lattice index theorem]\label{pf:linkN-sel:glit}
Assume global Artin reciprocity, the local Bloch--Kato theory, local Tate duality and the local Euler characteristic formula. Let $D$ be a
datum satisfying its design conditions. Let $\iota$ be a local index with $\ell=\ell_\iota\notin\linkNselG$, $r$ a good family and $\rho$ a
family at $\iota$, $\pi\in Z_\iota$, and let $\mathrm{dc}$ be a partition of the places of $E$ above $\ell$ into rows and columns, with row
lengths $\kappa_{\mathrm{row}}$, column exponents $e'$ and a block constant $c_B\in\N$ satisfying the twenty \emph{admitted conditions} (among
them: the coefficients of the characteristic polynomials of the $r_x$ lie in $T_\iota$; crystallinity and a gain bound at the places
above $\ell$; absolute irreducibility, a residual decomposition $\bar r_\pi\sim\bar\rho_\pi\oplus\bigoplus_{c=1}^t\bar\chi_{\iota,c}$ and
the semisimple decomposition $r_\pi\otimes K\simeq\rho_\pi\oplus\bigoplus_c\chi_{\iota,c}$; unramifiedness outside $\ell,S_\iota,Q_\iota$ and
conditions at $S_\iota\cup Q_\iota$; the row length bounds; the block condition with constant $c_B$; local Bloch--Kato and Tate duality
statements; and the bound $\linkNsellen(\mathrm{Hom}^{\mathrm{ur}}(G_E,K/\OO))\le v_\ell(h_E)$ for the unramified homomorphisms). If
$\linkNselSel(\iota,\rho,\pi)$ is finite and $\Sha^1(G_{E,S_\ell},\Ql/\Zl(-e))$ is finite for $0<e<t$, then
\begin{multline}\label{eq:pf:linkN-sel:glit}
\linkNselvcc(\iota,\pi)\le\sum_{w}\kappa(w)+\linkNselsP(\iota,\rho,\pi)+\sum_{0<e<t}\mathrm{sh}(\ell,e)+v_\ell(h_E)
+t\sum_{w\in\mathrm{row}}(n_w+2)c_B\\
+4n_V(\#\mathrm{col}+1)\delta_\iota+t\,\#\mathrm{col}\cdot h_0(\mathrm{col})+t\sum_{v\in Q_\iota}\linkNsellen(\OO/(\zeta_{\iota,v}-1)),
\end{multline}
where $\kappa(w)=\kappa_{\mathrm{row}}(w)$ for row places, $\kappa(w)$ is a column length for column places, and $h_0(\mathrm{col})$
is a column invariant. Moreover \cref{pf:linkN-sel:eib} holds for $D$.
\formallabel{GlobalTheory.T.lengthBound_of_rowConditions_and_dualityTheorems}
\formallabel{GlobalTheory.InertiaEigenvalues.LengthBoundClosing.stmt_closed}
\formallabel{GlobalTheory.T.integralityFromMetric_of_designConditionAll}
\formallabel{GlobalTheory.LocalisationBoundAtPoint.boundAtPoint_all_localisations_of_objectConditions}
\end{theorem}

\begin{proof}
The inequality is the theorem \lref{GlobalTheory.T.globalLatticeIndexTheorem_of_lengthBoundAtPlaces} (about $34{,}000$ lines, proved elsewhere in Part II),
which states \eqref{eq:pf:linkN-sel:glit} under the admitted conditions together with four further inputs: the Bloch--Kato annihilator
statements at all places, given by \lref{GlobalTheory.T.blochKatoAnnihilator_all_of_tateDuality_eulerChar} from the three
local literature statements; the finiteness of local cohomology, given by \lref{GlobalTheory.T.tateFinitenessAll_of_localTateDuality} from local Tate duality; the design condition at
$(\iota,r,\rho,\pi,\mathrm{dc})$, which is the hypothesis; and the triviality of unramified finite classes at the column places, given for
every partition by \lref{GlobalTheory.T.unramifiedTrivialH1f_of_decl}. The block condition is passed in the form required
there (it is only needed at row places). The last assertion is \cref{pf:linkN-sel:eib}.
\end{proof}

\begin{lemma}[Row and block bounds at $q$ and $d$]\label{pf:linkN-sel:rowblock}
Let $D$, $HW$, Fontaine data, families $r,\rho$ and size data $L$ be given, and assume \cref{pf:linkN-sel:qdfields}, the conclusion of
\cref{pf:linkN-sel:thmA} with a constant $C_A$, the Hodge jump windows and the Fontaine conditions at $q,d$, the admissibility conditions
(G2) and the Iwahori conditions at the indices of type $QD$, the literature hypotheses at $q,d$, and $S_c+\eta_{\mathrm{sz}}+t+2\le b$. Let
$\iota$ be of type $QD$, $\pi\in\linkNselZg_\iota$ and $\mathrm{dc}$ a partition.
\begin{enumerate}[label=(\alph*)]
\item There are row lengths satisfying the row length bounds at $\pi$ with, for every row place $w$,
$\kappa_{\mathrm{row}}(w)\le\max(C_A,0)\bigl((n_w+1)rY+3n_wm(2b+1)\bigr)+1$, where $Y=4b/\ell+2S_c+2\log_\ell(2b)+1+\eta_{\mathrm{sz}}+t$.
\item If moreover the Wach hypothesis and the local block statement hold, there is $c_B\in\N$ with the block condition at $\pi$ and
$c_B\le60tf(r+t)\log b$.
\end{enumerate}
\end{lemma}

Part (a) is \lref{GlobalTheory.T.rowPlaces_satisfy_rowHypothesis_at_q_and_d} and part (b) is
\lref{GlobalTheory.T.blockBound_qd_places_uniformConst}; their proofs (about $200$ and $250$ lines) are deferred.

\subsubsection{The closing inequality and Theorem B}

\begin{lemma}\label{pf:linkN-sel:closing}
Let $C_A,C,b,\ell,S,\eta,R_w,T_b,V_d\in\R$ and $f,r,t,m,p\in\N$ with $13\max(C_A,0)\le C$, $260f^2t^2(r+t)\le p$, $t,f\ge1$, $m\le r$,
$\ell\ge3$, $S,\eta\ge0$, $b>0$, $\log b>1$, and
\begin{gather*}
R_w\le K\bigl(3fr(\tfrac{4b}\ell+2S+2\log_\ell(2b)+1+\eta+t)+6fm(2b+1)\bigr)+f,\\
T_b\le240t^2f^2(r+t)\log b,\qquad V_d\le4(r+t)t,
\end{gather*}
where $K=\max(C_A,0)$. Then $R_w+T_b+V_d\le Cf\bigl(r(b/\ell+S+t+\eta+\log b)+mb\bigr)+p\log b$.
\formallabel{GlobalTheory.T.closing_numerical_inequality_theoremB}
\formallabel{GlobalTheory.closing_numerical_inequality}
\end{lemma}

\begin{proof}
Write $\lambda=\log_\ell(2b)$ and $\Lambda_b=\log b$. Since $\ell\ge3$, $\log\ell\ge1$, and $\log2<1<\Lambda_b$, so $\lambda\le\log(2b)\le
2\Lambda_b$. Then
\begin{align*}
&13\bigl(r(\tfrac b\ell+S+t+\eta+\Lambda_b)+mb\bigr)-\bigl(3r(\tfrac{4b}\ell+2S+2\lambda+1+\eta+t)+6m(2b+1)\bigr)\\
&\qquad=\tfrac{rb}\ell+7rS+10rt+10r\eta+13r\Lambda_b-6r\lambda-3r+mb-6m\\
&\qquad\ge\tfrac{rb}\ell+7rS+10r\eta+r\Lambda_b+mb+(7r-6m)\ \ge0,
\end{align*}
using $\lambda\le2\Lambda_b$, $t\ge1$ and $m\le r$. Multiplying by $Kf\ge0$ and using $13K\le C$ gives
$R_w\le Cf\bigl(r(\frac b\ell+S+t+\eta+\Lambda_b)+mb\bigr)+f$. Put $Q=f^2t^2(r+t)$; then $f\le Q$, $4(r+t)t\le4Q$ and $Q\le Q\Lambda_b$, so
$f+T_b+V_d\le Q\Lambda_b+240Q\Lambda_b+4Q\Lambda_b\le260Q\Lambda_b\le p\Lambda_b$.
\end{proof}

\begin{definition}[The bound of Theorem B]\label{pf:linkN-sel:TBdef}
For $c\in\R$ and $Q\in\N[X_0,X_1,X_2]$ put $\beta_B(L,c,Q,\ell)=cf\bigl(r(b/\ell+S_c+t+\eta_{\mathrm{sz}}+\log b)+mb\bigr)+Q(f,r,t)\log b$.
For a family $\rho$, $\mathrm{TB}(\rho,L,c,Q)$ is the statement: for every $\iota$ with $\ell=\ell_\iota\in\{q,d\}$ and $x^{2C_t+5}<\ell$,
every $\pi\in\linkNselZg_\iota$ with finite Selmer group, if $\Sha^1(G_{E,S_\ell},\Ql/\Zl(-e))$ is finite for $0<e<t$, then
\begin{equation}\label{eq:pf:linkN-sel:conc2}
\linkNselvcc(\iota,\pi)\le\beta_B(L,c,Q,\ell)+VV(\ell)+v_\ell(h_E)+\linkNselsP(\iota,\rho,\pi)+\eps_\ell .
\end{equation}
\formal{Formalism.MainNotions.EmbeddingTypeJumps.bndB}
\formal{Formalism.MainNotions.EmbeddingTypeJumps.Conclusion2}
\end{definition}

\begin{theorem}[Theorem B from admitted hypotheses]\label{pf:linkN-sel:thmBadm}
Let $D$, $HW$, Fontaine data, families $r,\rho$, size data $L$, $C_A,C\in\R$ and $Q$ be given with $13\max(C_A,0)\le C$ and
$Q(f,r,t)\ge260f^2t^2(r+t)$ for all $f,r,t$. Assume: the Wach hypothesis; the conclusion of \cref{pf:linkN-sel:thmA} with constant $C_A$;
inequality \eqref{eq:pf:linkN-sel:glit} for all admitted data at $(r,\rho)$ at the indices with $\ell_\iota\notin\linkNselG$; at the indices of
type $QD$ the admissibility conditions (G1)--(G6), the Iwahori conditions (continuity, semisimplicity, semisimple decomposition, cyclotomic
inertia image), the block statement, the Bloch--Kato statement, the Hodge jump windows, the Fontaine conditions and the literature
hypotheses; the \emph{killing condition at $q$ and $d$} (if all places above $\ell\in\{q,d\}$ are unramified or $tr<[E(\zeta_\ell):E]$, then
$\linkNselsP(\iota,\rho,\pi)\le\eps_\ell$ and $t\sum_{v\in Q_\iota}\linkNsellen(\OO/(\zeta_{\iota,v}-1))\le\eps_\ell$ for good $\pi$); and
$S_c+\eta_{\mathrm{sz}}+t+2\le b$. Then $\mathrm{TB}(\rho,L,C,Q)$.
\formallabel{GlobalTheory.TheoremB.theoremB_of_admitted}
\end{theorem}

\begin{proof}
Let $\iota$, $\ell\in\{q,d\}$ and $\pi$ be as in $\mathrm{TB}$; then $\iota$ is of type $QD$. Let $\mathrm{dc}$ be the partition with all
places above $\ell$ as rows and no columns. \Cref{pf:linkN-sel:rowblock} gives row lengths and $c_B$ with the stated bounds, and the
hypotheses provide the remaining admitted conditions (the column condition is empty). Since $\ell\notin\linkNselG$,
\eqref{eq:pf:linkN-sel:glit} applies; with no columns it reads
\[
\linkNselvcc\le R_w+\linkNselsP+\sum_{0<e<t}\mathrm{sh}(\ell,e)+v_\ell(h_E)+T_b+V_d+t\sum_{v\in Q_\iota}\linkNsellen(\OO/(\zeta_{\iota,v}-1)),
\]
with $R_w=\sum_{w}\kappa_{\mathrm{row}}(w)$, $T_b=t\sum_w(n_w+2)c_B$ and $V_d=4n_V\delta_\iota$. Let $\Sigma=\sum_wn_w$; then $\Sigma\le2f$ and
$2\#\mathrm{row}\le\Sigma$, hence $\#\mathrm{row}\le f$. With $K=\max(C_A,0)$, $Y$ as in \cref{pf:linkN-sel:rowblock} and $Y'=2b+1$,
\begin{align*}
R_w&\le\sum_w\bigl(K((n_w+1)rY+3n_wmY')+1\bigr)=(KrY+3KmY')\Sigma+\#\mathrm{row}\,(KrY+1)\\
&\le K(3frY+6fmY')+f.
\end{align*}
Next $T_b=t(\Sigma+2\#\mathrm{row})c_B\le t\cdot4f\cdot60tf(r+t)\log b=240t^2f^2(r+t)\log b$. Since $\ell\notin\linkNselP$, $s_0(\iota)=1$,
and $\ell>4tr+1>t$ gives $v_\ell(k)=0$ for $0<k<t$; so $\delta_\iota=t-1\le t$ and $V_d\le4(r+t)t$. As $tr<[E(\zeta_\ell):E]$, the killing
condition at $q,d$ gives $t\sum_v\linkNsellen(\OO/(\zeta_{\iota,v}-1))\le\eps_\ell$. Finally $\ell\ge3$, $\log b=\Lambda>1$, $b>0$, $f,t\ge1$
and $m\le r$, so \cref{pf:linkN-sel:closing} gives $R_w+T_b+V_d\le\beta_B(L,C,Q,\ell)$, which yields \eqref{eq:pf:linkN-sel:conc2}.
\end{proof}

\begin{theorem}[Twisted-place Selmer length bound]\label{pf:linkN-sel:tpb}
Assume global Artin reciprocity, the local Bloch--Kato theory, local Tate duality and the local Euler characteristic formula. Let $D$, $HW$,
Fontaine and comparison data, families $r,\rho$, size data $L$, a sign $\mathrm{ep}$, $C\ge2080$ and $Q$ with $Q(f,r,t)\ge260f^2t^2(r+t)$ be
given, and assume inequality \eqref{eq:pf:linkN-sel:glit} for all admitted data at $(r,\rho)$, the admissibility conditions (G1)--(G6) and
the Iwahori conditions at the indices of type $QD$, the killing condition at $q$ and $d$, the Hodge jump windows (for $\mathrm{ep}$) and the
Fontaine conditions at $q,d$, the tame generator conditions at the indices of type $QD$, and $S_c+\eta_{\mathrm{sz}}+t+2\le b$. Then the
statements of \cref{pf:linkN-sel:qdfields} and $\mathrm{TB}(\rho,L,C,Q)$ hold.
\formallabel{GlobalTheory.T.twistedPlace_selmerLengthBound_of_localDuality}
\formallabel{GlobalTheory.TwistedPlaceBound.closing}
\end{theorem}

\begin{proof}
The first statements are \cref{pf:linkN-sel:qdfields}. For $\mathrm{TB}$ we verify the hypotheses of \cref{pf:linkN-sel:thmBadm} with
$C_A=160$, so that $13\max(C_A,0)=2080\le C$. The Wach hypothesis follows from the statements
\lref{GlobalTheory.T.wachModuleLocalConclusions_usedForms} (the Wach module conclusions in the form used) and
\lref{GlobalTheory.T.wachHypothesis_crystallineLattices} (which deduces the Wach hypothesis for crystalline lattices from them). The
conclusion of \cref{pf:linkN-sel:thmA} holds with $C=160$. The block statement is the theorem
\lref{GlobalTheory.T.theorem_tb_of_local_blochKato_duality_euler} and the Bloch--Kato statement \lref{GlobalTheory.T.blochKatoLocal_all_of_tateDuality_eulerChar}, both derived from the three local literature statements. The
literature hypotheses at an index of type $QD$ consist of: the tame generator conditions (assumed); at every place above $\ell$ and every
coefficient extension, integral Tate duality in degree $2$ (from local Tate duality) and the Bloch--Kato dimension statement at the
auxiliary prime places (\lref{GlobalTheory.T.blochKatoDim_at_auxiliary_prime_places}, deferred); local coefficient
compatibility at all places above $\ell$ (\lref{GlobalTheory.BridgeTie.localCoefficientCompatibility}); the Fontaine properties at all
places, \lref{GlobalTheory.FontaineAdm.fontaineProperties_all_places}; integral Tate duality for $H^2$,
\lref{GlobalTheory.T.integral_tate_H2_top}; and the bound for the unramified homomorphisms by $v_\ell(h_E)$, deduced from
global Artin reciprocity in the lemma \lref{GlobalTheory.T.unramifiedClassFieldBound_of_artin}. The remaining hypotheses are
assumed.
\end{proof}

\begin{theorem}[Theorem B]\label{pf:linkN-sel:thmBmain}
Let $c_1=2080$ and $Q=260X_0^2X_2^2(X_1+X_2)$. Assume $\sfname{Hypotheses}_2$. Then for every datum $D$, size data $L$, parameter choice $P$
and every setup $(L,\mathrm{nch},HW,X,C_2)$ with Fontaine and comparison data, $\mathrm{TB}(X.\rho,L,c_1,Q)$ holds.
\formallabel{GlobalTheory.T.embedding_jumps_aux2_of_setup}
\formallabel{GlobalTheory.TheoremBQD.main_alt_alt}
\formallabel{GlobalTheory.TheoremBQD.main_of_alt}
\end{theorem}

\begin{proof}
We verify the hypotheses of \cref{pf:linkN-sel:tpb} for $(r,\rho)=(X.r,X.\rho)$, $C=c_1$ and $Q$; note $Q(f,r,t)=260f^2t^2(r+t)$. The
lemma \lref{GlobalTheory.T.split_bound_of_params_alt}, applied to the canonical hub datum of the setup (with the Fontaine and
finiteness fields of $\sfname{Hypotheses}_2$, Chebotarev, $\Lambda\ge8$ and the norm-choice facts), gives among others the Fontaine
conditions at $q,d$. The lemma \lref{GlobalTheory.T.set_pt_of_setup_not_twisted} gives at every $\iota$ with
$\ell_\iota\notin\linkNselG$ (in particular at the indices of type $QD$) and every good $\pi$ the admissibility conditions (G1)--(G6), the
Iwahori conditions and the tame generator conditions. The statement \lref{GlobalTheory.JumpStatement.holds_all_alt} gives a sign
$\mathrm{ep}$ for which the lemma \lref{GlobalTheory.T.hodgeJumpParameters_of_setup} (deferred) gives the Hodge jump windows
at $q$ and $d$. \Cref{pf:linkN-sel:glit}, whose design hypothesis holds by the lemma
\lref{GlobalTheory.DesignUnconditional.main} (deferred; it derives the design conditions from the largeness part of the norm-choice facts), gives
\eqref{eq:pf:linkN-sel:glit} for all admitted data. For the killing condition at $q$ and $d$: $\nu_L(\ell)=0$ for $\ell\in\{q,d\}$ by the
lemma \lref{UpperTheory.KStar.nuL_eq_zero_of_qd} (deferred), so the killing condition $\mathrm{Kill}(X.\rho)$ of the setup gives
$\linkNselsP\le\eps_\ell$ and the bound for the elements $\zeta_{\iota,v}-1$. Finally $S_c+\eta_{\mathrm{sz}}+t+2\le b$ is part of the setup.
\end{proof}

\begin{corollary}[The field \emph{ram}]\label{pf:linkN-sel:ram}
Let $c_1\in\R$, $Q$ a polynomial and $C\ge\max(c_1,\Sigma(Q),\deg Q,1)$. If $\mathrm{TB}(\rho,L,c_1,Q)$ and $\mathrm{Kill}(\rho)$ hold, $D$ has
finite odd Tate--Shafarevich groups, and the parameter choice satisfies $x^{2C_t+5}<q$ and $x^{2C_t+5}<d$, then
$\linkNselB(\rho,\iota,\pi;\linkNselbqd(L,C,\ell_\iota))$ for every $\iota$ with $\ell_\iota\in\{q,d\}$ and $\pi\in\linkNselZg_\iota$.
\formallabel{GlobalTheory.NLink.starG_of_thmB}
\end{corollary}

\begin{proof}
Let $\ell=\ell_\iota\in\{q,d\}$; it is an odd prime with $x^{2C_t+5}<\ell$. The Selmer group at $\pi$ is finite by $\mathrm{Kill}(\rho)$ and the
Tate--Shafarevich groups are finite, so \eqref{eq:pf:linkN-sel:conc2} holds. Since $f,r,t\ge1$, every monomial $f^ar^bt^c$ of $Q$ is at most
$(frt)^{\deg Q}$, so $Q(f,r,t)\le\Sigma(Q)(frt)^{\deg Q}\le C(frt)^C$. With $\log b=\Lambda$, $c_1\le C$ and the nonnegativity of
$r(b/\ell+S_c+t+\eta_{\mathrm{sz}}+\Lambda)+mb$ we get $\beta_B(L,c_1,Q,\ell)\le\linkNselbqd(L,C,\ell)$, and \eqref{eq:pf:linkN-sel:conc2}
becomes $\linkNselB(\rho,\iota,\pi;\linkNselbqd(L,C,\ell))$.
\end{proof}

% ---------------------------------------------------------------------------
\subsection{The congruence bound at large primes}\label{sec:pf:linkN-sel:large}

In this subsection lengths may be infinite: for a local index $\iota$ we write $\overline{\mathrm{len}}_\ell(N)=\len_\OO(N)/\len_\OO(\OO/\ell)
\in[0,\infty]$, and $\overline{\mathrm{sh}}_\ell(e)\in\N\cup\{\infty\}$ for the $\Zl$-length of $\Sha^1(G_{E,S_\ell},\Ql/\Zl(-e))$. A continuous
$\rho\colon G_E\to\GL_r(\OO_\iota)$ is a \emph{rank-$r$ part} of $r_\pi$ if $r_\pi\otimes K\simeq\rho\otimes K\oplus\bigoplus_{c=1}^t\chi_{\iota,c}$
as representations of $G_E$.

\begin{definition}[Congruence bound]\label{pf:linkN-sel:congbound}
The \emph{congruence bound} $\mathrm{CB}(D,L)$ means: for every local index $\iota$ with $\ell=\ell_\iota>\linkNsellst$, $\ell\notin\linkNselG$,
$\ell\notin\linkNselExc(L)$, whose eigensystems have the prescribed infinity types (the \emph{weight tie} at $\iota$) and which has a good
eigensystem, for every good family $r$ at $\iota$, every $\pi\in Z_\iota$ and every rank-$r$ part $\rho$ of $r_\pi$:
\[
\overline{\mathrm{len}}_\ell(\OO/\mathfrak c_\pi)\le\overline{\mathrm{len}}_\ell(\linkNselSel(\iota,\rho,\pi))+\sum_{0<e<t}\overline{\mathrm{sh}}_\ell(e)
+v_\ell(h_E)+t\sum_{v\in Q_\iota}\overline{\mathrm{len}}_\ell(\OO/(\zeta_{\iota,v}-1))\quad\text{in }[0,\infty].
\]
\formal{GlobalTheory.IndefiniteRowSideFullHecke.BaseSizes.CongBound_at}
\end{definition}

\begin{lemma}[The field \emph{large}]\label{pf:linkN-sel:large}
Let $D,L$ satisfy $\mathrm{CB}(D,L)$ and the weight parameter condition of the norm-choice facts (which gives the weight tie at every index),
let $(r,\rho)$ be a point family with $\mathrm{Kill}(\rho)$, assume that $D$ has finite odd Tate--Shafarevich groups and $3\le\linkNsellst$,
and let $\linkNselexc=\linkNselExc(L)$. Then $\linkNselB(\rho,\iota,\pi;0)$ holds for all $\iota$ with $\ell_\iota>\linkNsellst$,
$\ell_\iota\notin\linkNselG$, $\ell_\iota\notin\linkNselexc$, and all $\pi\in\linkNselZg_\iota$.
\formallabel{GlobalTheory.T.localBoundZeroAtLargePrimes_of_congruenceBound}
\formallabel{GlobalTheory.NLink.RowE.main}
\end{lemma}

\begin{proof}
Let $\iota,\pi$ be as stated and $\ell=\ell_\iota$. Then $\pi\in Z_\iota$ and $\pi$ is good, the weight tie holds at $\iota$, $r_\iota$ is good and
$\rho_{\iota,\pi}$ is a rank-$r$ part of $r_{\iota,\pi}$ (point family). The congruence bound gives the displayed inequality in $[0,\infty]$; all
terms on the right are finite. Indeed, the Selmer group is finite by $\mathrm{Kill}(\rho)$; since $\ell>\linkNsellst\ge3$, $\ell$ is odd and the
groups $\Sha^1(G_{E,S_\ell},\Ql/\Zl(-e))$, $0<e<t$, are finite; and for $v\in Q_\iota$, $\zeta_{\iota,v}\ne1$ by the tame design of the datum
(the lemma \lref{GlobalTheory.TameDesign.main} states the tame conditions of the datum at the places of $S_\iota\cup Q_\iota$, among them
$\zeta_{\iota,v}\ne1$ for $v\in Q_\iota$), so
$\OO/(\zeta_{\iota,v}-1)$ has finite length. Passing to real numbers: the left side is $\linkNselvcc(\iota,\pi)$, the Selmer term is
$\linkNselsP(\iota,\rho,\pi)$, the Tate--Shafarevich sum is $VV(\ell)$ (the data $\mathrm{sh}(\ell,e)$ equal the finite lengths), the class term
is $v_\ell(h_E)$, and the last term is at most $\eps_\ell$ by $\mathrm{Kill}(\rho)$. This is $\linkNselB(\rho,\iota,\pi;0)$.
\end{proof}

\subsubsection{The local input: Theorem TOP}

Let $K/\Ql$ be a finite extension with maximal unramified subextension $K_0$, $\mathit{LL}/\Ql$ a finite extension with ring of integers $\OO$
and a uniformiser $\lambda$, and $n=[K_0:\Ql]$. Consider a crystalline representation $V$ of $G_K$ over $\mathit{LL}$ with a $G_K$-stable
lattice $T_0$, and continuous characters $\chi_c\colon G_K\to\OO^\times$, $1\le c\le t$. The \emph{position hypothesis} with position data
$(h,w,\eta,\mathrm{typ},U)$ indexed by $\Z/n$ (the embeddings of $K_0$) says: $t\ge1$; the explicit separation and residue-avoidance
\emph{position conditions} hold for $(\ell,n,t,h,w,\eta,\mathrm{typ},U)$; the Hodge--Tate multiplicities of $V$ at the embeddings over the
$i$-th embedding of $K_0$ are given by the multiset $U_i$; and $\chi_c=\chi_1\chi_{\mathrm{cyc}}^{1-c}$, each $\chi_c$ being crystalline with the
single Hodge--Tate weight $w_i+c-1$ over the $i$-th embedding. A cocycle $z\colon G_K\to T\otimes_\OO X$ ($X$ a finite $\OO$-module, $T$ a lattice,
the action twisted by $\chi_1$) is \emph{admissible} if, for all large $a$ with $\lambda^aX=0$, the extension of $G_K$-modules defined by $z$
satisfies modulo $\lambda^a$ the Hodge--Tate and characteristic polynomial conditions determined by the position data and by
$\mathrm{charpoly}(V)\prod_c(Y-\chi_c)$.

\begin{definition}[Admitted clauses of Theorem TOP]\label{pf:linkN-sel:topadm}
For a datum $D$ with local data $HW$, Fontaine data and comparison data, the \emph{admitted clauses of Theorem TOP} are twenty-four statements
on the local theory of $D$: fifteen clauses on Wach lattices, two clauses on $A_{\inf}$, a Herr lattice clause and a clause on lattices; local Tate
duality and the local Euler characteristic formula for all finite extensions of $\Ql$; the Bloch--Kato dimension statement for the jumps of
unramified saturations at the places in case~A; the conclusions of \cref{pf:linkN-sel:crychar}; Sen multiplicities; crystalline literature
statements; and a Frobenius expansion congruence. They are verified in the proof of \cref{pf:linkN-sel:top}.
\formal{Formalism.MainNotions.tableLineDegreeTwo.admitted}
\end{definition}

\begin{theorem}[Core of Theorem TOP]\label{pf:linkN-sel:topcore}
Let $D$ be a datum with local data $HW$, Fontaine data and comparison data satisfying the admitted clauses of Theorem TOP. Then for every prime
$\ell\ge5$, every pair $(K,\mathit{LL})$ in case~A and every uniformiser $\lambda$ of $\OO_{\mathit{LL}}$, the conclusions (i)--(iii) of
\cref{pf:linkN-sel:top} hold at $(K,\mathit{LL},\lambda)$, for the embeddings given by the Wach structure of level $1$ and the inertia group
$I_K$.
\formallabel{GlobalTheory.TopBridge.twistedLatticeCohomology_theoremTOP}
\end{theorem}

\begin{proof}
Fix $\ell$; the period rings of the local theory at $\ell$ are identified with $B_{\mathrm{cris}}$ and $B_{\mathrm{dR}}$. A \emph{situation} for
Theorem TOP consists of $K$, $\mathit{LL}$, $n,r\in\N$, a map $\mathrm{emb}$ from $\Z/n$ to the embeddings of $K_0$ into $\mathit{LL}$, and top data
of size $(n,r)$ (a lattice $T_0$ of rank $r$, characters $\chi_c$, position data and a uniformiser, as in \cref{pf:linkN-sel:top}). A situation is
\emph{realised} if for every residue field $k$ of $\OO_{\mathit{LL}}$ there are a ring with a cyclotomic Frobenius lift on $k((X))$, a functor to
\'etale $\varphi$-modules over it and a model of the top data whose characters reduce to those of the situation, such that $T_0$ maps onto the
model modulo $\lambda$ compatibly with the twist by $\chi_1^{-1}$, and every admitted residual representation has a model. The proof has three
steps; all statements cited are deferred.

\emph{Step 1: every situation is realised.} This is \lref{GlobalTheory.BasePairs.real_of_presented} (about $2{,}800$ lines), which derives
realisation from three inputs: presentation data for the situation, which exist by \lref{GlobalTheory.RedFun.realD} (about $2{,}600$ lines,
through reductions of tensor products and cyclotomic rings); the used properties of the Wach lattices of the Wach structures of level $1$ in
case~A, which follow from the admitted clauses by \lref{GlobalTheory.T.wachLattice_usedProperties_of_admitted} (about $1{,}500$ lines); and the
crystalline calculus of the situation (direct sums, twists and characters have the expected jumps), \lref{GlobalTheory.TopSteps.CrysRules.crysCalc}.

\emph{Step 2: three statements on admissible classes.} From realisation, the admitted local Tate duality and Euler characteristic formula and the
admitted Bloch--Kato dimension statement, \lref{GlobalTheory.TopBridge.stmtAt_of_real} (about $5{,}900$ lines, through bounds for the
residual parts and the Euler characteristic) proves the \emph{error exponent statement}: in every situation satisfying the standing
hypotheses, the twisted lattice has no invariants and vanishing $H^2$, its $H^1$ is free of rank $rn$ with the finite part a direct summand of the
expected rank, and for every finite module $X$ the admissible classes with values in $T_0\otimes X$ are exactly the finite classes. From realisation alone,
\lref{GlobalTheory.TopSteps.CharVanish.stmt_of_real} proves that the admissible classes of the characters $\chi_c$, $2\le c\le t$, with values
in a finite module vanish, and \lref{GlobalTheory.TopSteps.admissibleUntwistedCharUnramified_of_residueRealisation} (about $1{,}050$ lines)
proves that the admissible classes of $\chi_1$ are the classes of homomorphisms trivial on the inertia group.

\emph{Step 3.} \lref{GlobalTheory.T.exactLocalConditions_topCore} (about $2{,}100$ lines, through the comparison of admissibility of extensions
and of their quotients and the transport of cocycles along equivariant isomorphisms) deduces (i)--(iii) at every pair in case~A and every
uniformiser from the three statements of Step 2.
\end{proof}

\begin{theorem}[Theorem TOP]\label{pf:linkN-sel:top}
Assume the local Bloch--Kato theory, local Tate duality and the local Euler characteristic formula. Let $D$ be a datum with local data $HW$,
Fontaine data and comparison data. For every prime $\ell\ge5$, all pairs $(K,\mathit{LL})$ of finite extensions of $\Ql$ singled out by the
local theory of $D$, every uniformiser $\lambda$ of $\OO=\OO_{\mathit{LL}}$, all $(V,T_0,\chi)$ satisfying the position hypothesis and every
finite $\OO$-module $X$:
\begin{enumerate}[label=(\roman*)]
\item a cocycle of $T_0\otimes X$ (twisted by $\chi_1$) is admissible if and only if it lies in the finite part $H^1_f$ of the twisted lattice
$\chi_1^{-1}T_0$;
\item for $2\le c\le t$ and every lattice $L_c$ of $\chi_c$, a cocycle of $L_c\otimes X$ is admissible if and only if it is a coboundary;
\item for every lattice $L_1$ of $\chi_1$, a cocycle of $L_1\otimes X$ is admissible if and only if it is unramified.
\end{enumerate}
\formallabel{GlobalTheory.T.blochKatoTop_of_localTateDuality_eulerChar}
\formallabel{GlobalTheory.TopClosing.blochKatoTop_of_localTateDuality_eulerChar_alt}
\end{theorem}

\begin{proof}
By \cref{pf:linkN-sel:topcore}, the conclusion holds for every datum satisfying the admitted clauses of Theorem TOP
(\cref{pf:linkN-sel:topadm}); note that the pairs singled out by the local theory are the pairs in case~A. The clauses are provided as follows. Fifteen clauses on Wach lattices, and a norm statement, are given by the lemma
\lref{GlobalTheory.T.wachLattice_laterClauses_hold}; two clauses on $A_{\inf}$ follow from the norm statement by the lemma
\lref{Formalism.MainNotions.tableLineDegreeTwo.ainf_clauses_of_theorem}. The Herr lattice clause follows
(\lref{Formalism.MainNotions.InvariantsGenerators.herr_lattice_of_fon}) from the Herr statement
\lref{GlobalTheory.HerrFull.herr_part_a_every_record}, from the invariants statement for fields of norms
\lref{GlobalTheory.T.invariants_statement_holds} (which uses the Fontaine--decompletion statement
\lref{GlobalTheory.Decompl.fontaine_and_decompletion}) and from the statement
\lref{GlobalTheory.T.fontaine_fullFaithful_exact_psi} (full faithfulness and exactness of Fontaine's functor and existence of the
operator $\psi$). The clause on lattices follows (\lref{Formalism.MainNotions.InvariantsGenerators.herr_lattice_rest_iff}) from the
existence of sublattices \lref{GlobalTheory.EtaleSub.exists_lattice_of_pow_mem} and an extension rule
\lref{GlobalTheory.ExtDict.extension_rule}. The remaining clauses are local Tate duality and the Euler characteristic formula (hypotheses),
a Bloch--Kato dimension statement for unramified saturations \lref{GlobalTheory.FiltrationJumpMultiplicity.blochKato_dimension_saturation_unramified},
the existence of crystalline characters with given jumps \lref{GlobalTheory.T.crystalline_characters_given_jumps_unramified_twists}, Sen
multiplicities \lref{GlobalTheory.SenMult.stmtSen}, and crystalline literature statements \lref{GlobalTheory.CrysFormalism.crysLit_all}.
The statements \lref{GlobalTheory.T.fontaine_fullFaithful_exact_psi} and
\lref{GlobalTheory.T.crystalline_characters_given_jumps_unramified_twists} are \cref{pf:linkN-sel:fontaine,pf:linkN-sel:crychar};
\lref{GlobalTheory.T.wachLattice_laterClauses_hold}, \lref{GlobalTheory.HerrFull.herr_part_a_every_record} and
\lref{GlobalTheory.Decompl.fontaine_and_decompletion} are proved elsewhere in Part II; the proofs of the other statements cited in this
proof are deferred.
\end{proof}

\begin{remark}\label{pf:linkN-sel:topremark}
The local Bloch--Kato hypothesis in the statement of \cref{pf:linkN-sel:top} is not used in its proof: the dimension statement that is needed is
proved in the development (\lref{GlobalTheory.FiltrationJumpMultiplicity.blochKato_dimension_saturation_unramified}).
\end{remark}

For a local index $\iota$, a place $w\mid\ell$ of $E$ and a Galois representation $M$ into $M_k(\OO_\iota)$, the \emph{cut cocycles} are the local
cocycles at $w$ of the form $g\mapsto\varphi(gE-E)$, where $Y$ is a finitely generated $G_E$-stable $\OO_\iota$-lattice in a finite sum of the
twisted representations $r_x$ ($x\in X_\iota$), $E$ is a vector with $gE-E\in Y$ for all $g$, and $\varphi\colon Y\to(K/\OO)^k$ is
$G_E$-equivariant for the action through $M$. We write $\mathrm{LB}_w(M,e)$ for the statement that the span of the cut cocycles has length at most
$e\cdot\len(\OO/\ell)$ modulo the finite part $H^1_f$ of the local lattice of $M$.

\begin{lemma}[Cut cocycles after enlarging the coefficients]\label{pf:linkN-sel:enlarge}
Let $D$ and $HW$ be given, $\iota$ a local index with $\ell=\ell_\iota\ge5$, $r$ a good family, $\pi\in Z_\iota$, $\rho$ a family with the
semisimple decomposition at $(\iota,\pi)$ and continuity, and $w\mid\ell$ a place of $E$. Let $\mathit{LL}_C$ be a finite extension of the
coefficient field at $w$ with ring of integers $\OO_C$ such that $(K_w,\mathit{LL}_C)$ is in case~A, $\lambda$ a uniformiser of $\OO_C$,
$\theta\colon G_{K_w}\to\OO_C^\times$ a character and $(h,w',\eta,\mathrm{typ},U)$ position data (we use the notation $H_f$ and $Z^1_{\mathrm{ur}}$ of
\cref{sec:pf:linkN-sel:local}). Assume that the conclusions (i)--(iii) of
\cref{pf:linkN-sel:top} hold at $(K_w,\mathit{LL}_C,\lambda)$, and the \emph{setting condition}: the position hypothesis holds for the
$\theta$-twist of $W_\rho\otimes\OO_C$ and the characters attached to $\theta$, and for every eigensystem $x$ the $\theta$-twist of
$r_x\otimes\OO_C$ satisfies the Hodge--Tate and characteristic polynomial conditions determined by the position data. For a cut cocycle $z$ let
$\tilde z$ be its image under $K/\OO\to\mathit{LL}_C/\OO_C$. Then:
\begin{enumerate}[label=(\alph*)]
\item for every cut cocycle $z$ for $W_\rho$, $\tilde z$ lies in $H_f$ of the lattice $W_\rho\otimes\OO_C$;
\item for $2\le c\le t$ and every cut cocycle $z$ for $\chi_{\iota,c}\chi_{\iota,1}^{-1}$, $\tilde z$ is a coboundary;
\item for every cut cocycle $z$ for the trivial character, $\tilde z$ lies in $Z^1_{\mathrm{ur}}(\OO_C,I_{K_w})$.
\end{enumerate}
\formallabel{GlobalTheory.T.unramifiedEnlargement_cut_cocycles_blochKato}
\formallabel{GlobalTheory.UseLoc.Cases.main}
\formallabel{GlobalTheory.UseLoc.Cases.partW}
\formallabel{GlobalTheory.UseLoc.Cases.partC}
\formallabel{GlobalTheory.UseLoc.Cases.partU}
\end{lemma}

\begin{proof}
The three parts follow the same pattern; we use the following statements, all with deferred proofs.
\begin{itemize}
\item \lref{GlobalTheory.T.cutCocycles_locallyConstant_killed_by_uniformiserPow}: for every cut cocycle $z$ (for any of the representations
above) there is $a\in\N$ such that $\tilde z$ is locally constant with values killed by $\lambda^a$; and $X_a=\lambda^{-a}\OO_C/\OO_C$ is finite
(\lref{GlobalTheory.T.finite_lambdaPowerTorsionFieldModInt}).
\item \lref{GlobalTheory.T.exists_integralTensorTorsion_cocycle_of_torsion_localCocycle} and
\lref{GlobalTheory.T.exists_twistedCharacter_cocycle_of_torsion_localCocycle}: such a $\tilde z$ is the image of a cocycle $z'$ of the
$\theta$-twisted lattice (of $W_\rho\otimes\OO_C$, resp.\ of the character) tensored with $X_a$, the action being twisted by the first character
attached to $\theta$, as in \cref{pf:linkN-sel:top}.
\item \lref{GlobalTheory.T.cocycleExtension_twistedWeilCoefficients_admissible} and
\lref{GlobalTheory.T.cocycleExtension_conductorTwistedCharacter_admissible}: under the setting condition and the hypotheses on
$r,\rho,\pi$, this $z'$ is admissible in the sense of \cref{pf:linkN-sel:top}.
\item \lref{GlobalTheory.T.isBlochKato_read_of_tensor_selmer_add_coboundary},
\lref{GlobalTheory.T.isCoboundary_read_of_twistedCharacter_coboundary} and
\lref{GlobalTheory.T.isUnramified_read_of_trivialTwist_coboundary_on_subgroup}: if $z'$ lies in the finite part, resp.\ is a coboundary,
resp.\ is unramified, then $\tilde z$ lies in $H_f$, resp.\ is a coboundary, resp.\ lies in $Z^1_{\mathrm{ur}}(\,\cdot\,,I_{K_w})$.
\end{itemize}
For (a), the conclusion (i) of \cref{pf:linkN-sel:top} with $X=X_a$, applied to the $\theta$-twist of $W_\rho\otimes\OO_C$ and to the representation
$W_\rho\otimes\OO_C$, which is its twist by the inverse of the first character
(\lref{GlobalTheory.T.representationOfIntegralMatrices_eq_inv_twistedChar_smul}), shows that the admissible cocycle $z'$ lies in the
finite part; hence $\tilde z\in H_f$. For (b), the conclusion (ii) applied to the lattice $\OO_C$ of the $c$-th character shows that $z'$ is a
coboundary, and for (c) the conclusion (iii) applied to the lattice $\OO_C$ of the first character shows that $z'$ is unramified; reading back
gives the claims.
\end{proof}

\begin{proposition}[Local conditions above $\ell$]\label{pf:linkN-sel:useloc}
Let $D$, $HW$ and Fontaine data be given, and assume the conclusion of \cref{pf:linkN-sel:top}, the existence of crystalline characters with given
jumps, and the jump conditions of $D$. Let $\iota$ be a local index with $\ell=\ell_\iota\ge5$, $\ell\notin\linkNselP\cup\{q,d\}$, $r$ a good family,
$\pi\in Z_\iota$ and $\rho$ a family satisfying at $(\iota,\pi)$ the admissibility conditions (G2), (G3), continuity, the semisimple decomposition and
the cyclotomic inertia image. Suppose that for every $n\ge1$, every $\mathrm{up}\colon\Z/n\to\{\pm\}$ and every $\tau\colon\Z/n\to\mathcal T_D$
($\mathcal T_D$ a finite index set of $D$) there are position data satisfying the position conditions for $(\ell,n,t)$ whose multisets describe the
jump sets $J(\mathrm{up},j_\iota,\tau)$ of $D$. Then at every place $w\mid\ell$ of $E$: $\mathrm{LB}_w(W_\rho,0)$ holds for
$W_\rho=\rho_\pi\chi_{\iota,1}^{-1}$ (the cut cocycles for $W_\rho$ lie in $H^1_f$); the cut cocycles for the characters
$\chi_{\iota,c}\chi_{\iota,1}^{-1}$, $2\le c\le t$, are coboundaries; and the cut cocycles for the trivial character are unramified.
\formallabel{GlobalTheory.UseLoc.main}
\end{proposition}

\begin{proof}
We use three descent statements for the passage from $\OO_\iota$ to a coefficient extension $\OO_C$ (deferred):
\lref{GlobalTheory.T.mem_finite_part_of_lift}, \lref{GlobalTheory.T.coboundary_descent_enlargement_coefficients} and
\lref{GlobalTheory.T.unramifiedCocycle_descent_enlargement} state that a cocycle with values in $(K/\OO)^k$ whose image in
$(\mathit{LL}_C/\OO_C)^k$ lies in $H_f$, resp.\ is a coboundary, resp.\ is unramified on a subgroup, has the same property. By
\lref{GlobalTheory.T.exists_exactLocalConditionData_at_place} (deferred), which uses the existence of crystalline characters with given
jumps, the jump conditions, the hypotheses on $r,\rho,\pi$ and the position data, there are $\mathit{LL}_C$, $\lambda$, $\theta$ and position
data at $w$ as in \cref{pf:linkN-sel:enlarge}, satisfying the setting condition, with $(K_w,\mathit{LL}_C)$ in case~A. The conclusion of
\cref{pf:linkN-sel:top} at $(K_w,\mathit{LL}_C,\lambda)$ is a hypothesis. \Cref{pf:linkN-sel:enlarge} and the descent statements give: every cut
cocycle for $W_\rho$ lies in $H^1_f$ of the local lattice of $W_\rho$; the cut cocycles for $\chi_{\iota,c}\chi_{\iota,1}^{-1}$, $2\le c\le t$, are
coboundaries; and the cut cocycles for the trivial character are unramified on $I_{K_w}$, which is the inertia group at $w$
(\lref{GlobalTheory.BridgeInertia.inertEq}). Finally, since all cut cocycles for $W_\rho$ lie in $H^1_f$, so does their span, and the length
in $\mathrm{LB}_w(W_\rho,0)$ is $0$.
\end{proof}

\begin{lemma}[Position data off the exceptional primes]\label{pf:linkN-sel:posdata}
For every family $\mathcal F$ there is $x_0$ such that for all $x\ge x_0$, $x>x_1$, all $i\in I_x$ satisfying the jump conditions, and every prime
$\ell>\linkNsellst(L_{x,i})$: $\ell\ge5$, $t+1<\ell$, $\ell\notin\{q,d\}$ and $\ell\notin\linkNselP$; and if moreover $\ell\notin\linkNselExc(L_{x,i})$,
then for every $1\le j\le m$, $n\ge1$, $\mathrm{up}$ and $\tau$ there are position data satisfying the position conditions for $(\ell,n,t)$ whose
multisets have multiplicities at most one and describe the jump sets $J(\mathrm{up},j,\tau)$.
\formallabel{GlobalTheory.T.exists_localPositionData_off_exceptional_primes}
\formallabel{GlobalTheory.LocalPositionData.position_data_exist}
\end{lemma}

\begin{proof}
The first assertions follow from two numerical lemmas (deferred): $q$, $d$ and the small primes are eventually below $\linkNsellst$
(\lref{GlobalTheory.T.conductorPrimes_lt_specialPrimeThreshold_eventually}), and $\linkNsellst>7$ and $\linkNsellst>t+1$ along the family
(\lref{GlobalTheory.T.specialPrimeThreshold_numerics_along_family}). For the second, let $F_\ell$ and $F_c$ be the integer points of the ``low''
and ``cloud'' parts of the jump set. By the lemma \lref{GlobalTheory.T.lowJumps_and_cloud_in_windows_along_family} (deferred), the
low jumps lie at distance at most $\eta_{\mathrm{sz}}$ and at least $2b/\linkNsellst+t/2+3$ from $(t-1)/2$, and the cloud lies, after the reflection
prescribed by $\mathrm{up}$, between $W_0$ and $W_0+S_c+\eta_{\mathrm{sz}}+1$ for some base size $W_0$. The position data are defined explicitly
from $F_\ell$ and $F_c$; then $h_i\le b+S_c+2\eta_{\mathrm{sz}}+3$, so the carry quantity $\sigma_i$ of the position conditions is at most
$2b/\linkNsellst$ by a geometric series estimate, hence $\sigma_i\le\sigma^+$. Since $\ell\notin\linkNselExc(L)$, $\ell$ divides no integer of the
windows $\mathrm{win}(W_0)$ of \cref{pf:linkN-sel:excdef}; this is the residue avoidance required by the position conditions, and the remaining
separation conditions follow from the window bounds by the lemma \lref{GlobalTheory.T.jump_position_conditions_at_letter}
(deferred). The multiplicity statement holds by construction.
\end{proof}

\subsubsection{The global congruence count}

Let $\iota$ be a local index and $r$ a good family at $\iota$. An \emph{idempotent choice datum} $\mathcal S$ at $\iota$ consists of a finite set
$X'$ with a map $p\colon X'\to X_\iota$ and a point $\pi'\in X'$, a subalgebra $T'\subseteq\OO^{X'}$, idempotent-type matrices, an isomorphism
$K^{n_V}\simeq K^r\times K^t$, and block homomorphisms $\psi_\rho$ (to $M_r(\OO)$) and $\psi_{\chi,c}$ (to $\OO$, $1\le c\le t$). It is an
\emph{orthogonal idempotent decomposition} for $(r,\rho,\pi)$ if, in particular, these define an orthogonal decomposition compatible with
$\rho_\pi\oplus\bigoplus_c\chi_{\iota,c}$ and a certain element $\delta_{\mathcal S}\in\OO$ is a unit. From $\mathcal S$ one builds modules of global
cocycles $\mathrm{Coc}_\rho\subseteq\{G_E\to(K/\OO)^r\}$ and $\mathrm{Coc}_c\subseteq\{G_E\to K/\OO\}$ ($1\le c\le t$), spanned by the functions
$\gamma\mapsto\varphi(y_\gamma)$, where $y_\gamma$ is the matrix-valued function attached to $\gamma$ by $\mathcal S$ and $\varphi$ runs over the
$\OO$-linear maps which are equivariant for $\psi_\rho$, resp.\ $\psi_{\chi,c}$.

\begin{lemma}[Inertia blocks]\label{pf:linkN-sel:quad}
Let $\iota$ be a local index, $r$ a good family, $\rho$ a family, $\pi\in Z_\iota$, and $\mathcal S$ an orthogonal idempotent decomposition for
$(r,\rho,\pi)$. Let $v\in S_\iota\cup Q_\iota$ be a place at which inertia acts through the quadratic relation
$(r_x(\gamma)-1)(r_x(\gamma)-\chi_{\iota,1}(\gamma))=0$ ($x\in X_\iota$, $\gamma\in I_v$). Then for $1\le c\le t$ the module of restrictions of the
elements of $\mathrm{Coc}_c$ to the inertia group $I_v$ has length at most $\len_\OO(\OO/(\zeta_{\iota,v}-1))$.
\formallabel{GlobalTheory.T.inertiaCocycleBlock_length_le_of_quadraticRelation}
\formallabel{GlobalTheory.Cut.QuadN.main}
\formallabel{GlobalTheory.Cut.QuadN.psiChi_group}
\end{lemma}

\begin{proof}
Let $\gamma_0\in I_v$ be the tame generator attached to $(\iota,v)$ by the datum, so that $\zeta_{\iota,v}=\chi_{\iota,1}(\gamma_0)$. Since
$v\nmid\ell$, the cyclotomic character is unramified at $v$, so $\chi_{\iota,c}=\chi_{\iota,1}$ on $I_v$
(\lref{GlobalTheory.T.cyclotomicChar_unramified_away_from_ell}, another section). The compatibility conditions of the orthogonal
decomposition, applied to the vector corresponding to $(0,(1,\dots,1))\in K^r\times K^t$, show that the block homomorphism $\psi_{\chi,c}$
takes the value $\chi_{\iota,c}(g)$ on the element of $\mathcal S$ attached to $g\in G_E$. By \lref{GlobalTheory.Cut.Check.cocW_quad}
(another section), which deduces it from the quadratic relation, every $f\in\mathrm{Coc}_c$ satisfies $\psi_{\chi,c}(\cdot)f(\gamma)=f(\gamma)$
at the elements attached to $\gamma\in I_v$; for $\gamma=\gamma_0$ this reads $\zeta_{\iota,v}f(\gamma_0)=f(\gamma_0)$, i.e.\
$(\zeta_{\iota,v}-1)f(\gamma_0)=0$. On the other hand $f$ is a continuous cocycle with values in the discrete module $K/\OO$ (statement (C) in
the proof of \cref{pf:linkN-sel:gcc}), hence locally constant, and it is additive on $I_v$, because its coefficient character
$\chi_{\iota,c}\chi_{\iota,1}^{-1}$ is trivial on $I_v$. By \lref{GlobalTheory.T.wildInertia_killed_by_additive_locallyConstant} (another
section) such a function vanishes on the wild inertia subgroup (as $v\nmid\ell$), and since $\gamma_0$ topologically generates the tame
quotient, $f|_{I_v}=0$ as soon as $f(\gamma_0)=0$ (\lref{GlobalTheory.Cut.Quad.eq_zero_of_tameGen}). Hence evaluation at $\gamma_0$ is an
injective $\OO$-linear map from the module of restrictions into the $(\zeta_{\iota,v}-1)$-torsion of $K/\OO$, whose length is at most
$\len_\OO(\OO/(\zeta_{\iota,v}-1))$ (\lref{GlobalTheory.Cut.Quad.length_torsionBy_le}).
\end{proof}

\begin{theorem}[Global congruence count]\label{pf:linkN-sel:gcc}
Assume global Artin reciprocity. Let $\iota$ be a local index, $r$ a good family, $\rho$ a family, $\pi\in Z_\iota$, and $\mathcal S$ an
orthogonal idempotent decomposition for $(r,\rho,\pi)$. Assume:
\begin{enumerate}[label=(\alph*)]
\item for $x\notin Z_\iota$ the representation $r_x\otimes\chi_{\iota,1}^{-1}$ has no nonzero invariant vector over $K$, and the $r_x$ are
unramified outside $\ell$, $S_\iota$ and $Q_\iota$;
\item at every $v\in S_\iota$, either inertia acts through the quadratic relation $(r_x(\gamma)-1)(r_x(\gamma)-\chi_{\iota,1}(\gamma))=0$ and
$\zeta_{\iota,v}-1$ is a unit, or a polynomial condition on the action of a tame generator holds; at every $v\in Q_\iota$ the quadratic relation holds
and $\zeta_{\iota,v}\ne1$;
\item at every place $w\mid\ell$: $\mathrm{LB}_w(W_\rho,0)$ with $W_\rho=\rho_\pi\chi_{\iota,1}^{-1}$, the cut cocycles for
$\chi_{\iota,c}\chi_{\iota,1}^{-1}$ ($2\le c\le t$) are coboundaries, and those for the trivial character are unramified.
\end{enumerate}
If $\linkNselSel(\iota,\rho,\pi)$ and $\Sha^1(G_{E,S_\ell},\Ql/\Zl(-e))$ ($0<e<t$) are finite, then
\[
\linkNselvcc(\iota,\pi)\le\linkNselsP(\iota,\rho,\pi)+\sum_{0<e<t}\mathrm{sh}(\ell,e)+v_\ell(h_E)+t\sum_{v\in Q_\iota}\linkNsellen(\OO/(\zeta_{\iota,v}-1)).
\]
\formallabel{GlobalTheory.T.globalCongruenceCount_of_rootDifference_unit}
\formallabel{GlobalTheory.CutCoreUnit.main}
\end{theorem}

\begin{proof}
Let $L_0=\len_\OO(\OO/\ell\OO)$, a positive integer since $\ell\ne0$ is not a unit in $\OO$; let $s=\len_\OO\linkNselSel(\iota,\rho,\pi)<\infty$,
$h=v_\ell(h_E)$, $g_e=\overline{\mathrm{sh}}_\ell(e)<\infty$, and $z_v=\len_\OO(\OO/(\zeta_{\iota,v}-1))$ for $v\in Q_\iota$ (finite since $\zeta_{\iota,v}\ne1$).
We use the following statements; the proofs of (D), (B), (C) and of the second half of (I) are deferred:
\begin{itemize}
\item[(D)] the \emph{determinant bound} \lref{GlobalTheory.T.length_heckeValueIdeal_le_determinant_bound}: under the first condition of (a),
$\len_\OO(\OO/\mathfrak c'_\pi)\le\len\mathrm{Coc}_\rho+\sum_{c=1}^t\len\mathrm{Coc}_c$, where $\mathfrak c'_\pi$ is the ideal generated by the
values at $\pi$ of the elements of $T_\iota$ vanishing at the eigensystems outside $Z_\iota$; and $\mathfrak c'_\pi=\mathfrak c_\pi$ by the
definitions;
\item[(B)] the \emph{coboundary bound} \lref{GlobalTheory.T.length_cutCoboundaryModule_le_unitDelta}: the coboundary parts of $\mathrm{Coc}_\rho$
and $\mathrm{Coc}_c$ have length at most a multiple of $\len(\OO/\delta_{\mathcal S})=0$;
\item[(C)] \lref{GlobalTheory.T.finiteQuotientCocycleSpans_le_continuousCocycles} and
\lref{GlobalTheory.T.localCocycles_aboveEll_inCut_unitDelta}: the elements of $\mathrm{Coc}_\rho$ and $\mathrm{Coc}_c$ are continuous cocycles,
and their restrictions to the places above $\ell$ are cut cocycles;
\item[(I)] \cref{pf:linkN-sel:quad}: at $v\in S_\iota\cup Q_\iota$ with the quadratic relation, the module of inertia restrictions of $\mathrm{Coc}_c$
has length at most $\len(\OO/(\zeta_{\iota,v}-1))$; and \lref{GlobalTheory.T.inertia_restriction_polynomial_condition_second_form}: under the
polynomial condition the inertia restrictions vanish.
\end{itemize}
We also use results of another section of Part II: \lref{GlobalTheory.T.length_global_cocycles_le_sum_places_over} (a module of global cocycles
whose restrictions above $\ell$ are cut cocycles with $\mathrm{LB}_w(\cdot,e_w)$ contains a submodule of cocycles satisfying the local conditions,
of colength at most $\sum_we_wL_0$); \lref{GlobalTheory.T.length_cocycles_le_coboundaries_add_selmer} (a module of continuous cocycles for
$W_\rho$ satisfying the local conditions at $\ell$, with vanishing inertia restrictions outside $\ell,S_\iota,Q_\iota$, has length at most the length of
its coboundary part plus $s$); \lref{GlobalTheory.Proof.LocalGlobalSelmer.main} (under the second condition of (a), the inertia restrictions
outside $\ell,S_\iota,Q_\iota$ vanish); \lref{GlobalTheory.Cut.Peel.main} (a module of cocycles contains a submodule with trivial inertia
restrictions at the places of $Q_\iota$, of colength at most the sum over $v\in Q_\iota$ of the lengths of the modules of inertia restrictions at $v$);
\lref{GlobalTheory.UnramHom.length_le_classNumber} (by global Artin reciprocity, a module of everywhere unramified homomorphisms
$G_E\to K/\OO$ has length at most $L_0h$); \lref{GlobalTheory.CutTrivIne.main} (unramified cut cocycles at $w\mid\ell$ have trivial inertia
restrictions); and \lref{GlobalTheory.T.length_cyclotomic_block_cocycles_le_sha} (a module of cocycles for $\chi_{\mathrm{cyc}}^{1-c}$, $c\ge2$,
which are coboundaries at the places above $\ell$ and unramified elsewhere, has length at most its coboundary part plus $L_0\,g_{c-1}$).

\emph{The $\rho$-block.} By (C) and (c), \lref{GlobalTheory.T.length_global_cocycles_le_sum_places_over} with $e_w=0$ shows that
$\mathrm{Coc}_\rho$ satisfies the local conditions at $\ell$; its inertia restrictions outside $\ell,S_\iota,Q_\iota$ vanish by (a); hence
$\len\mathrm{Coc}_\rho\le(\text{coboundary part})+s=s$ by \lref{GlobalTheory.T.length_cocycles_le_coboundaries_add_selmer} and (B).

\emph{The $\chi$-blocks.} Let $1\le c\le t$; the elements of $\mathrm{Coc}_c$ are cocycles for $\chi_{\iota,c}\chi_{\iota,1}^{-1}=\chi_{\mathrm{cyc}}^{1-c}$.
By \lref{GlobalTheory.Cut.Peel.main} and (I) at the places of $Q_\iota$ there is $C'\subseteq\mathrm{Coc}_c$ of colength $\le\sum_{v\in Q_\iota}z_v$
with trivial inertia restrictions at $Q_\iota$. At $v\in S_\iota$ the inertia restrictions of $C'$ vanish by (I) (in the quadratic case
$\len(\OO/(\zeta_{\iota,v}-1))=0$ since $\zeta_{\iota,v}-1$ is a unit), and outside $\ell,S_\iota,Q_\iota$ by (a). If $c=1$, the elements of $C'$ are
homomorphisms, unramified at the places above $\ell$ by (c) and \lref{GlobalTheory.CutTrivIne.main}; so $\len C'\le L_0h$. If $c\ge2$, they are
coboundaries at the places above $\ell$ by (c), and $\len C'\le L_0g_{c-1}$ by \lref{GlobalTheory.T.length_cyclotomic_block_cocycles_le_sha}
and (B). Hence $\len\mathrm{Coc}_1\le L_0h+\sum_vz_v$ and $\len\mathrm{Coc}_c\le L_0g_{c-1}+\sum_vz_v$ for $c\ge2$.

\emph{Conclusion.} By (D),
\[
\len_\OO(\OO/\mathfrak c_\pi)\le s+L_0\Bigl(h+\sum_{0<e<t}g_e\Bigr)+t\sum_{v\in Q_\iota}z_v .
\]
Divide by $L_0$: since $g_e=\mathrm{sh}(\ell,e)$, $s/L_0=\linkNselsP(\iota,\rho,\pi)$ and $z_v/L_0=\linkNsellen(\OO/(\zeta_{\iota,v}-1))$, this is the
claim.
\end{proof}

\begin{lemma}\label{pf:linkN-sel:extended}
Let $\iota$, $\pi\in Z_\iota$ and $\rho$ be such that $\zeta_{\iota,v}\ne1$ for $v\in Q_\iota$, and suppose that the conclusion of
\cref{pf:linkN-sel:gcc} holds whenever the Selmer group and the Tate--Shafarevich groups are finite. Then the inequality of
\cref{pf:linkN-sel:congbound} holds in $[0,\infty]$.
\formallabel{GlobalTheory.T.congruenceBound_extended_of_rationalBound}
\formallabel{GlobalTheory.Proof.congruenceBound_extended_of_rationalBound.congBound_of_rational}
\end{lemma}

\begin{proof}
If the Selmer group or some $\Sha^1(G_{E,S_\ell},\Ql/\Zl(-e))$ is infinite, the right-hand side is $\infty$. Otherwise all terms are finite
($\OO/\mathfrak c_\pi$ has finite length since $\pi\in Z_\iota$, and $\OO/(\zeta_{\iota,v}-1)$ since $\zeta_{\iota,v}\ne1$), they coincide with the
normalised rational lengths (the elements $\zeta_{\iota,v}$ entering the two formulations agree), and the claim is the rational inequality.
\end{proof}

\begin{theorem}[Congruence bound along families]\label{pf:linkN-sel:congfam}
Assume $\sfname{Hypotheses}_2$. For every family $\mathcal F$ there is $x_E$ such that $\mathrm{CB}(D_{x,i},L_{x,i})$ holds for all $x\ge x_E$ with
$x>x_1$ and all $i\in I_x$ for which the norm-choice facts hold for $(D_{x,i},\mathrm{nch}_{x,i})$.
\formallabel{GlobalTheory.T.eventually_congBound_of_facts}
\formallabel{GlobalTheory.NLinkE.main_alt_alt}
\end{theorem}

\begin{proof}
Let $x_E$ be the maximum of the threshold of \cref{pf:linkN-sel:posdata}, of the threshold of the lemma
\lref{GlobalTheory.T.eventually_prvJ_of_facts} (deferred; the jump conditions hold eventually for members with the norm-choice facts), and of a
threshold beyond which the parameters are admissible. Fix a member with the norm-choice facts, and $\iota,r,\pi,\rho$ as in
\cref{pf:linkN-sel:congbound}; let $\ell=\ell_\iota$.
\begin{enumerate}[label=\textit{Step \arabic*.},leftmargin=*]
\item By \cref{pf:linkN-sel:posdata}, $\ell\ge5$, $t+1<\ell$, $\ell\notin\{q,d\}\cup\linkNselP$, and (as $\ell\notin\linkNselExc(L)$) position data
exist for every $(n,\mathrm{up},\tau)$ and the lift index $j_\iota$. Hence $E$ is unramified at $\ell$ (every prime outside $\linkNselP\cup\{q,d\}$
is unramified in $E$, a property of the host), and the cyclotomic inertia image holds at $\iota$ (a property of every local index).
\item From the norm-choice facts at the large index $\iota$ with a good eigensystem, six lemmas (deferred) give: (G1), continuity and
semisimplicity of the $r_x$ (\lref{GlobalTheory.T.admG1_continuous_admSS_of_rank_r_summand}); (G3)
(\lref{GlobalTheory.T.admG3_of_isRankR_of_lStar_lt}); (G2), using Fontaine's period ring and the admissibility of the parameters
(\lref{GlobalTheory.T.admG2_of_isRankR_of_lStar_lt}); condition (a) of \cref{pf:linkN-sel:gcc}
(\lref{GlobalTheory.T.noChiOff_and_unrOut_of_large_good_place}); and condition (b) at $S_\iota$ and at $Q_\iota$
(\lref{GlobalTheory.T.quadAt_unit_zeta_sub_one_or_polyAt}, \lref{GlobalTheory.T.quadAt_and_zeta_ne_one_of_large_place}). The
semisimple decomposition at $\pi$ is the hypothesis that $\rho$ is a rank-$r$ part of $r_\pi$.
\item An orthogonal idempotent decomposition $\mathcal S$ for $(r,\rho,\pi)$ exists by the lemma
\lref{GlobalTheory.T.exists_globalStep_objects_of_cyclotomic_generator_inertia} (deferred), which constructs it from Steps 1--2, $t+1<\ell$ and
the unramifiedness of $E$ at $\ell$.
\item The Bhatt--Scholze field of $\sfname{Hypotheses}_2$ gives local data $HW$, Fontaine data and comparison data
(\lref{GlobalTheory.T.exists_wachFunctor_comparison_of_bhattScholze}), and \cref{pf:linkN-sel:top} applies to them.
\Cref{pf:linkN-sel:useloc}, with the crystalline characters with given jumps, the jump conditions, Steps 1--2 and the position data, gives condition
(c) of \cref{pf:linkN-sel:gcc} at every $w\mid\ell$.
\item \Cref{pf:linkN-sel:gcc} (with global Artin reciprocity) gives the rational inequality whenever the Selmer and Tate--Shafarevich groups are
finite, and \cref{pf:linkN-sel:extended} (with $\zeta_{\iota,v}\ne1$ for $v\in Q_\iota$, Step 2) gives the inequality of $\mathrm{CB}(D,L)$.\qedhere
\end{enumerate}
\end{proof}

% ---------------------------------------------------------------------------
\subsection{Theorem E and Corollary A}\label{sec:pf:linkN-sel:thmE}

The middle-range bound \lref{GlobalTheory.T.star_ad_and_bb_le_of_joint_cond} (another section) uses a lifting corollary for the local theory
of the datum with a constant $C_2\ge C_0$. We prove the existence of $C_0$.

\begin{definition}[Lifting theorem and lifting corollary]\label{pf:linkN-sel:liftdef}
Let $D$ be a datum. For finite extensions $K,\mathit{LL}$ of $\Ql$, a surjective morphism $F\colon T_1\to T_2$ of crystalline
$\OO_{\mathit{LL}}$-lattices of $G_K$ and $e\in\N$, we say that $F$ \emph{lifts with exponent} $e$ if $\ell^e c$ lies in the image of the finite
cocycles of $T_1$ for every finite cocycle $c$ of $T_2$. The local theory of $D$ satisfies the \emph{lifting theorem with constant $C$} if:
(i) for every $\ell\ge5$, all $K,\mathit{LL}$, every $n\ge1$ and every map $\mathrm{emb}$ from $\Z/n$ to the embeddings of $K_0$ into
$\mathit{LL}$ satisfying the letter condition, and every such $F$ with good coefficients whose source has its Hodge--Tate jumps in windows of
prescribed kinds and sizes (parameters $j,G,S,\eta,h$, types, cuts and arc lengths subject to explicit size, cut and word conditions), $F$ lifts with
an explicit exponent $e=e(\ell,n,j,g,h,\dots)$, where $g=\max(G,S,\eta,1)$; and (ii) under the same size conditions
$e\le C\bigl(h/(\ell-1)^{j+1}+g\,\ell_{\max}+\log_\ell h\bigr)$, $\ell_{\max}$ being the maximal arc length of the cuts. It satisfies the
\emph{lifting corollary with constant $C$} if for every $\ell\ge11$, all $K,\mathit{LL},n,\mathrm{emb}$ as above, and every such $F$ with
$\rank T_2=r$, $1\le j$, $j+1\le n$, the jumps of $T_1$ over the $i$-th embedding in an upper or lower window of size $\mathrm{sz}(i)$ with
parameters $(S,\eta)$, every $i$ either thick (of size near $b$) or of size $\le w$, $C(j+1)(w+S+\eta+1)\le b$, and the thick indices forming
clusters of length $\le j$:
\[
\mathrm{col}(F)\le C(n+1)r\Bigl(\frac{h_{\mathrm{cor}}}{(\ell-1)^{j+1}}+(w+S+\eta+1)(j+1)+\log_\ell h_{\mathrm{cor}}\Bigr)\len(\OO_{\mathit{LL}}/\ell),
\]
where $h_{\mathrm{cor}}=b+\eta+S$ if some $\mathrm{sz}(i)$ lies in $[b,b+\eta]$ and $h_{\mathrm{cor}}=w+S+\eta$ otherwise.
\formal{Formalism.MainNotions.WachCoordinates.LiftingTheoremLocalTheory}
\formal{Formalism.MainNotions.WachCoordinates.ThmELiftAt}
\formal{Formalism.MainNotions.WachCoordinates.ThmEBound}
\formal{Formalism.MainNotions.WachCoordinates.LocalCorollaryLifting}
\formal{Formalism.MainNotions.WachCoordinates.CorAAt}
\end{definition}

\begin{theorem}[Theorem E]\label{pf:linkN-sel:thmEcore}
Let $D$ be a datum with local data $HW$, Fontaine data and comparison data, and $C\in\R$. Assume: the clauses (P3), (P4), (P5), (P6), (P9) and two
further conditions on $\Delta$-invariant lattices; the admitted local hypotheses of Theorem E (the conditions on $\delta$-invariant lattices, the
Archimedean signature matching with constant $C$, and the Hodge, filtration and saturation hypotheses at the places in case~A with $\ell\ge5$);
$D_{\mathrm{dR}}=D_{\mathrm{cris}}$ at the pairs in case~A; Fontaine's theory (\cref{pf:linkN-sel:fontaine}); and, for $\ell\ge5$ and every pair in
case~A, the existence of a unit generator of the Wach module of the trivial lattice. Then the local theory of $D$ satisfies the lifting theorem
with constant $800$.
\formallabel{Formalism.MainNotions.N.Ln03.local_theorem_of_inputs}
\end{theorem}

\begin{proof}
All statements cited in this proof are deferred, except \lref{Formalism.MainNotions.LiftingTheory.gen_one_crys}.

\emph{Reduction.} By \lref{Formalism.MainNotions.N.Ln03.local_theorem_of_cyclotomic_lift} (about $700$ lines), together with the base change of
the local theory along coefficient extensions (\lref{Formalism.MainNotions.UnusedSlot80.base_change_every_theory}), the lifting theorem with
constant $800$ holds as soon as the lifting (i) of \cref{pf:linkN-sel:liftdef} holds, for the embeddings of the Wach structure $W$ of level $1$, at
every place $(\ell,K,\mathit{LL})$ in case~A with $\ell\ge5$ for which $\mathit{LL}$ contains a primitive $\ell^{[K_0:\Ql]}$-th root of unity. Its proof
passes to such coefficient fields by base change and proves the bound (ii) with constant $800$ for the explicit exponent.

\emph{The lifting at $W$.} Fix such a place. By \lref{Formalism.MainNotions.WachCoordinates.Br.crystalline_lifting_second_form}
(about $9{,}800$ lines, through the primitive and final forms of Theorem E, arc lifting and realisation in coordinates) the lifting (i) holds
at the place if the hypotheses of the second form hold and there is a unit generator $e$ of $D(\mathbf 1)$ satisfying the Bloch--Kato extension property
and the realisation property. The hypotheses of the second form are assembled from (P5), (P6), (P9), the further lattice conditions and the
admitted hypotheses (\lref{Formalism.MainNotions.WachCoordinates.Br.brHyp2_of_rec}). A unit generator $e$ exists by assumption; it has the
Bloch--Kato extension property by \lref{UpperTheory.Prune.bkExt_holds}, from the admitted saturation hypothesis and Fontaine's theory. For
the realisation property we apply \lref{Formalism.MainNotions.N.Ln03.realised2_pf} (about $8{,}700$ lines, through level kernel detection
and the coordinate, splitting and coefficient steps). Its hypotheses are: the compatibilities of the comparison data with the levels, with
$\varphi$ and with $\Gamma$ (fields of the comparison data); the admitted hypotheses; (P3), (P4), (P5), (P9) and the further lattice condition;
$D_{\mathrm{dR}}=D_{\mathrm{cris}}$; the crystallinity of the generator of the trivial lattice
(\lref{Formalism.MainNotions.LiftingTheory.gen_one_crys}); the fact that $\OO_{\mathit{LL}}$ is a complete discrete valuation ring in which
$\ell$ is not a unit (\lref{Formalism.MainNotions.N.Ln03.Pf.complMax_holds}); a unit relating the generator to the cyclotomic data
(\lref{Formalism.MainNotions.N.Ln03.Pf.uTieRec_holds}); that $K$ is unramified; the primitive root of unity; and $\ell\ge5$.
\end{proof}

\begin{theorem}[Corollary of Theorem E]\label{pf:linkN-sel:corE}
There is $C'\ge1$ such that for every datum $D$ and every $C_1\ge0$: if the local theory of $D$ satisfies the lifting theorem with constant $C_1$,
the generator bound for $H^1_f$ at all places, and the jump statement for surjections (if $F\colon T_1\to T_2$ is a surjection of crystalline
lattices and $A$ is a set of integers containing the jumps of $T_1$ at a letter, then $A$ contains the jumps of $T_2$ there), then it satisfies
the lifting corollary with constant $\max(C',4C_1)$.
\formallabel{Formalism.MainNotions.WachCoordinates.corollary_of_Theorem_E}
\end{theorem}

\begin{proof}
By \lref{Formalism.MainNotions.WachCoordinates.aprimeConfig_holds} (deferred; about $740$ lines, through the existence of cuts with
exponent bounds for words with sparse thick letters) there is $C'\ge1$ satisfying the \emph{cut configuration statement}: for all
$\ell,n,j,b,w,S,\eta$, types and sizes with $j\ge1$, $j+1\le n$, the window and thickness conditions, $C'(j+1)(w+S+\eta+1)\le b$ and clustered
thick letters, there are cuts whose exponents satisfy the bounds needed below. Fix $D$, $C_1$, a place and $n,\mathrm{emb}$. The lemma
\lref{Formalism.MainNotions.WachCoordinates.corAAt_of} (deferred; about $430$ lines) states that the cut configuration statement for $C'$,
the bound (ii) of the lifting theorem with $C_1$, the lifting (i) at the place, the jump statement, the generator bound and the finiteness of
$\len(\OO/\ell\OO)$ (\lref{Formalism.MainNotions.LocalTheoryWindows.lenOl_finite}, deferred) give the lifting corollary at the place with
constant $\max(C',4C_1)$; its proof obtains a lifting exponent from (i) and the cuts, and bounds $\mathrm{col}(F)$ by the number of generators of
$H^1_f(T_2)$ times this exponent.
\end{proof}

\begin{theorem}[Theorem E and Corollary A]\label{pf:linkN-sel:thmE}
Assume the local Bloch--Kato theory, local Tate duality and the local Euler characteristic formula. There is $C_0\in\R$ such that for every datum
$D$ with local data $HW$, Fontaine data and comparison data, the local theory of $D$ satisfies the lifting theorem with constant $800$ and, for every
$C_2\ge C_0$, the lifting corollary with constant $C_2$.
\formallabel{GlobalTheory.T.theoremE_corollaryA_closing_all_data}
\formallabel{GlobalTheory.ThmECorA.closing}
\formallabel{Formalism.MainNotions.N.Ln03.corollary_local_exists_constant}
\formallabel{Formalism.MainNotions.N.Ln03.local_corollary_of_de_rham}
\end{theorem}

\begin{proof}
Let $C'$ be the constant of \cref{pf:linkN-sel:corE} and put $C_0=\max(C',3200)$.

Fix $D$ with local and comparison data. The following inputs hold: the equality $D_{\mathrm{dR}}=D_{\mathrm{cris}}$ for the pairs $(K,\mathit{LL})$
singled out by the local theory (\lref{GlobalTheory.Per.DdREqDcris.record_caseA}); the Fontaine and decompletion statements
(\lref{GlobalTheory.Decompl.fontaine_and_decompletion}); Fontaine's theory (\cref{pf:linkN-sel:fontaine}); the existence of unit generators of the Wach modules for $\ell\ge5$
(\lref{GlobalTheory.WachUnit.unitGen_exists}); and sixteen admitted local hypotheses, derived from the three literature statements by
\lref{GlobalTheory.T.localLifting_theoremE_hypotheses_admitted} (the proofs of the inputs other than \cref{pf:linkN-sel:fontaine} and the Fontaine and decompletion statements are deferred). From the admitted
hypotheses one extracts the conditions on $\delta$-invariant lattices, the generator bound for $H^1_f$ and the Hodge exactness of all
extensions. Then \cref{pf:linkN-sel:thmEcore} gives the lifting theorem with constant $800$.
Applying \cref{pf:linkN-sel:corE} with $C_1=800$, the generator bound, and the jump statement
deduced from Hodge exactness (if $F$ is onto and $T_1,T_2$ are crystalline, every jump of $T_2$ is a jump of $T_1$), we obtain the lifting corollary
with constant $\max(C',3200)=C_0$; by \lref{Formalism.MainNotions.N.Ln03.corollary_monotone_in_constant} (the
corollary with constant $C$ implies the corollary with every larger constant) it holds for every $C_2\ge C_0$.
\end{proof}

% ---------------------------------------------------------------------------
\subsection{Joint hub conditions at inert primes}\label{sec:pf:linkN-sel:bcd}

\begin{definition}[Joint hub conditions]\label{pf:linkN-sel:jointdef}
Let $R=(C_2,\dots)$ be row constants and $X$ a hub datum. The \emph{joint hub conditions} for $(L,HW,X,R)$ require, for every local index $\iota$ with
$\ell_\iota$ inert and every $\pi\in\linkNselZg_\iota$: if $\ell_\iota<\ell_A$ and $4L_w<\mathrm{ord}_H(\ell_\iota)$ (the range of (R2)), a local
representation bound condition for $(X.r_\iota,X.\rho_\iota,\pi)$ at the places above $\ell_\iota$ with constant $C_2$; otherwise a local representation
congruence condition for $(X.r_\iota,X.\rho_\iota,\pi)$.
\formal{Formalism.MainNotions.N.KillingConditions.JointParametersHubConditions}
\end{definition}

\begin{theorem}[Joint conditions along families]\label{pf:linkN-sel:famjoints}
Let $\mathcal F\in\sfname{FamP}(\eps)$ and $C_2\in\R$, and assume, for all data, Fontaine's period ring statement at every prime and the Fontaine
conditions at the places above $\ell_\iota$ for all local indices (as in \cref{pf:linkN-sel:jointprop}). Then there is $x_0$ such that for all $x\ge x_0$ with $x>x_1$, all $i\in I_x$ and all $HW$,
if $D_{x,i}$ satisfies the jump conclusion, then:
\begin{enumerate}[label=(\alph*)]
\item (\emph{joint condition B}) at every inert index $\iota$ with $\ell=\ell_\iota$ unramified in the host and in row~B, for every point $(r,\rho,\pi)$
at $\iota$ and every place $w\mid\ell$: $K_w$ is unramified over $\Ql$, and for every good coefficient extension at $w$ there is a word on the
letters of $w$ with the word base conditions and the thick-letter lattice surjection hypotheses, such that $\ell\ge11$, $j_A+1\le n_w$,
$C_2(j_A+1)(w+S+\eta+1)\le b$ and the thick letters form clusters of length at most $j_A$, where $j_A=j_A(L_{x,i},\ell)$;
\item (\emph{joint condition CD}) at every inert index with $\ell$ unramified, for every point and every $w\mid\ell$: $K_w$ is unramified, and for
every good coefficient extension at $w$ there is a word with the word base conditions, the thick-letter hypotheses, $\ell\ge11$ and the word
box condition, which either has a non-thick letter or is of alternating type with the alternating lattice surjection hypotheses.
\end{enumerate}
\formallabel{GlobalTheory.HTW.RecS.fam_joints}
\end{theorem}

\begin{proof}
Let $\eps_z>0$ be the second scale parameter of $\mathcal F$ and $A$ the bound of $\mathcal F$ for the parameter constants $C_{Lw}$; the scale
parameters of all members are $(\eps,\eps_z)$, both positive. By \lref{GlobalTheory.T.inertLongOrbitWordConditions_of_large_scale}
(deferred; about $310$ lines, through a dichotomy for twisted dual size words at long orbits) there is $x_B$ such that (a) holds for every datum
$D$, size data $L$ of scale $(x,\eps,\eps_z)$ with $x\ge x_B$, parameter choice with $C_{Lw}\le A$ and $HW$, under the two Fontaine
assumptions and the jump conclusion; by \lref{GlobalTheory.T.inertPrimeWordConditions_of_large_scale} (deferred; about $570$ lines) there is
$x_C$ with the same property for (b). Put $x_0=\max(x_B,x_C)$; for $x\ge x_0$, $x>x_1$ the members of $\mathcal F$ satisfy these assumptions,
since their scale is $x$ and their parameter constant equals that of the family.
\end{proof}

\begin{proposition}\label{pf:linkN-sel:jointprop}
Assume Fontaine's period ring, the Fontaine bounds and the Fontaine properties at the places above $\ell$ for all data, and the closure of the
jump conditions for all data. For every family $\mathcal F$ and all row constants $R$ there is $x_0$ such that for all $x\ge x_0$, $x>x_1$, $i\in I_x$,
all $HW$ and every hub datum $X$ whose representations form a point family and for which the jump parameters of $L_{x,i}$ are admissible, the joint
hub conditions hold for $(L_{x,i},HW,X,R)$.
\formallabel{GlobalTheory.HTW.RecS.fam_jointBCD_closed_alt}
\formallabel{GlobalTheory.HTW.RecS.fam_jointBCD}
\end{proposition}

\begin{proof}
By \cref{pf:linkN-sel:famjoints}, for every family and every $C_2$ there is
$x_0$ beyond which, at every member satisfying the jump conclusion, the word and Hodge--Tate data at the inert indices satisfy two families of joint
conditions; the lemma \lref{GlobalTheory.Survey.ClusteredLetterwiseBoxes.frozenBCD_of} (deferred) turns these, for a hub datum whose
representations satisfy a point condition at the inert indices, into the joint hub conditions. The jump conclusion follows from the admissibility of
the jump parameters by the closure of the jump conditions. The point condition at an inert index $\iota$ with good $\pi$ follows from the point
family by the lemma \lref{GlobalTheory.T.ptIn_of_ptFam_deRham_away}, which needs Fontaine's period ring, the de Rham property away from
the generator primes contained in the admissibility of the jump parameters, and $\ell_\iota\notin\linkNselG$; the last holds because an inert prime
is not a generator prime (\lref{GlobalTheory.T.inert_prime_ne_typePrime}, deferred).
\end{proof}

% ---------------------------------------------------------------------------
\subsection{Proof of the main theorem}\label{sec:pf:linkN-sel:proof}

\begin{proof}[Proof of \cref{pf:linkN-sel:main}]
\emph{Constants.} Let $C_0$ be given by \cref{pf:linkN-sel:thmE}, and $c_1=2080$, $Q=260X_0^2X_2^2(X_1+X_2)$ by \cref{pf:linkN-sel:thmBmain}. Let
$C_1\ge1$ and $x_p$ be the constant and threshold of the parameter bound \lref{GlobalTheory.T.admissible_parameters_family_const}
(eventually along $\mathcal F$, the field \emph{par} holds with constant $C_1$). Let $(C_a,x_a)$ be given by $(\mathrm B_{23})$ and $(C_b,x_b)$ by
\cref{pf:linkN-sel:starS}. Put $C_{23}=\max(0,C_a,C_b)$, $x_{23}=\max(1,x_a,x_b)$, $C_2=\max(C_0,1)$, and let $C_s$ be as in \eqref{eq:pf:linkN-sel:Cs}.
Then $C_s\ge1$, and the row constants $R(C_2)=(C_2,320,25,23,32,32)$ are admissible for $C_s$ in the sense of \cref{pf:linkN-sel:rownum}: $1\le C_2$,
$3C_2+1\le C_s$, $K_1=368$, $K_2+1=2083$, $K_3=3\cdot23\cdot320+12\cdot23+9=22365$, $6\cdot320+1=1921$ and $2(3\cdot23+1)=140$ are all at most
$C_s$, since $C_s\ge22365$.

\emph{Thresholds.} Let $x_{\mathrm{set}}$ be the threshold of \cref{pf:linkN-sel:setup} (for $C$, $C_2$, $C_T$), $x_{\mathrm{exc}}=x_1$ that of
\cref{pf:linkN-sel:exc}, $x_{\mathrm{sm}}$ that of \cref{pf:linkN-sel:small} for the constant $C_s\ge1$, $x_E$ that of \cref{pf:linkN-sel:congfam}, and
$x_{\mathrm{bcd}}$ that of \cref{pf:linkN-sel:jointprop} for $R(C_2)$, whose hypotheses are supplied by the Fontaine field of
$\sfname{Hypotheses}_2$ and the statements \lref{GlobalTheory.Per.FontaineAll.fontaine_bounds} (the Fontaine bounds for all data),
\lref{GlobalTheory.FontaineAdm.fontaineProperties_all_places} (the Fontaine properties at all places) and
\lref{GlobalTheory.JumpStatement.holds_all_alt} (the closure of the jump conditions). Put
$x_0=\max(x_{\mathrm{set}},x_{\mathrm{exc}},x_{\mathrm{sm}},x_E,x_{\mathrm{bcd}},x_p,x_{23})$.

\emph{A member.} Let $x\ge x_0$ with $x>x_1$, and $i\in I_x$ such that $D=D_{x,i}$ has finite odd Tate--Shafarevich groups; write $L=L_{x,i}$,
$P=P_{x,i}$, $\mathrm{nch}=\mathrm{nch}_{x,i}$. By \cref{pf:linkN-sel:exc} the set $\linkNselexc=\linkNselExc(L)$ is finite and satisfies the field
\emph{exc}. By \cref{pf:linkN-sel:setup} there are $HW$, Fontaine and comparison data and a hub datum $X$ with $\linkNselexc(X)=\linkNselExc(L)$,
admissible jump parameters, and a setup $(L,\mathrm{nch},HW,X,C_2)$. By $(\mathrm F)_{C,C_T}$, $D$ satisfies the norm-choice facts, the commuting
endomorphism condition and $\mathrm{PC}(C,C_T)$. By \cref{pf:linkN-sel:small}, the field \emph{small} holds for $(L,C_s)$ and $C_s\le b$; hence
$C_T\le C_s\le b$. We verify the fifteen fields of $\linkNselT(L,P,X,C_s)$.
\begin{description}[style=nextline,leftmargin=1.5em,font=\normalfont\itshape]
\item[pos] $C_s\ge0$.
\item[pass] \Cref{pf:linkN-sel:pass}, from the commuting endomorphism condition.
\item[mult] The result \lref{GlobalTheory.T.goodPointMultiplicityLeRank_proof} of another section, which derives the field \emph{mult} from three
conditions of the norm-choice facts (two conclusions on the attached representations and an isotypicity statement) and the uniqueness of
localisations.
\item[rank] A condition of the norm-choice facts.
\item[kill] Part of the setup.
\item[num] \Cref{pf:linkN-sel:rownum}, since $X$ has the canonical tables for $R(C_2)$, $R(C_2)$ is admissible for $C_s$, and
$S_c+\eta_{\mathrm{sz}}+t+2\le b$ (setup).
\item[mid, bb] The joint hub conditions hold for $(L,HW,X,R(C_2))$ by \cref{pf:linkN-sel:jointprop} (the representations of $X$ form a point family
since $X$ is canonical). By \cref{pf:linkN-sel:thmE}, the lifting corollary holds with constant $C_2\ge C_0$. Since $C_2\ge1$, the middle-range
theorem \lref{GlobalTheory.T.star_ad_and_bb_le_of_joint_cond} of another section, which states that under $\sfname{Hypotheses}_2$, for a setup
$(L,\mathrm{nch},HW,X,C_2)$ with $C_2\ge1$, the lifting corollary with constant $C_2$ and the joint hub conditions for $R(C_2)$, the field \emph{mid}
holds and the field \emph{bb} holds with constant $7200$, gives \emph{mid}, and \emph{bb} with constant $7200\le C_s$; \cref{pf:linkN-sel:mono}
gives \emph{bb} with $C_s$.
\item[large] \Cref{pf:linkN-sel:congfam} gives $\mathrm{CB}(D,L)$; the weight parameter condition is part of the norm-choice facts, $(X.r,X.\rho)$ is a
point family with $\mathrm{Kill}(X.\rho)$, and $3\le\linkNsellst$ by \emph{small}; so \cref{pf:linkN-sel:large} with $\linkNselexc(X)=\linkNselExc(L)$
gives \emph{large}.
\item[exc] \Cref{pf:linkN-sel:exc}, as $\linkNselexc(X)=\linkNselExc(L)$.
\item[two] Since $x\ge x_a$ and $x\ge x_b$, $(\mathrm B_{23})$ and \cref{pf:linkN-sel:starS}, applied to the setup, give the local bounds with
$\linkNselbsm(L,C_a)$ at $\ell=2$ and at $\ell=3\notin\linkNselP$, and with $\linkNselbsm(L,C_b)$ at the odd small primes. As $C_a,C_b\le C_{23}$ and
$C_{23}\ge0$, \cref{pf:linkN-sel:mono} replaces both by $\linkNselbsm(L,C_{23})$, and \cref{pf:linkN-sel:starF} gives \emph{two} with $C_{23}$; then
\cref{pf:linkN-sel:mono} gives it with $C_s\ge C_{23}$.
\item[ram] \Cref{pf:linkN-sel:thmBmain} gives $\mathrm{TB}(X.\rho,L,c_1,Q)$; the parameter choice satisfies $x^{2C_t+5}<q,d$; and $C_s\ge C_G$. So
\cref{pf:linkN-sel:ram} gives \emph{ram} with $C_s$.
\item[typ] \Cref{pf:linkN-sel:typprop} with $C'=C\ge0$, the constant $C_1\ge1$ (the field \emph{par} with $C_1$ holds as $x\ge x_p$), $C_T\le b$,
$\mathrm{PC}(C,C_T)$, the field \emph{mult} and the size condition on the generator primes (a size condition of the member) gives \emph{typ} with
$2000(C+1)+5C_1\le C_s$; \cref{pf:linkN-sel:mono} gives it with $C_s$.
\item[par] The field \emph{par} with $C_1\le C_s$ and \cref{pf:linkN-sel:mono}.
\item[small] \Cref{pf:linkN-sel:small}.
\end{description}
Hence $\linkNselT(L,P,X,C_s)$ holds for all such members, which is the claim with the constants $C_s$ and $x_0$.
\end{proof}

% ---------------------------------------------------------------------------
\subsection{Local results}\label{sec:pf:linkN-sel:local}

This subsection proves local results used above and in \cref{sec:pf:linkN-adm}. Throughout, $p$ is a prime, $K$ a finite extension
of $\Q_p$ with maximal unramified subfield $K_0$, $G$ a closed subgroup of $G_K$ (mostly $G=G_K$), $\mathit{LL}$ a finite extension of
$\Q_p$ with ring of integers $\OO$, $\chi$ the cyclotomic character and $B_{\mathrm{cris}}$ the crystalline period ring of the local
theory. A \emph{lattice} $T$ is a finite-dimensional continuous $\mathit{LL}$-representation $V=V(T)$ of $G$ together with a
$G$-stable finitely generated $\OO$-submodule of $V$ spanning it, also denoted $T$; a \emph{lattice homomorphism} is a continuous
$G$-equivariant $\mathit{LL}$-linear map carrying lattice into lattice, and a \emph{short exact sequence of lattices}
$0\to T_0\xrightarrow{\iota}T_1\xrightarrow{\pi}T_2\to0$ is a pair of lattice homomorphisms which is exact on the representations,
with $\pi(T_1)=T_2$ and $\iota^{-1}(T_1)=T_0$. For a lattice $T\subseteq V$ we use the following $\OO$-modules of maps $G\to V/T$:
the \emph{divisible cocycles} $Z^1_{\mathrm{div}}(T)$, the locally constant maps $c$ with $c(gh)=c(g)+gc(h)$; the
\emph{coboundaries} $B(T)=\{g\mapsto gx-x:x\in V/T\}$; the \emph{finite classes} $H_f(T)=\overline{Z^1_f(V)}+B(T)$, where $Z^1_f(V)$
is the space of continuous cocycles $c\colon G\to V$ with $c(g)\otimes1=gx-x$ for some $x\in V\otimes_{\Q_p}B_{\mathrm{cris}}$ and the
bar denotes reduction modulo $T$; and, for a subgroup $I\subseteq G$, the \emph{unramified classes} $Z^1_{\mathrm{ur}}(T,I)$, the maps
$c$ for which there is $x\in V/T$ with $c(g)=gx-x$ for all $g\in I$. Put
$H^1_{\mathrm{div}}(T)=Z^1_{\mathrm{div}}(T)/(Z^1_{\mathrm{div}}(T)\cap B(T))$ and let $H^1_{\mathrm{div},f}(T)$ be the image of
$Z^1_{\mathrm{div}}(T)\cap H_f(T)$ in it. Finally $H^1(G,T)$ is the module of continuous cocycles with values in $T$ modulo the
coboundaries of elements of $T$, and $H^1_f(T)\subseteq H^1(G,T)$ the submodule of classes of cocycles in $Z^1_f(V)$. Lengths are
$\OO$-lengths with values in $\N\cup\{\infty\}$.

\subsubsection{Two length estimates}

\begin{lemma}[Torsion lengths]\label{pf:linkN-sel:torslen}
Let $\OO$ be a commutative ring and $a\in\OO$. If an $\OO$-module $P$ is generated by $n$ elements each killed by $a$, then
$\len(P)\le n\,\len(\OO/a\OO)$. Moreover $\len(\OO/a^e\OO)\le e\,\len(\OO/a\OO)$ for all $e\in\N$.
\formallabel{Formalism.WachCohomologyFacts.PLenBound.length_le_of_span}
\formallabel{Formalism.WachCohomologyFacts.PLenBound.length_quot_pow_le}
\end{lemma}

\begin{proof}
$P$ is a quotient of $(\OO/a\OO)^n$, and length is monotone under quotients and additive in short exact sequences. For the second
claim argue by induction on $e$: the exact sequence $0\to a^e\OO/a^{e+1}\OO\to\OO/a^{e+1}\OO\to\OO/a^e\OO\to0$ and the surjection
$\OO/a\OO\to a^e\OO/a^{e+1}\OO$, $x\mapsto a^ex$, give $\len(\OO/a^{e+1}\OO)\le\len(\OO/a\OO)+\len(\OO/a^e\OO)$.
\end{proof}

\begin{lemma}[Two-step count]\label{pf:linkN-sel:twostep}
Let $\OO$ be a commutative ring, $M$ an $\OO$-module, $\ell,g,g_0,e,e'\in\N$, $I\subseteq Z\subseteq M$ submodules and $s,s_0\subseteq M$
finite sets with $\#s\le g$, $\#s_0\le g_0$, $s_0\subseteq Z$, $\OO s+I=Z$, $\ell^eZ\subseteq\OO s_0+I$ and $\ell^{e'}\OO s_0\subseteq I$. Then
\[
\len(Z/I)\le(ge+g_0e')\,\len(\OO/\ell\OO).
\]
\formallabel{Formalism.TruncatedSeriesCounting.TruncatedSeriesCountingBounds.length_quotient_le_of_two_step}
\end{lemma}

\begin{proof}
Since $s_0\subseteq Z$ and $I\subseteq Z$, $Z'=\OO s_0+I$ is a submodule of $Z$ containing $I$, and $0\to Z'/I\to Z/I\to Z/Z'\to0$ is exact,
so $\len(Z/I)=\len(Z'/I)+\len(Z/Z')$. The module $Z'/I$ is generated by the images of the elements of $s_0$, each killed by $\ell^{e'}$;
by \cref{pf:linkN-sel:torslen}, $\len(Z'/I)\le g_0\len(\OO/\ell^{e'}\OO)\le g_0e'\len(\OO/\ell\OO)$. The module $Z/Z'$ is generated by the
images of the elements of $s$ (as $\OO s+I=Z$), each killed by $\ell^e$ (as $\ell^eZ\subseteq Z'$); hence $\len(Z/Z')\le ge\len(\OO/\ell\OO)$.
\end{proof}

\subsubsection{The local Bloch--Kato condition}

A \emph{local situation} $S$ consists of the period data of $G=G_K$ at $p$, a coefficient field $\mathit{LL}$ and a Wach structure $W$ of
level $s$ (\cref{pf:linkN-adm:wach}), and $G_{K_s}\subseteq G$ denotes the subgroup on which $\chi\equiv1\bmod p^s$. A lattice $Y$ is
\emph{crystalline on} a subgroup $H\subseteq G$ if $\dim_{\mathit{LL}}(V(Y)\otimes B_{\mathrm{cris}})^H=[K_0:\Q_p]\dim V(Y)$. Given a
comparison datum $E$ for $S$ and $M_0,h_0\in\N$, the \emph{Wach hypothesis} $\mathrm{HypW}(Y,M_0,h_0)$ requires that $Y$ is crystalline on
$G_{K_s}$, that its jumps lie in $[-M_0,h_0]$, and that $D(Y)$ contains an $A$-submodule satisfying the Wach clauses of $E$.

\begin{definition}[The local theorem TB]\label{pf:linkN-sel:TBloc}
The \emph{local theorem TB} holds at $(S,E)$ if the following holds for all $e,M_0,h_0\in\N$. Let
$0\to Y\xrightarrow{\iota}\tilde Y\xrightarrow{\pi}Z\to0$ be a short exact sequence of lattices with $Z$ of rank one and trivial action, let
$T_e$ be a lattice of rank one on which $g$ acts by $\chi(g)^{-e}$, assume $\tilde Y$ crystalline on $G_{K_s}$ and $\mathrm{HypW}(Y,M_0,h_0)$,
and let $\tilde y\in\tilde Y$ with $\OO\pi(\tilde y)=Z$, $\varphi\colon Y\to V(T_e)/T_e$ an $\OO$-linear $G$-equivariant map and
$f\colon G\to Y$ with $\iota(f(g))=g\tilde y-\tilde y$. Then $p^{c_{\mathrm{TB}}(Y)}\,\varphi\circ f\in H_f(T_e)$, where
$c_{\mathrm{TB}}(Y)=c_{\mathrm{TB}}(Y;M_0,h_0,e)\in\N$ is an explicit constant depending on $p$, $s$, $e$, $K$, the jump set of $Y$ and
$M_0+h_0$.
\formal{Formalism.MainNotions.LocalSituation.TBat}
\formal{Formalism.MainNotions.LocalSituation.cTBY}
\end{definition}

\begin{lemma}[TB for a vector of value one]\label{pf:linkN-sel:tbpart1}
Assume the local theorem TB at $(S,E)$. Let $\tilde Y$ be a lattice which is crystalline on $G_{K_s}$, $\lambda\colon V(\tilde Y)\to\mathit{LL}$ a
$G$-invariant linear form with $\lambda(\tilde Y)\subseteq\OO$, and $e_0\in\tilde Y$ with $\lambda(e_0)=1$; let $Y$ be the lattice
$\tilde Y\cap\ker\lambda$ of $\ker\lambda$. Let $e,M_0,h_0,c_B\in\N$ with $\mathrm{HypW}(Y,M_0,h_0)$ and $c_{\mathrm{TB}}(Y;M_0,h_0,e)\le c_B$,
let $T_e$ be as in \cref{pf:linkN-sel:TBloc}, and let $\varphi\colon V(\tilde Y)\to V(T_e)/T_e$ be $\OO$-linear with
$\varphi(gy)=g\varphi(y)$ for $y\in Y$. Then $p^{c_B}\cdot\bigl(g\mapsto\varphi(ge_0-e_0)\bigr)\in H_f(T_e)$.
\formallabel{GlobalTheory.TBLocal.tb_part1}
\formallabel{GlobalTheory.TBLocal.ses}
\end{lemma}

\begin{proof}
Let $Z=\OO\subseteq\mathit{LL}$ with trivial action. The inclusion $\iota\colon\ker\lambda\to V(\tilde Y)$ and $\lambda\colon V(\tilde Y)\to\mathit{LL}$
form a short exact sequence of lattices $0\to Y\to\tilde Y\to Z\to0$: $\lambda$ is onto since $\lambda(ze_0)=z$; $\lambda(\tilde Y)=\OO$ since
$\lambda(\tilde Y)\subseteq\OO$ and $\lambda(ae_0)=a$; the kernel of $\lambda$ is the image of $\iota$; and $\iota^{-1}(\tilde Y)=Y$ by
definition. As $\lambda$ is $G$-invariant, $f(g)=ge_0-e_0$ lies in $Y$, and $\OO\lambda(e_0)=Z$. Apply TB with $\tilde y=e_0$, the restriction
of $\varphi$ to $Y$ and this $f$: $p^{c_{\mathrm{TB}}(Y)}\bigl(g\mapsto\varphi(ge_0-e_0)\bigr)\in H_f(T_e)$. Multiplying by
$p^{c_B-c_{\mathrm{TB}}(Y)}$ gives the claim, $H_f(T_e)$ being an $\OO$-module.
\end{proof}

\begin{lemma}[The clauses of TBe]\label{pf:linkN-sel:tbe}
Let $T$ be a lattice, $I=J$ subgroups of $G$, and $e,c_B\in\N$. Assume: if $e>0$, then $H_f(T)=B(T)$; if $e=0$, then a map $c$ lies in
$H_f(T)$ if and only if $c\in Z^1_{\mathrm{div}}(T)$ and $c(g)=gx-x$ for all $g\in J$, for some $x$; and the $p^{c_B}$-torsion of
$H^1_{\mathrm{div}}(T)/H^1_{\mathrm{div},f}(T)$ has length at most $([K:\Q_p]+1)c_B\len(\OO/p\OO)$. Then:
\begin{enumerate}[label=(\alph*)]
\item if $e>0$, then $Z^1_{\mathrm{div}}(T)\cap H_f(T)\subseteq B(T)$;
\item if $e=0$, then $Z^1_{\mathrm{div}}(T)\cap H_f(T)=Z^1_{\mathrm{div}}(T)\cap\bigl(Z^1_{\mathrm{ur}}(T,I)+B(T)\bigr)$;
\item with $A=\{c\in Z^1_{\mathrm{div}}(T):p^{c_B}c\in H_f(T)\}$, $\len\bigl(A/(Z^1_{\mathrm{div}}(T)\cap H_f(T))\bigr)\le([K:\Q_p]+1)c_B\len(\OO/p\OO)$.
\end{enumerate}
\formallabel{GlobalTheory.TBLocal.tbe_clauses}
\formallabel{GlobalTheory.TBLocal.len_le_torsion}
\end{lemma}

\begin{proof}
(a) is immediate. (b) If $c\in Z^1_{\mathrm{div}}\cap H_f$, the hypothesis for $e=0$ gives $x$ with $c(g)=gx-x$ on $J=I$, so $c\in Z^1_{\mathrm{ur}}(T,I)$.
Conversely let $c=y+b\in Z^1_{\mathrm{div}}$ with $y(g)=gx-x$ on $I$ and $b(g)=gu-u$; then $c(g)=g(x+u)-(x+u)$ on $J$, so $c\in H_f$ by the
hypothesis. (c) For $c\in A$ the class of $c$ in $Q=H^1_{\mathrm{div}}(T)/H^1_{\mathrm{div},f}(T)$ is killed by $p^{c_B}$, since $p^{c_B}c\in
Z^1_{\mathrm{div}}\cap H_f$. This gives an $\OO$-linear map from $A$ to the $p^{c_B}$-torsion of $Q$. Its kernel is $Z^1_{\mathrm{div}}\cap H_f$:
if the class of $c$ is that of some $w\in Z^1_{\mathrm{div}}\cap H_f$, then $c-w\in B(T)\subseteq H_f(T)$, so $c\in H_f$; the converse is clear.
Hence $A/(Z^1_{\mathrm{div}}\cap H_f)$ embeds into the torsion submodule, and the length bound follows from the hypothesis.
\end{proof}

\subsubsection{Generators of finite classes}

\begin{proposition}[Generator bound]\label{pf:linkN-sel:gen}
For every datum $D$, every prime $\ell$ and all finite extensions $K,\mathit{LL}$ of $\Ql$ with $K=K_0$, every crystalline
$\OO_{\mathit{LL}}$-lattice $T$ of $G_K$ has a generating set of $H^1_f(T)$ with at most $([K_0:\Ql]+1)\rank T$ elements.
\formallabel{GlobalTheory.T.h1fGeneratorBoundAll_loc}
\formallabel{GlobalTheory.GenCountUnram.admGenL_all}
\end{proposition}

\begin{proof}
Let $V=V(T)$. Crystalline representations are de Rham (\lref{GlobalTheory.Per.CrisDR.isDR_of_isCrys}, which states this for the
period data of the datum). For de Rham $V$ the lemma \lref{GlobalTheory.T.h1fT_span_card_bound} (deferred) gives a generating set $s$
of $H^1_f(T)$ with $\#s+\dim H^0(G_K,V)\le\dim H^1_f(V)+\dim V$, where $H^1_f(V)$ is the space of classes of $Z^1_f(V)$; the dimension
formula of Bloch--Kato, \lref{GlobalTheory.SeqDim.h1f_finrank}, gives
$\dim H^1_f(V)=\dim H^0(G_K,V)+\dim D_{\mathrm{dR}}(V)-\dim\mathrm{Fil}^0D_{\mathrm{dR}}(V)$ \cite{BlochKato1990}; and
$\dim D_{\mathrm{dR}}(V)\le[K:\Ql]\dim V$ (\lref{GlobalTheory.Per.TiedDdR.ddR_finrank_le}). Since $K=K_0$,
$\#s\le\dim D_{\mathrm{dR}}(V)+\dim V\le([K_0:\Ql]+1)\dim V$. The ring $\OO_{\mathit{LL}}$ satisfies the integrality conditions required
there by \lref{GlobalTheory.LocalIntegers.good_integers}.
\end{proof}

\subsubsection{Herr classes}

For a Wach structure $W$ of level $s$ and a lattice $T$, a \emph{divisible Wach cocycle} is a pair $(y,z)\in D(T)^2$ satisfying the cocycle
condition of the Herr complex \cite{Herr1998}; $W$ provides a class map $\mathrm{cl}_T$ from divisible Wach cocycles to
$H^1(G_{K_s},T)$ and, for the $\Delta$-invariant version (pairs in the $\Delta$-invariant Wach cocycles), a class map $\mathrm{cl}^K_T$ to
$H^1(G_K,T)$. We say that $W$ has \emph{natural Herr classes} if:
\begin{enumerate}[label=(N\arabic*)]
\item for every lattice homomorphism $f\colon T\to T'$ and divisible Wach cocycles $(y,z)$ of $T$ and $(y',z')=(D(f)y,D(f)z)$ of $T'$,
$f_*\mathrm{cl}_T(y,z)=\mathrm{cl}_{T'}(y',z')$ in $H^1(G_{K_s},T')$;
\item if $p=2$ implies $\sqrt{-1}\in K$, the same holds for $\mathrm{cl}^K$;
\item (extension formula) for lattice homomorphisms $\iota\colon T\to T_p$ and $\pi\colon T_p\to\mathbf 1$ between crystalline lattices ($\mathbf 1$
the trivial lattice of $W$), every lattice cocycle $c$ of $T$ which is an extension cocycle of $(\iota,\pi)$ and every unit generator $e$ of
$D(\mathbf 1)$, there is $u\in\OO^\times$ such that
for all $x\in D(T_p)$, $a\in\OO$ with $D(\pi)x=ae$ and all divisible Wach cocycles $(y,z)$ of $T$ with $D(\iota)y=\varphi x-x$ and
$D(\iota)z=\gamma x-x$, one has $\mathrm{cl}_T(y,z)=ua\cdot\mathrm{res}_{G_{K_s}}[c]$;
\item if $s=1$ and $p\ne2$, the analogue of (N3) for $\Delta$-invariant $x$, $\Delta$-invariant Wach cocycles and $\mathrm{cl}^K$, with
$\gamma_K$ in place of $\gamma$ and without restriction.
\end{enumerate}

\begin{proposition}[Natural Herr classes]\label{pf:linkN-sel:herr}
For every datum $D$ with local Hodge--Tate data $HW$ and Fontaine data, the Wach structure of every good triple $(K,s,\mathit{LL})$ at every
prime has natural Herr classes.
\formallabel{GlobalTheory.AdmHerrLocal.main}
\formallabel{Formalism.MainNotions.HerrComplexNaturality.herr_naturality_theorem_holds}
\formallabel{Formalism.MainNotions.HerrComplexNaturality.herr_class_maps_of_naturality}
\end{proposition}

\begin{proof}
Herr's theorem in the form \lref{GlobalTheory.HerrFull.herr_a_and_natural} states, at every good triple, two forms of part (a) of
Herr's theorem (injectivity and bijectivity properties of the class maps) together with (N1); the first form contains, under the hypothesis
of (N2), the naturality of $\mathrm{cl}^K$, which is (N2). The extension formulas (N3) and (N4), with a unit
$u$, are \lref{GlobalTheory.HerrExtUnit.main}.
\end{proof}

\subsubsection{Fontaine's functor}

Let $W$ be a Wach structure (\cref{pf:linkN-adm:wach}) with ring $A=\OO[[X]]^{\Z/n}$, \'etale ring $\mathcal E\supseteq A$, Frobenius
$\varphi$ and $\Gamma$-action, and \'etale $(\varphi,\Gamma)$-modules $D(T)$ over $\mathcal E$ for the lattices $T$. We say that \emph{Fontaine's
theory holds at $W$} if:
\begin{enumerate}[label=(F\arabic*)]
\item (full faithfulness) for lattices $T,T'$ and every $\mathcal E$-linear map $f\colon D(T)\to D(T')$ commuting with $\varphi$ and $\Gamma$
there is a unique lattice homomorphism $F\colon T\to T'$ with $D(F)=f$;
\item (exactness) for every short exact sequence of lattices $0\to T_0\xrightarrow{\iota}T_1\xrightarrow{\pi}T_2\to0$, $D(\iota)$ is injective,
$D(\pi)$ is surjective and $\ker D(\pi)=\operatorname{im}D(\iota)$;
\item ($\psi$ on $\mathcal E$) there is an $\OO$-linear $\psi\colon\mathcal E\to\mathcal E$ commuting with $\Gamma$, with $\psi(\varphi(a)x)=a\psi(x)$,
and with $\psi\bigl((1+X)^j\varphi(a)\bigr)=a$ for $j=0$ and $=0$ for $0<j<p$; and every $x\in\mathcal E$ can be written
$x=\sum_{j=0}^{p-1}(1+X)^j\varphi(c_j)$ with $c_j\in\mathcal E$;
\item ($\psi$ on modules) for every lattice $T$ there is an $\OO$-linear $\psi\colon D(T)\to D(T)$ commuting with $\Gamma$, with $\psi(\varphi x)=x$
and $\psi(\varphi(a)x)=a\psi(x)$ for $a\in\mathcal E$.
\end{enumerate}
These are the basic properties of Fontaine's equivalence between lattices and \'etale $(\varphi,\Gamma)$-modules and of the operator $\psi$
\cite{Fontaine1990,Herr1998}.

\begin{theorem}[Fontaine's functor]\label{pf:linkN-sel:fontaine}
For every datum $D$ with local Hodge--Tate data $HW$ and Fontaine data, Fontaine's theory holds at the Wach structure of level $1$ of every
place $(\ell,K,\mathit{LL})$ in case~A.
\formallabel{GlobalTheory.T.fontaine_fullFaithful_exact_psi}
\formallabel{GlobalTheory.fontaine_litFont_assembled}
\end{theorem}

\begin{proof}
Let $W$ be the Wach structure; the residue characteristic $p$ of its period data is $\ell$. (F1) is
\lref{GlobalTheory.T.wachFunctor_fullyFaithful} (deferred), which states (F1) at every good triple. (F2) is
\cref{pf:linkN-sel:exact} below. For (F3), \lref{GlobalTheory.T.wachFunctorRing_psiOperator_exists_of_good} (deferred) gives an
$\OO$-linear $\psi$ on $\mathcal E$ with the three displayed properties, and \lref{GlobalTheory.T.etaleRing_free_over_frobenius_image}
(deferred) states that every $x\in\mathcal E$ has a unique expansion $x=\sum_{j<p}(1+X)^j\varphi(c_j)$, i.e.\ that $\mathcal E$ is free over
$\varphi(\mathcal E)$ with basis $(1+X)^j$, $0\le j<p$; only the existence of the expansion is used. (F4) is
\lref{GlobalTheory.T.exists_psi_operator_etale_phiGamma_module} (deferred), which constructs $\psi$ on every $D(T)$.
\end{proof}

\begin{lemma}[Exactness of the \'etale functor]\label{pf:linkN-sel:exact}
For every datum with local Hodge--Tate and Fontaine data, the functor $D$ of the Wach structure of every good triple satisfies (F2).
\formallabel{GlobalTheory.LatExact.W_exact}
\formallabel{GlobalTheory.T.etaleFunctor_exact_on_lattices}
\end{lemma}

\begin{proof}
By the Fontaine and decompletion statements \lref{GlobalTheory.Decompl.fontaine_and_decompletion}, the hypotheses of
\lref{GlobalTheory.EtaleExact.DD_exact} are satisfied; the latter states that $D$ carries short exact sequences of lattices to exact
sequences of $\mathcal E$-modules: $D(\iota)$ injective, $\ker D(\pi)=\operatorname{im}D(\iota)$ and $D(\pi)$ surjective.
\end{proof}

\subsubsection{The arc tools}

\begin{definition}[The arc tools]\label{pf:linkN-sel:arcdef}
Let $D$ be a datum with local Hodge--Tate data, Fontaine data and comparison data. The \emph{arc tools} are the following thirteen
statements on the Wach structures of level $1$ at the places in case~A with $\ell\ge5$ (in the notation of
\cref{pf:linkN-adm:def-AB}): the saturation statement (sat) for arc lifting and its \'etale form (sat$_E$); the saturation statement (satP)
for Frobenius pullbacks; the statement $\gamma(\pi_0)\in\pi_0A^\Delta$; the killer statement at levels (kLev), the killer bound with
constant $8$, and the second form (kLev$_E$) of (kLev); the first and second filtration statements (F1), (F2) for arc lifting and their
\'etale forms (F1$_E$), (F2$_E$); the filtration properties (FP$_E$) in \'etale coordinates, which involve the comparison data; and the
freeness (bc) of finite cohomology under unramified base change.
\formal{Formalism.MainNotions.ArcLiftingSetup.adm5}
\formal{Formalism.MainNotions.WachCoordinates.adm8}
\formal{Formalism.MainNotions.ArcLiftingSetup.SaturationFrobeniusPullback}
\formal{Formalism.MainNotions.UnusedSlot80.GammaPreservesUniformiser}
\formal{Formalism.MainNotions.ArcLiftingSetup.KillerLevL}
\formal{Formalism.MainNotions.ArcLiftingSetup.adm8}
\formal{Formalism.MainNotions.WachCoordinates.adm11}
\formal{Formalism.MainNotions.ArcLiftingSetup.adm6}
\formal{Formalism.MainNotions.ArcLiftingSetup.adm7}
\formal{Formalism.MainNotions.WachCoordinates.adm9}
\formal{Formalism.MainNotions.WachCoordinates.adm10}
\formal{Formalism.MainNotions.WachCoordinates.adm7}
\formal{Formalism.MainNotions.WachCoordinates.adm13}
\end{definition}

\begin{theorem}[The arc tools]\label{pf:linkN-sel:arc}
For every datum $D$ with local Hodge--Tate data, Fontaine data and comparison data, the arc tools hold.
\formallabel{GlobalTheory.T.arcTheorem_tools_of_wachModuleConclusions}
\formallabel{GlobalTheory.ArcTools.arc_tools_from_clauses}
\end{theorem}

\begin{proof}
The later clauses of the Wach lattice theory \lref{GlobalTheory.T.wachLattice_laterClauses_hold} give the clauses (P1)--(P10) and
further clauses, among them two clauses (a), (b) on Wach cohomology classes; \cref{pf:linkN-sel:herr} gives the naturality (N1), (N2) of the
Herr classes. Each arc tool is then derived by a lemma of the development (all deferred):
\begin{itemize}
\item (sat) and (sat$_E$) from (P1) and (P9) (\lref{Formalism.MainNotions.UnusedSlot80.saturation_general_statement},
\lref{Formalism.MainNotions.UnusedSlot80.saturation_unramified_statement}); (satP) from (P9)
(\lref{Formalism.MainNotions.UnusedSlot80.saturation_pullback_statement}); and $\gamma(\pi_0)\in\pi_0A^\Delta$ from (P9)
(\lref{Formalism.MainNotions.UnusedSlot80.gamma_uniformizer_divisible_statement});
\item (kLev) from the comparison data, (P4) and (P9) (\lref{Formalism.MainNotions.UnusedSlot80.annihilator_levels_unramified_statement},
about $700$ lines, through a grading of the annihilator ideals by levels); the killer bound with constant $8$ from (kLev) and
$\gamma(\pi_0)\in\pi_0A^\Delta$ (\lref{Formalism.MainNotions.UnusedSlot80.killer_first_theorem_holds}, about $280$ lines, through a
$\pi$-stability statement for contents); and (kLev$_E$) from $\gamma(\pi_0)\in\pi_0A^\Delta$
(\lref{Formalism.MainNotions.UnusedSlot80.killer_second_theorem_holds});
\item (F1) and (F1$_E$) from (P2), (P5), (P9), the clause (a) and (N1), (N2)
(\lref{Formalism.MainNotions.UnusedSlot80.firstFiltrationTheoremCaseA},
\lref{Formalism.MainNotions.UnusedSlot80.firstFiltrationTheoremEtale}); (F2) and (F2$_E$) from (P7), (P9), the clause (b) and
(F1), resp.\ (F1$_E$) (\lref{Formalism.MainNotions.UnusedSlot80.second_filtration_theorem_arc_lifting},
\lref{Formalism.MainNotions.UnusedSlot80.cohomology_class_unramified_statement});
\item (FP$_E$) from the comparison data, (P3), (P4) and (P9) (\lref{Formalism.MainNotions.UnusedSlot80.tools_unramified_statement});
and (bc) for every datum (\lref{Formalism.MainNotions.UnusedSlot80.base_change_statement}).\qedhere
\end{itemize}
\end{proof}

\subsubsection{Crystalline characters with given jumps}

\begin{proposition}[Crystalline characters with given jumps]\label{pf:linkN-sel:crychar}
Let $D$ be a datum with local Hodge--Tate data and Fontaine data, $\ell$ a prime and $(K,\mathit{LL})$ a pair in case~A, so that the letters of the
Wach structure of level $1$ correspond to the embeddings $\tau\colon K\to\mathit{LL}$.
\begin{enumerate}[label=(\alph*)]
\item For every function $a$ from the letters to $\Z$ there is a continuous crystalline character $\psi\colon G_K\to\OO^\times$ which has, at every
embedding $\tau$ over the letter $i$, the single Hodge--Tate jump $a(i)$, with multiplicity one.
\item If $\rho$ is a crystalline representation of $G_K$ and $\mu\colon G_K\to\mathit{LL}^\times$ a continuous character trivial on the inertia group,
then $\mu\otimes\rho$ is crystalline and has the same Hodge--Tate multiplicities as $\rho$ at every embedding.
\end{enumerate}
\formallabel{GlobalTheory.T.crystalline_characters_given_jumps_unramified_twists}
\formallabel{GlobalTheory.CrysCharJumps.adm20_of_fonOf}
\formallabel{GlobalTheory.CrysCharJumps.hasJumps_prod}
\end{proposition}

\begin{proof}
(a) For every embedding $\tau_0$, \lref{GlobalTheory.T.exists_crystalline_char_jump_one_alt} (deferred; a character of Lubin--Tate type) gives a continuous crystalline character $\psi_{\tau_0}\colon G_K\to\OO^\times$ whose only Hodge--Tate jump is $1$
at $\tau_0$ and $0$ at the other embeddings, all with multiplicity one. Put $c(\tau)=a(i)$ for $\tau$ over the letter $i$ and
$\psi=\prod_\tau\psi_\tau^{c(\tau)}$. By induction over finite sets $s$ of embeddings, $\prod_{\tau\in s}\psi_\tau^{c(\tau)}$ is crystalline with the
single jump $c(\tau)$ at $\tau\in s$ and $0$ at $\tau\notin s$: for $s=\emptyset$ this is the trivial character, and the induction step uses that
integral powers and products of crystalline characters are crystalline, with jumps multiplied, resp.\ added
(\lref{GlobalTheory.CrysCharJumps.hasJumps_one}, \lref{GlobalTheory.CrysCharJumps.hasJumps_zpow},
\lref{GlobalTheory.CrysCharJumps.hasJumps_mul}). For $s$ the set of all embeddings this is (a).

(b) A continuous unramified character is crystalline with the single Hodge--Tate weight $0$
(\lref{GlobalTheory.T.unramifiedChar_isCrystalline_hodgeTate_zero}), and the twist of a crystalline representation by a crystalline
character is crystalline, with multiplicities shifted by the weight of the character
(\lref{GlobalTheory.TwistRepCrys.isCrys_and_multV_twist}); with weight $0$ the multiplicities are unchanged.
\end{proof}
